\documentclass[a4paper,12pt]{article}
\usepackage[spanish]{babel}
\usepackage[utf8]{inputenc}
\usepackage[T1]{fontenc}
\usepackage{graphicx} % Para insertar imágenes
\usepackage{amsmath, amssymb} % Símbolos matemáticos
\usepackage{hyperref} % Enlaces y referencias clicables
\usepackage{lmodern} % Fuente moderna
\usepackage{geometry} % Márgenes
\usepackage{polynom}
\usepackage[most]{tcolorbox} % Cajitas para ejemplos y procedimientos
\usepackage{newunicodechar}
\usepackage{cancel}
\usepackage[svgnames]{xcolor}
\newunicodechar{⋆}{\star}
\geometry{margin=2.5cm}

% ==== Definición de estilos de caja ====
\tcbset{
  colback=white,
  colframe=black!60,
  sharp corners,
  boxrule=0.8pt,
  left=6pt,right=6pt,top=6pt,bottom=6pt
}

\newtcolorbox{procedimiento}[1][]{
  colback=green!5!white,
  colframe=green!60!black,
  title=\textbf{Procedimiento},
  fonttitle=\bfseries,
  #1
}

\newtcolorbox{ejemplo}[1][]{
  colback=blue!5!white,
  colframe=blue!60!black,
  title=\textbf{Ejemplo},
  fonttitle=\bfseries,
  #1
}

\newtcolorbox{ejercicio}[1][]{
  colback=red!5!white,
  colframe=red!60!black,
  title=\textbf{Ejercicio},
  fonttitle=\bfseries,
  sharp corners,
  boxrule=0.8pt,
  left=6pt,right=6pt,top=6pt,bottom=6pt,
  #1
}

% ==== Documento principal ====
\title{Álgebra Lineal y Estructuras Matemáticas (ALEM) Ejercicios Resueltos Tema 3}
\author{Alonso Hernández Robles (Grupo E)}
\date{23/12/2025}

\begin{document}

\maketitle

\tableofcontents
\newpage

\section{Operaciones con Matrices}

\begin{ejercicio}[title=Ejercicio 1]
Sean $A,B \in \mathfrak{M}_2(\mathbb{Z}_7)$ tales que
\[
A + B = 
\begin{pmatrix}
4 & 3\\
2 & 6
\end{pmatrix}, 
\qquad
A - B =
\begin{pmatrix}
5 & 1\\
4 & 2
\end{pmatrix}.
\]
Calcule la matriz $A^2 - B^2$.

\tcblower
\textbf{Resolución:}

Planteamos las dos ecuaciones como un sistema:
\[
\begin{cases}
A + B = \begin{pmatrix}4 & 3\\[2pt] 2 & 6\end{pmatrix}\\[4pt]
A - B = \begin{pmatrix}5 & 1\\[2pt] 4 & 2\end{pmatrix}
\end{cases}
\]

Sumando las dos ecuaciones:
\[
2A =
\begin{pmatrix}
4+5 & 3+1\\[2pt]
2+4 & 6+2
\end{pmatrix}
=
\begin{pmatrix}
2 & 4\\[2pt]
6 & 1
\end{pmatrix}.
\]

Como $2^{-1}=4$ en $\mathbb{Z}_7$, multiplicamos por $4$:
\[
A = 4 \cdot 
\begin{pmatrix}
2 & 4\\[2pt]
6 & 1
\end{pmatrix}
=
\begin{pmatrix}
1 & 2\\[2pt]
3 & 4
\end{pmatrix}.
\]

Sustituyendo en la primera ecuación para obtener \(B\):
\[
\begin{pmatrix}1 & 2\\[2pt]3 & 4\end{pmatrix} + B =
\begin{pmatrix}4 & 3\\[2pt]2 & 6\end{pmatrix}
\implies
B =
\begin{pmatrix}4 & 3\\[2pt]2 & 6\end{pmatrix}
-
\begin{pmatrix}1 & 2\\[2pt]3 & 4\end{pmatrix}
=
\begin{pmatrix}3 & 1\\[2pt]6 & 2\end{pmatrix}.
\]

Calculamos los cuadrados:

\[
A^2 =
\begin{pmatrix}1 & 2\\[2pt]3 & 4\end{pmatrix}
\begin{pmatrix}1 & 2\\[2pt]3 & 4\end{pmatrix}
=
\begin{pmatrix}
1\cdot1+2\cdot3 & 1\cdot2+2\cdot4\\[2pt]
3\cdot1+4\cdot3 & 3\cdot2+4\cdot4
\end{pmatrix}
=
\begin{pmatrix}
1+6 & 2+8\\[2pt]
3+12 & 6+16
\end{pmatrix}
\equiv
\begin{pmatrix}
0 & 3\\[2pt]
1 & 1
\end{pmatrix}.
\]

\[
B^2 =
\begin{pmatrix}3 & 1\\[2pt]6 & 2\end{pmatrix}
\begin{pmatrix}3 & 1\\[2pt]6 & 2\end{pmatrix}
=
\begin{pmatrix}
3\cdot3+1\cdot6 & 3\cdot1+1\cdot2\\[2pt]
6\cdot3+2\cdot6 & 6\cdot1+2\cdot2
\end{pmatrix}
=
\begin{pmatrix}
9+6 & 3+2\\[2pt]
18+12 & 6+4
\end{pmatrix}
\equiv
\begin{pmatrix}
1 & 5\\[2pt]
2 & 3
\end{pmatrix}.
\]

Finalmente:
\[
A^2 - B^2 =
\begin{pmatrix}0 & 3\\[2pt]1 & 1\end{pmatrix}
-
\begin{pmatrix}1 & 5\\[2pt]2 & 3\end{pmatrix}
=
\boxed{\begin{pmatrix}6 & 5\\[2pt]6 & 5\end{pmatrix}}.
\]
\end{ejercicio}

\begin{ejercicio}[title=Ejercicio 2]
Sean las matrices 
\[
A = 
\begin{pmatrix}
1 & a\\
2 & b
\end{pmatrix},
\qquad
B =
\begin{pmatrix}
4 & 6\\
10 & 2
\end{pmatrix}
\]
con coeficientes en $\mathbb{Z}_{11}$.
Determine todos los pares ordenados $(a,b)$ para los cuales se verifica que $(A+B)\cdot(A-B)=A^2-B^2$.

\tcblower
\textbf{Resolución:}

\vspace{0.5em}

\textbf{Método 1. Construir las matrices correspondientes e igualarlas.}

\begin{itemize}
    \item Calculamos:
    \[
    A+B =
    \begin{pmatrix}
    5 & a+6\\
    1 & b+2
    \end{pmatrix},\qquad
    A-B =
    \begin{pmatrix}
    8 & a+5\\
    3 & b+9
    \end{pmatrix}
    \]
    \item Multiplicamos:
    \[
    (A+B)(A-B) =
    \begin{pmatrix}
    7+3a+7 & 5a+3+(a+b)(b+9)\\
    8+3b+6 & a+5+(b+2)(b+9)
    \end{pmatrix} =
    \]
    \[ =
    \begin{pmatrix}
    3a+3 & ab+3a+6b+4\\
    3b+3 & b^2+a+1
    \end{pmatrix}
    \]
    \item Calculamos $A^2$:
    \[
    A^2 = 
    \begin{pmatrix}
    1 & a\\
    2 & b
    \end{pmatrix}
    \begin{pmatrix}
    1 & a\\
    2 & b
    \end{pmatrix}
    =
    \begin{pmatrix}
    2a+1 & a+ab\\
    2b+2 & 2a+b^2
    \end{pmatrix}
    \]
    \item Calculamos $B^2$:
    \[
    B^2 =
    \begin{pmatrix}
    4 & 6\\
    10 & 2
    \end{pmatrix}
    \begin{pmatrix}
    4 & 6\\
    10 & 2
    \end{pmatrix}
    =
    \begin{pmatrix}
    10 & 3\\
    5 & 9
    \end{pmatrix}
    \]
    \item Restamos:
    \[
    A^2-B^2 =
    \begin{pmatrix}
    2a+1-10 & a+ab-3\\
    2b+2-5 & 2a+b^2-9
    \end{pmatrix}
    =
    \begin{pmatrix}
    2a+2 & ab+a+8\\
    2b+8 & b^2+2a+2
    \end{pmatrix}
    \]
    \item Igualamos $(A+B)(A-B) = A^2-B^2$:
    \[
    \begin{pmatrix}
    3a+3 & ab+3a+6b+4\\
    3b+3 & b^2+a+1
    \end{pmatrix} = 
    \begin{pmatrix}
    2a+2 & ab+a+8\\
    2b+8 & b^2+2a+2
    \end{pmatrix}
    \]
    \item Los elementos ${\cdot}_{11}$ han de ser iguales. Análogamente con los elementos ${\cdot}_{21}$:
    \[
    3a+3 = 2a+2 \implies \boxed{a=10}, \quad
    3b+3 = 2b+8 \implies \boxed{b=5}.
    \]
    \item Estos valores funcionan para el resto de las ecuaciones (se han omitido las comprobaciones).
\end{itemize}

\end{ejercicio}
\begin{ejercicio}[title=Ejercicio 2]

\textbf{Método 2. Usando $(A+B)(A-B) = A^2-B^2 \iff AB = BA$.}

\begin{itemize}
    \item Calculamos:
    \[
    AB = 
    \begin{pmatrix}
    1 & a\\ 2 & b
    \end{pmatrix}
    \begin{pmatrix}
    4 & 6\\ 10 & 2
    \end{pmatrix}
    =
    \begin{pmatrix}
    10a+4 & 2a+6\\
    10b+8 & 2b+1
    \end{pmatrix},
    \]
    \[
    BA = 
    \begin{pmatrix}
    4 & 6\\ 10 & 2
    \end{pmatrix}
    \begin{pmatrix}
    1 & a\\ 2 & b
    \end{pmatrix}
    =
    \begin{pmatrix}
    5 & 4a+6b\\
    3 & 10a+2b
    \end{pmatrix}
    \]
    \item Los elementos ${\cdot}_{11}$ han de ser iguales. Análogamente con los elementos ${\cdot}_{21}$:
    \[
    10a+4 = 5 \implies \boxed{a = 10}, \quad
    10b+8 = 3 \implies \boxed{b = 5}.
    \]
    \item Comprobamos la igualdad de los otros elementos:
    \[
    2(10)+6 = 4(10)+6(5) = 4. \text{ Sí}.
    \]\[
    2(5)+1 = 10(10)+2(5) = 0. \text{ Sí}.
    \]
\end{itemize}

\end{ejercicio}

\begin{ejercicio}[title=Ejercicio 3]
Dadas $A,B \in \mathfrak{M}_n(K)$, ambas simétricas, pruebe que $A \cdot B$ es simétrica si y sólo si $A \cdot B = B \cdot A$.

\tcblower
\textbf{Resolución:}

Debemos efectuar la demostración en ambos sentidos:

\begin{itemize}
    \item[$\boxed\Rightarrow$] Supongamos que $A\cdot B$ es simétrica.  
    Por definición, esto significa
    \[
    A\cdot B = (A\cdot B)^T.
    \] 
    Pero sabemos que por propiedades del producto de matrices se cumple
    \[
    (A\cdot B)^T = B^T \cdot A^T.
    \] 
    Como $A$ y $B$ son simétricas ($A^T = A$, $B^T = B$), obtenemos
    \[
    B^T \cdot A^T  = B \cdot A.
    \] 
    Por tanto, $A \cdot B = B \cdot A$.

    \item[$\boxed{\Leftarrow}$] Supongamos que $A \cdot B = B \cdot A$.  
    Aplicando la traspuesta a ambos lados:
    \[
    (A \cdot B)^T = (B \cdot A)^T.
    \] 
    Por propiedades del producto de matrices:
    \[
    (B \cdot A)^T = A^T \cdot B^T.
    \] 
    Como $A$ y $B$ son simétricas ($A^T = A$, $B^T = B$), obtenemos:
    \[
    (A \cdot B)^T = A \cdot B.
    \] 
    Por la definición, esto significa que $A \cdot B$ es simétrica.
\end{itemize}
\[
\boxed{A \cdot B \text{ simétrica } \quad \Longleftrightarrow \quad A \cdot B = B \cdot A}
\]

\end{ejercicio}

\begin{ejercicio}[title=Ejercicio 4]
Si 
\[
A = 
\begin{pmatrix}
ab & b^2\\
- a^2 & -ab
\end{pmatrix},
\]
compruebe que $A^2$ es la matriz nula.

\tcblower
\textbf{Resolución:}

Calculamos $A^2 = A \cdot A$:

\[
A^2 =
\begin{pmatrix}
ab & b^2\\
- a^2 & -ab
\end{pmatrix}
\begin{pmatrix}
ab & b^2\\
- a^2 & -ab
\end{pmatrix}
=\begin{pmatrix}
a^2 b^2 - a^2 b^2 & a b^3 - a b^3 \\
- a^3 b + a^3 b & - a^2 b^2 + a^2 b^2
\end{pmatrix}
=
\boxed{\begin{pmatrix}
0 & 0 \\
0 & 0
\end{pmatrix}}.
\]
\end{ejercicio}

\begin{ejercicio}[title=Ejercicio 5]
En cada apartado, resuelva el sistema propuesto en $\mathfrak{M}_2(\mathbb{Z}_7)$:

\[
\begin{tabular}{c @{\hspace{2cm}} c}
$a)\begin{cases}
4A + 5B =
\begin{pmatrix}
3 & 2\\
6 & 3
\end{pmatrix}\\[1em]
2A + 3B =
\begin{pmatrix}
1 & 2\\
0 & 0
\end{pmatrix}
\end{cases}$ &
$b)\begin{cases}
5A + 2B^T =
\begin{pmatrix}
5 & 2\\
2 & 3
\end{pmatrix}\\[1em]
3A^T + 5B =
\begin{pmatrix}
4 & 1\\
6 & 2
\end{pmatrix}
\end{cases}$
\end{tabular}
\]

\tcblower
\textbf{Resolución:}

\textbf{a)}

\vspace{0.5em}

Multiplicamos por 2 la segunda ecuación:

\[
\begin{cases}
    4A + 5B = 
    \begin{pmatrix}
        3 & 2 \\
        6 & 3
    \end{pmatrix} \\
    4A + 6B = 
    \begin{pmatrix}
        2 & 4 \\
        0 & 0
    \end{pmatrix}
\end{cases}
\]

Restando la segunda a la primera:

\[
-B = 
\begin{pmatrix}
1 & -2\\
6 & 3
\end{pmatrix}
\implies B = 
\begin{pmatrix}
-1 & 2\\
-6 & -3
\end{pmatrix}
=
\begin{pmatrix}
6 & 2\\
1 & 4
\end{pmatrix}.
\]

Sustituyendo en la segunda ecuación:

\[
2A + 3B = 
\begin{pmatrix}
1 & 2\\
0 & 0
\end{pmatrix} 
\implies
2A = 
\begin{pmatrix}
1 & 2\\
0 & 0
\end{pmatrix} - 3B 
=
\]
\[=
\begin{pmatrix}
1 & 2\\
0 & 0
\end{pmatrix} - 
3 \begin{pmatrix}
6 & 2\\
1 & 4
\end{pmatrix}
=
\begin{pmatrix}
1 & 2\\
0 & 0
\end{pmatrix} - 
\begin{pmatrix}
4 & 6\\
3 & 5
\end{pmatrix}
=
\begin{pmatrix}
4 & 3\\
4 & 2
\end{pmatrix}.
\]

\vspace{0.5em}

Multiplicando por $2^{-1} = 4$ en $\mathbb{Z}_7$:

\[
A = 4 \cdot 
\begin{pmatrix}
4 & 3\\
4 & 2
\end{pmatrix} =
\begin{pmatrix}
2 & 5\\
2 & 1
\end{pmatrix}.
\]

\[
\boxed{A = \begin{pmatrix}2 & 5\\ 2 & 1\end{pmatrix}}, \quad 
\boxed{B = \begin{pmatrix}6 & 2\\ 1 & 4\end{pmatrix}}.
\]

\textbf{b)}

\vspace{0.5em}

Por propiedades de la traspuesta y la suma, aplicando la traspuesta a la segunda ecuación original obtenemos un sistema equivalente con \(B^T\):

\[
(3A^T + 5B)^T = 
\begin{pmatrix}
4 & 1\\
6 & 2
\end{pmatrix}^T
\implies
3A + 5B^T =
\begin{pmatrix}
4 & 6\\
1 & 2
\end{pmatrix}.
\]

\end{ejercicio}
\begin{ejercicio}[title=Ejercicio 5]

Así, el sistema equivalente en las incógnitas \(A\) y \(B^T\) es

\[
\begin{cases}
5A + 2B^T =
\begin{pmatrix}
5 & 2\\
2 & 3
\end{pmatrix},\\[4pt]
3A + 5B^T =
\begin{pmatrix}
4 & 6\\
1 & 2
\end{pmatrix}.
\end{cases}
\]

Sumamos ambas ecuaciones en $\mathbb{Z}_7$ y queda:
\[
A = \begin{pmatrix}2 & 1\\[2pt] 3 & 5\end{pmatrix}
\]

Calculamos la traspuesta:
\[
A^T = \begin{pmatrix}2 & 3\\[2pt] 1 & 5\end{pmatrix}.
\]

Para poder sustituirla en la segunda ecuación original:
\[
3A^T + 5B = \begin{pmatrix}4 & 1\\[2pt] 6 & 2\end{pmatrix}
\quad\Rightarrow\quad
3\begin{pmatrix}2 & 3\\[2pt] 1 & 5\end{pmatrix} + 5B =
\begin{pmatrix}4 & 1\\[2pt] 6 & 2\end{pmatrix}.
\]

Calculamos \(3A^T\) en \(\mathbb{Z}_7\):
\[
3\begin{pmatrix}2 & 3\\[2pt] 1 & 5\end{pmatrix} =
\begin{pmatrix}6 & 2\\[2pt] 3 & 1\end{pmatrix}.
\]

Despejamos \(5B\):
\[
5B =
\begin{pmatrix}4 & 1\\[2pt] 6 & 2\end{pmatrix}
-
\begin{pmatrix}6 & 2\\[2pt] 3 & 1\end{pmatrix}
=
\begin{pmatrix}5 & 6\\[2pt] 3 & 1\end{pmatrix}.
\]

Como \(5^{-1}=3\) en \(\mathbb{Z}_7\), multiplicamos por \(3\) para obtener \(B\):
\[
B = 3 \begin{pmatrix}5 & 6\\[2pt] 3 & 1\end{pmatrix}
=
\begin{pmatrix}1 & 4\\[2pt] 2 & 3\end{pmatrix}.
\]

\[
\boxed{A = \begin{pmatrix}2 & 1\\ 3 & 5\end{pmatrix}}, \qquad
\boxed{B = \begin{pmatrix}1 & 4\\ 2 & 3\end{pmatrix}}.
\]

\end{ejercicio}

\begin{ejercicio}[title=Ejercicio 6]
Demuestre que toda matriz de $\mathfrak{M}_n(\mathbb{R})$ se expresa de manera única como suma de una matriz simétrica y una matriz antisimétrica.  
Aplique lo que acaba de probar a la matriz siguiente:
\[
B =
\begin{pmatrix}
1 & 7\\
2 & -4
\end{pmatrix}.
\]

\tcblower
\textbf{Resolución:}

Sea \(M\in\mathfrak{M}_n(\mathbb{R})\). Queremos demostrar que existen \(S,A\) tales que
\[
M = S + A,
\]
con \(S\) simétrica y \(A\) antisimétrica.

Trasponemos ambos miembros. Usamos propidades de la suma y la definición de simétrica ($S=S^T$) y antisimétrica ($A=-A^T$):
\[
M^T = (S+A)^T = S^T + A^T = S - A.
\]

Por tanto obtenemos el sistema
\[
\begin{cases}
M = S + A\\[4pt]
M^T = S - A
\end{cases}
\]

Sumando y restando estas dos ecuaciones despejamos:
\[
M + M^T = 2S \quad\Longrightarrow\quad S = \frac{M+M^T}{2},
\]
\[
M - M^T = 2A \quad\Longrightarrow\quad A = \frac{M-M^T}{2}.
\]

Si S es simétrica, entonces por definición \(S = S^T\). Igualando y aplicando propiedades:
\[
\frac{M+M^T}{2}
= \left(\frac{M+M^T}{2}\right)^T
= \frac{(M+M^T)^T}{2}
= \frac{M^T+M}{2}
= \frac{M+M^T}{2}
\]
\[
\implies
\text{S es simétrica.}
\]

Si A es antisimétrica, entonces por definición \(A = -A^T\). Igualando y aplicando propiedades:
\[
\frac{M-M^T}{2}
= -\left(\frac{M-M^T}{2}\right)^T
= -\frac{(M-M^T)^T}{2}
= -\frac{M^T-M}{2}
= \frac{M-M^T}{2}
\]
\[
\implies
\text{A es antisimétrica.}
\]

Entonces
\[
\boxed{\forall M \in \mathfrak{M}_n(\mathbb{R})\quad M = S+A, \qquad S\text{ simétrica,} \quad A\text{ antisimétrica.}}
\]
\end{ejercicio}
\begin{ejercicio}[title=Ejercicio 6]
\textbf{Aplicación al caso concreto.}
\[
B=
\begin{pmatrix}
1 & 7\\
2 & -4
\end{pmatrix}, \qquad
B^T=
\begin{pmatrix}
1 & 2\\
7 & -4
\end{pmatrix}.
\]

Calculamos la parte simétrica:
\[
S=\frac{B+B^T}{2}=\frac{1}{2}
\begin{pmatrix}
1+1 & 7+2\\[4pt]
2+7 & -4-4
\end{pmatrix}
=\frac{1}{2}
\begin{pmatrix}
2 & 9\\[4pt]
9 & -8
\end{pmatrix}
=
\begin{pmatrix}
1 & \tfrac{9}{2}\\[4pt]
\tfrac{9}{2} & -4
\end{pmatrix}.
\]

Calculamos la parte antisimétrica:
\[
A=\frac{B-B^T}{2}=\frac{1}{2}
\begin{pmatrix}
1-1 & 7-2\\[4pt]
2-7 & -4+4
\end{pmatrix}
=\frac{1}{2}
\begin{pmatrix}
0 & 5\\[4pt]
-5 & 0
\end{pmatrix}
=
\begin{pmatrix}
0 & \tfrac{5}{2}\\[4pt]
-\tfrac{5}{2} & 0
\end{pmatrix}.
\]

Comprobación:
\[
S+A =
\begin{pmatrix}
1 & \tfrac{9}{2}\\[4pt]
\tfrac{9}{2} & -4
\end{pmatrix}
+
\begin{pmatrix}
0 & \tfrac{5}{2}\\[4pt]
-\tfrac{5}{2} & 0
\end{pmatrix}
=
\begin{pmatrix}
1 & 7\\[4pt]
2 & -4
\end{pmatrix}=B.
\]

\[
\boxed{B = S+A, \qquad S\text{ simétrica,} \quad A\text{ antisimétrica.}}
\]
\end{ejercicio}

\section{Operaciones o Transformaciones Elementales sobre Matrices}

\begin{ejercicio}[title=Ejercicio 8]
Calcule el rango de la matriz siguiente con coeficientes en $\mathbb{R}$:
\[
A =
\begin{pmatrix}
1 & 2 & -1 & 2\\
1 & 3 & 2 & -1\\
1 & 1 & -4 & 5\\
1 & 4 & 5 & -4
\end{pmatrix}.
\]

\tcblower
\textbf{Resolución:}

Aplicamos operaciones elementales de fila para obtener una matriz en forma escalonada por filas y observar el número de filas no nulas, que es el rango.

\[
A =
\begin{pmatrix}
1 & 2 & -1 & 2\\
1 & 3 & 2 & -1\\
1 & 1 & -4 & 5\\
1 & 4 & 5 & -4
\end{pmatrix}
\overset{\substack{F_4 - F_1 \\ F_3 - F_1 \\ F_2 - F_1 \\{}}}{\sim_f}
\begin{pmatrix}
1 & 2 & -1 & 2\\
0 & -1 & -3 & 3\\
0 & 1 & 3 & -3\\
0 & -2 & -6 & 6
\end{pmatrix}
\overset{\substack{F_4 - 2F_2 \\ F_3 + 2F_2 \\{}}}{\sim_f}
\begin{pmatrix}
1 & 2 & -1 & 2\\
0 & -1 & -3 & 3\\
0 & -1 & -3 & 3\\
0 & 0 & 0 & 0
\end{pmatrix}
\]

\[
\overset{\substack{F_3 - F_2 \\{}}}{\sim_f}
\begin{pmatrix}
1 & 2 & -1 & 2\\
0 & -1 & -3 & 3\\
0 & 0 & 0 & 0\\
0 & 0 & 0 & 0
\end{pmatrix}
\text{ está en forma escalonada por filas, entonces:}
\]

\[
\operatorname{Rg}(A) = \text{nº filas no nulas} = \boxed{2}.
\]
\end{ejercicio}

\begin{ejercicio}[title=Ejercicio 9]
Calcule el rango de la matriz siguiente con coeficientes en $\mathbb{Z}_7$:
\[
A =
\begin{pmatrix}
3 & 2 & 5 & 4 & 1\\
2 & 5 & 6 & 1 & 4\\
6 & 3 & 1 & 4 & 3\\
1 & 0 & 5 & 1 & 1\\
5 & 3 & 3 & 3 & 2
\end{pmatrix}.
\]

\tcblower
\textbf{Resolución:}

Aplicamos operaciones elementales de fila :

\[
A =
\begin{pmatrix}
3 & 2 & 5 & 4 & 1\\
2 & 5 & 6 & 1 & 4\\
6 & 3 & 1 & 4 & 3\\
1 & 0 & 5 & 1 & 1\\
5 & 3 & 3 & 3 & 2
\end{pmatrix}
\overset{\substack{5F_1\\4F_2\\6F_3\\3F_5\\{}}}{\sim_f}
\begin{pmatrix}
1 & 3 & 4 & 6 & 5\\
1 & 6 & 3 & 4 & 2\\
1 & 4 & 6 & 3 & 4\\
1 & 0 & 5 & 1 & 1\\
1 & 2 & 2 & 2 & 6
\end{pmatrix}
\overset{\substack{F_2 - F_1\\F_3 - F_1\\F_4 - F_1\\F_5 - F_1\\{}}}{\sim_f}
\begin{pmatrix}
1 & 3 & 4 & 6 & 5\\
0 & 3 & 6 & 5 & 4\\
0 & 1 & 2 & 4 & 6\\
0 & 4 & 1 & 2 & 3\\
0 & 6 & 5 & 3 & 1
\end{pmatrix}
\]

\[
\overset{\substack{F_3 + 2F_2\\F_4 + F_2\\F_5 - 2F_2\\{}}}{\sim_f}
\begin{pmatrix}
1 & 3 & 4 & 6 & 5\\
0 & 3 & 6 & 5 & 4\\
0 & 0 & 0 & 0 & 0\\
0 & 0 & 0 & 0 & 0\\
0 & 0 & 0 & 0 & 0
\end{pmatrix}
\text{ está en forma escalonada por filas, entonces:}
\]

\[
\operatorname{Rg}(A) = \text{nº filas no nulas} = \boxed{2}.
\]

\end{ejercicio}

\begin{ejercicio}[title=Ejercicio 10]
La forma normal de Hermite por filas de la matriz 
\[
A =
\begin{pmatrix}
3 & 1 & 1 & 2\\
1 & 4 & 2 & 1\\
5 & 2 & 1 & 3
\end{pmatrix}
\in \mathfrak{M}_{3\times 4}(\mathbb{Z}_7)
\]
es:

\[
\begin{tabular}{clclclcl}
(a) &
$\begin{pmatrix}
1 & 0 & 0 & 6\\
0 & 1 & 0 & 3\\
0 & 0 & 1 & 2
\end{pmatrix}$ &
(b) &
$\begin{pmatrix}
1 & 0 & 0 & 2\\
0 & 1 & 0 & 5\\
0 & 0 & 1 & 2
\end{pmatrix}$ &
(c) &
$\begin{pmatrix}
1 & 0 & 1 & 1\\
0 & 1 & 3 & 2\\
0 & 0 & 0 & 0
\end{pmatrix}$ &
(d) &
$\begin{pmatrix}
1 & 0 & 1 & 0\\
0 & 1 & 3 & 0\\
0 & 0 & 0 & 1
\end{pmatrix}$
\end{tabular}
\]

\tcblower
\textbf{Resolución:}

Aplicamos operaciones elementales de fila (en $\mathbb{Z}_7$):

\[
A =
\begin{pmatrix}
3 & 1 & 1 & 2\\
1 & 4 & 2 & 1\\
5 & 2 & 1 & 3
\end{pmatrix}
\overset{\substack{5F_1 \\ 3F_3 \\{}}}{\sim_f}
\begin{pmatrix}
1 & 5 & 5 & 3\\
1 & 4 & 2 & 1\\
1 & 6 & 3 & 2
\end{pmatrix}
\overset{\substack{F_2 - F_1 \\ F_3 - F_1 \\{}}}{\sim_f}
\begin{pmatrix}
1 & 5 & 5 & 3\\
0 & 6 & 4 & 5\\
0 & 1 & 5 & 6
\end{pmatrix}
\overset{\substack{F_2 + F_3 \\{}}}{\sim_f}
\begin{pmatrix}
1 & 5 & 5 & 3\\
0 & 0 & 2 & 4\\
0 & 1 & 5 & 6
\end{pmatrix}
\overset{\substack{4F_2 \\{}}}{\sim_f}
\]

\[
\begin{pmatrix}
1 & 5 & 5 & 3\\
0 & 0 & 1 & 2\\
0 & 1 & 5 & 6
\end{pmatrix}
\overset{\substack{F_2 \leftrightarrow F_3 \\{}}}{\sim_f}
\begin{pmatrix}
1 & 5 & 5 & 3\\
0 & 1 & 5 & 6\\
0 & 0 & 1 & 2
\end{pmatrix}
\overset{\substack{F_1 - 5F_3 \\ F_2 - 5F_3 \\{}}}{\sim_f}
\begin{pmatrix}
1 & 5 & 0 & 0\\
0 & 1 & 0 & 3\\
0 & 0 & 1 & 2
\end{pmatrix}
\overset{\substack{F_1 - 5F_2 \\{}}}{\sim_f}
\begin{pmatrix}
1 & 0 & 0 & 6\\
0 & 1 & 0 & 3\\
0 & 0 & 1 & 2
\end{pmatrix}
= H_f(A)
\]

Por tanto la forma normal de Hermite por filas de \(A\) es

\[
\boxed{\text{a)}\;
\begin{pmatrix}
1 & 0 & 0 & 6\\
0 & 1 & 0 & 3\\
0 & 0 & 1 & 2
\end{pmatrix}
\;}
\]
\end{ejercicio}

\begin{ejercicio}[title=Ejercicio 11]
Estudie cuáles de las matrices siguientes con coeficientes en $\mathbb{Z}_7$ son equivalentes por filas:

\[
A =
\begin{pmatrix}
1 & 2 & 3 & 3\\
3 & 4 & 1 & 2\\
1 & 1 & 6 & 5
\end{pmatrix}, \quad
B =
\begin{pmatrix}
3 & 1 & 0 & 2\\
2 & 4 & 2 & 1\\
6 & 1 & 5 & 5
\end{pmatrix}, \quad
C =
\begin{pmatrix}
3 & 1 & 3 & 2\\
2 & 4 & 6 & 1\\
6 & 1 & 2 & 5
\end{pmatrix}.
\]

\tcblower
\textbf{Resolución:}

Sabemos que dos matrices $A$ y $B$ serán equivalentes por filas si su forma normal de Hermite por filas es la misma, es decir, $H_f(A)=H_f(B)$.

\vspace{0.5em}

Calculamos $H_f(A)$:
\[
A =
\begin{pmatrix}
1 & 2 & 3 & 3\\
3 & 4 & 1 & 2\\
1 & 1 & 6 & 5
\end{pmatrix}
\overset{\substack{F_2 - 3F_1 \\ F_3 - F_1 \\ {}}}{\sim_f}
\begin{pmatrix}
1 & 2 & 3 & 3\\
0 & 5 & 6 & 0\\
0 & 6 & 3 & 2
\end{pmatrix}
\overset{\substack{3F_2 \\ {}}}{\sim_f}
\begin{pmatrix}
1 & 2 & 3 & 3\\
0 & 1 & 4 & 0\\
0 & 6 & 3 & 2
\end{pmatrix}
\overset{\substack{F_1 - 2F_2 \\ F_3 + F_2 \\ {}}}{\sim_f}
\begin{pmatrix}
1 & 0 & 2 & 3\\
0 & 1 & 4 & 0\\
0 & 0 & 0 & 2
\end{pmatrix}
\]
\[
\overset{\substack{4F_3 \\ {}}}{\sim_f}
\begin{pmatrix}
1 & 0 & 2 & 3\\
0 & 1 & 4 & 0\\
0 & 0 & 0 & 1
\end{pmatrix}
\overset{\substack{F_1 - 3F_3 \\ {}}}{\sim_f}
\begin{pmatrix}
1 & 0 & 2 & 0\\
0 & 1 & 4 & 0\\
0 & 0 & 0 & 1
\end{pmatrix}
= H_f(A).
\]

\vspace{0.5em}

Calculamos $H_f(B)$:
\[
B =
\begin{pmatrix}
3 & 1 & 0 & 2\\
2 & 4 & 2 & 1\\
6 & 1 & 5 & 5
\end{pmatrix}
\overset{\substack{5F_1 \\ 4F_2 \\ 6F_3 \\ {}}}{\sim_f}
\begin{pmatrix}
1 & 5 & 0 & 3\\
1 & 2 & 1 & 4\\
1 & 6 & 2 & 2
\end{pmatrix}
\overset{\substack{F_2-F_1 \\ F_3-F_1 \\ {}}}{\sim_f}
\begin{pmatrix}
1 & 5 & 0 & 3\\
0 & 4 & 1 & 1\\
0 & 1 & 2 & 6
\end{pmatrix}
\overset{\substack{2F_2 \\ {}}}{\sim_f}
\begin{pmatrix}
1 & 5 & 0 & 3\\
0 & 1 & 2 & 2\\
0 & 1 & 2 & 6
\end{pmatrix}
\]
\[
\overset{\substack{F_3 - F_2 \\ {}}}{\sim_f}
\begin{pmatrix}
1 & 5 & 0 & 3\\
0 & 1 & 2 & 2\\
0 & 0 & 0 & 4
\end{pmatrix}
\overset{\substack{2F_3 \\ F_1 - 5F_2 \\ {}}}{\sim_f}
\begin{pmatrix}
1 & 0 & 4 & 0\\
0 & 1 & 2 & 2\\
0 & 0 & 0 & 1
\end{pmatrix}
\overset{\substack{F_2 - 2F_3 \\ {}}}{\sim_f}
\begin{pmatrix}
1 & 0 & 4 & 0\\
0 & 1 & 2 & 0\\
0 & 0 & 0 & 1
\end{pmatrix}
= H_f(B).
\]

\vspace{0.5em}

Calculamos $H_f(C)$:
\[
C =
\begin{pmatrix}
3 & 1 & 3 & 2\\
2 & 4 & 6 & 1\\
6 & 1 & 2 & 5
\end{pmatrix}
\overset{\substack{5F_1 \\ 4F_2 \\ -F_3 \\ {}}}{\sim_f}
\begin{pmatrix}
1 & 5 & 1 & 3\\
1 & 2 & 3 & 4\\
1 & 6 & 5 & 2
\end{pmatrix}
\overset{\substack{F_2-F_1 \\ F_3-F_1 \\ {}}}{\sim_f}
\begin{pmatrix}
1 & 5 & 1 & 3\\
0 & 4 & 2 & 1\\
0 & 1 & 4 & 6
\end{pmatrix}
\overset{\substack{2F_2 \\ {}}}{\sim_f}
\begin{pmatrix}
1 & 5 & 1 & 3\\
0 & 1 & 4 & 2\\
0 & 1 & 4 & 6
\end{pmatrix}
\]
\[
\overset{F_3-F_2}{\sim_f}
\begin{pmatrix}
1 & 5 & 1 & 3\\
0 & 1 & 4 & 2\\
0 & 0 & 0 & 4
\end{pmatrix}
\overset{\substack{2F_3 \\ F_1 - 5F_2 \\ {}}}{\sim_f}
\begin{pmatrix}
1 & 0 & 2 & 0\\
0 & 1 & 4 & 2\\
0 & 0 & 0 & 1
\end{pmatrix}
\overset{\substack{F_2 - 2F_3 \\ {}}}{\sim_f}
\begin{pmatrix}
1 & 0 & 2 & 0\\
0 & 1 & 4 & 0\\
0 & 0 & 0 & 1
\end{pmatrix}
= H_f(C).
\]

\vspace{0.5em}

Observamos que $H_f(A) = H_f(C)$, entonces \boxed{\text{$A$ y $C$ son equivalentes por filas}}.

\end{ejercicio}

\begin{ejercicio}[title=Ejercicio 12]
Calcule la forma normal de Hermite por columnas de la matriz siguiente con coeficientes en $\mathbb{Z}_7$:
\[
A =
\begin{pmatrix}
1 & 2 & 3 & 4\\
5 & 3 & 1 & 6\\
3 & 4 & 2 & 3
\end{pmatrix}.
\]

\tcblower
\textbf{Resolución:}

\[
A =
\begin{pmatrix}
1 & 2 & 3 & 4\\
5 & 3 & 1 & 6\\
3 & 4 & 2 & 3
\end{pmatrix}
\overset{\substack{4C_2 \\ 5C_3 \\ 2C_4 \\ {}}}{\sim_c}
\begin{pmatrix}
1 & 1 & 1 & 1\\
5 & 5 & 5 & 5\\
3 & 2 & 3 & 6
\end{pmatrix}
\overset{\substack{C_4 - C_1 \\ C_3 - C_1 \\ C_2 - C_1 \\ {}}}{\sim_c}
\begin{pmatrix}
1 & 0 & 0 & 0\\
5 & 0 & 0 & 0\\
3 & 6 & 0 & 3
\end{pmatrix}
\]

\[
\overset{\substack{-C_2 \\ 5C_4\\ {}}}{\sim_c}
\begin{pmatrix}
1 & 0 & 0 & 0\\
5 & 0 & 0 & 0\\
3 & 1 & 0 & 1
\end{pmatrix}
\overset{\substack{C_4 - C_2 \\ C_1 - 3C_2 \\ {}}}{\sim_c}
\begin{pmatrix}
1 & 0 & 0 & 0\\
5 & 0 & 0 & 0\\
0 & 1 & 0 & 0
\end{pmatrix}
= H_c(A).
\]

Entonces:

\[
\boxed{
H_c(A) =
\begin{pmatrix}
1 & 0 & 0 & 0\\
5 & 0 & 0 & 0\\
0 & 1 & 0 & 0
\end{pmatrix}
}
\]
\end{ejercicio}

\section{Sistemas de Ecuaciones Lineales}

\subsection{Teorema de Rouché--Frobenius}

\begin{ejercicio}[title=Ejercicio 27]
Acerca del sistema siguiente con coeficientes en $\mathbb{R}$:
\[
\begin{cases}
x + ay - z = 1\\
x + y + z = a
\end{cases}
\]
podemos afirmar que:

\begin{enumerate}
\item[a)] Independientemente del valor de $a$, es compatible determinado.
\item[b)] Independientemente del valor de $a$, es compatible indeterminado.
\item[c)] Es siempre incompatible.
\item[d)] La compatibilidad o incompatibilidad depende del valor de $a$.
\end{enumerate}

\tcblower
\textbf{Resolución:}

Usamos el teorema de Rouché–Frobenius para averiguar la compatibilidad del sistema. Esto consiste en analizar los valores de los rangos de $A$ y $(A\,|\,B)$. Calculamos una matriz en forma escalonada por filas para averiguar los rangos:

\[
(A\,|\,B)=
\begin{pmatrix}
1 & a & -1 &\Big|& 1\\
1 & 1 &  1 &\Big|& a
\end{pmatrix}
\overset{\substack{F_2 - F_1\\{}}}{\sim_f}
\begin{pmatrix}
1 & a & -1 &\Big|& 1\\
0 & 1-a & 2 &\Big|& a-1
\end{pmatrix} \quad \text{\shortstack{ está en forma\\escalonada por filas.}}
\]
\[
\Rightarrow \operatorname{Rg}(A)=\text{nº filas no nulas en }A = 2
\]
porque, sea cual sea el valor de $a$, las filas no pueden ser nulas por ya haber elementos no nulos en ellas. 
De la misma forma,
\[
\Rightarrow \operatorname{Rg}(A\,|\,B)=\text{nº filas no nulas en }(A\,|\,B)=2.
\]

Como
\[
\operatorname{Rg}(A)=\operatorname{Rg}(A\,|\,B)=2 \neq 3=\text{nº incógnitas},
\]
por el teorema de Rouché-Frobenius, el sistema es SCI sea cual sea el valor de \(a\).

\[
\boxed{\text{b) Independientemente del valor de }a,\ \text{es compatible indeterminado.}}
\]
\end{ejercicio}

\begin{ejercicio}[title=Ejercicio 28]
Sea el sistema siguiente de ecuaciones lineales en $\mathbb{R}$:
\[
\begin{cases}
x + ay + 3az = a + 3\\
ay + 2az = a + 1
\end{cases}
\]
Entonces:

\begin{enumerate}
\item[a)] El sistema es compatible determinado independientemente de $a$.
\item[b)] El sistema es compatible indeterminado independientemente de $a$.
\item[c)] El sistema es incompatible independientemente de $a$.
\item[d)] La compatibilidad del sistema depende de $a$.
\end{enumerate}

\tcblower
\textbf{Resolución:}

Usamos el teorema de Rouché–Frobenius para averiguar la compatibilidad del sistema. Esto consiste en analizar los valores de los rangos de $A$ y $(A\,|\,B)$. Calculamos una matriz en forma escalonada por filas para averiguar los rangos. $(A\,|\,B)$ ya está en forma escalonada por filas, pero aplicamos una operación elemental de filas para simplificar la discriminación de valores de $a$:

\[
(A\,|\,B)=
\begin{pmatrix}
1 & a & 3a &\Big|& a+3\\
0 & a & 2a &\Big|& a+1
\end{pmatrix}
\overset{\substack{F_1 - F_2\\{}}}{\sim_f}
\begin{pmatrix}
1 & 0 & a &\Big|& 2\\
0 & a & 2a &\Big|& a+1
\end{pmatrix}
\]

\vspace{0.5em}

Analizamos rangos en función de \(a\).

\vspace{0.5em}

Para \(A\):
\begin{itemize}
\item La primera fila de \(A\) nunca puede ser nula (tiene un 1).
\item La segunda fila de \(A\) es \((0,a,2a)\), que es nula si y sólo si \(a=0\).
\end{itemize}

Por tanto
\[
\operatorname{Rg}(A)=
\begin{cases}
1 & \text{si } a=0,\\[4pt]
2 & \text{si } a\neq 0,
\end{cases}
\qquad \neq \text{nº incógnitas} = 3.
\]

\vspace{0.5em}

Para \((A\,|\,B)\):

\begin{itemize}
\item La primera fila nunca puede ser nula (contiene un 1 y también un 2).
\item La segunda fila de \((A\,|\,B)\) es \((0,a,2a\,|\,a+1)\). Para que esa fila sea nula necesitaríamos simultáneamente \(a=0\), \(2a=0\) y \(a+1=0\), es decir \(a=0\) y \(a=-1\) a la vez, lo cual es imposible. Por tanto la segunda fila nunca es nula.
\end{itemize}

Así
\[
\operatorname{Rg}(A\,|\,B)=2 \neq \text{nº incógnitas} = 3 \quad\text{para todo }a.
\]
\end{ejercicio}
\begin{ejercicio}[title=Ejercicio 28]
Comparando rangos:
\[
\begin{cases}
\operatorname{Rg}(A)=1\ne\operatorname{Rg}(A\,|\,B)= 2 \quad &\text{si } a=0,\\[4pt]
\operatorname{Rg}(A)= \operatorname{Rg}(A\,|\,B) = 2 \neq \text{nº incógnitas} = 3\quad &\text{si } a\neq 0.
\end{cases}
\]

Por el teorema de Rouché–Frobenius eso significa que:
\[
\text{El sistema es }
\begin{cases}
\text{SI} \quad &\text{si } a=0,\\[4pt]
\text{SCI}\quad &\text{si } a\neq 0.
\end{cases}
\]

Por tanto la compatibilidad depende de \(a\). La respuesta es:
\[
\boxed{\text{d) La compatibilidad del sistema depende de }a.}
\]
\end{ejercicio}

\begin{ejercicio}[title=Ejercicio 29]
Para el sistema de ecuaciones siguiente con coeficientes en $\mathbb{Z}_{11}$:
\[
\begin{cases}
3x + y + a z = 1\\
2x + 4y + (8a + 1)z = 2\\
x + 8y + (4a + 10)z = 10
\end{cases}
\]
se verifica que:

\begin{enumerate}
\item[a)] Es incompatible para cualquier valor de $a$.
\item[b)] Hay valores de $a$ para los cuales es compatible y valores de $a$ para los cuales es incompatible.
\item[c)] Es compatible indeterminado para cualquier valor de $a$.
\item[d)] Es compatible para cualquier valor de $a$. Además, hay valores de $a$ para los cuales es compatible determinado y otros para los cuales es compatible indeterminado.
\end{enumerate}

\tcblower
\textbf{Resolución:}

Usamos el teorema de Rouché–Frobenius para averiguar la compatibilidad del sistema. Esto consiste en analizar los valores de los rangos de $A$ y $(A\,|\,B)$. Calculamos una matriz en forma escalonada por filas.

\[
(A\,|\,B)=
\begin{pmatrix}
3 & 1 & a &\Big|& 1\\
2 & 4 & 8a+1 &\Big|& 2\\
1 & 8 & 4a+10 &\Big|& 10
\end{pmatrix}
\overset{\substack{F_2+3F_1\\F_3+7F_1}}{\sim_f}
\begin{pmatrix}
3 & 1 & a &\Big|& 1\\
0 & 7 & 1 &\Big|& 5\\
0 & 4 & 10 &\Big|& 6
\end{pmatrix}
\overset{\substack{F_3+F_2\\{}}}{\sim_f}
\begin{pmatrix}
3 & 1 & a &\Big|& 1\\
0 & 7 & 1 &\Big|& 5\\
0 & 0 & 0 &\Big|& 0
\end{pmatrix},
\]
que está en forma escalonada por filas.

\vspace{1em}

Tanto para $A$ como para $(A\,|\,B)$:
\begin{itemize}
    \item La primera fila no puede ser nula porque tiene elementos no nulos.
    \item La segunda fila es no nula.
    \item La tercera fila es nula.
\end{itemize}

Por tanto
\[
\operatorname{Rg}(A)=\operatorname{Rg}(A\,|\,B)=2 \neq 3 = \text{ nº de incógnitas.}
\]

Por el teorema de Rouché–Frobenius, es un SCI sea cual sea el valor de $a$.

\[
\boxed{\text{c) Es compatible indeterminado para cualquier valor de }a.}
\]
\end{ejercicio}

\begin{ejercicio}[title=Ejercicio 37 (⋆)]
Sean $\mathbb{K}$ un cuerpo, $n \in \mathbb{N}$, $x_0,x_1,\ldots,x_n \in \mathbb{K}$ distintos entre sí, e $y_0,y_1,\ldots,y_n \in \mathbb{K}$ no necesariamente distintos.  
Demuestre que existe un único polinomio $f(x) \in \mathbb{K}[x]$ tal que $\operatorname{grado}(f(x)) \le n$ y $f(x_i) = y_i$ para todo $i \in \{0,1,\ldots,n\}$.  
El único polinomio $f(x)$ que verifica las condiciones anteriores se denomina \textit{polinomio de interpolación}.

\tcblower
\textbf{Resolución:}

Queremos encontrar un polinomio 
\[
f(x) = a_0 + a_1 x + a_2 x^2 + \cdots + a_n x^n
\]
de grado como máximo $n$ que pase por los puntos $(x_i, y_i)$, es decir, que cumpla $f(x_i) = y_i$ para $i=0,1,\dots,n$. Es decir,
\[
\begin{cases}
f(x_0) = y_0\\
f(x_1) = y_1\\
\vdots\\
f(x_n) = y_n.
\end{cases}
\]

Esto nos da un sistema de $n+1$ ecuaciones lineales con $n+1$ incógnitas $(a_0, a_1, \dots, a_n)$:
\[
\begin{cases}
a_0 + a_1 x_0 + a_2 x_0^2 + \cdots + a_n x_0^n = y_0\\
a_0 + a_1 x_1 + a_2 x_1^2 + \cdots + a_n x_1^n = y_1\\
\vdots\\
a_0 + a_1 x_n + a_2 x_n^2 + \cdots + a_n x_n^n = y_n
\end{cases}
\]

La matriz de coeficientes de este sistema es:
\[
A =
\begin{pmatrix}
1 & x_0 & x_0^2 & \cdots & x_0^n \\
1 & x_1 & x_1^2 & \cdots & x_1^n \\
\vdots & \vdots & \vdots & \ddots & \vdots \\
1 & x_n & x_n^2 & \cdots & x_n^n
\end{pmatrix}.
\]

Su determinante es, por propiedades de los determinantes, el determinante de la traspuesta. Pero el determinante de la traspuesta es el de \textit{Vandermonde}. Por el ejercicio 34, sabemos que entonces vale:
\[
|A| = |A^T| = \prod_{0 \le j < k \le n} (x_j - x_k).
\]
Como los $x_i$ son todos distintos, cada factor $x_j - x_k$ es distinto de cero, y por tanto:
\[
\scalebox{1}{$
|A| \neq 0 \quad \Rightarrow \quad \text{SCD} \quad \Rightarrow \quad \text{1 solución}
$}
\]
En otras palabras, hay un único conjunto de coeficientes $(a_0, a_1, \dots, a_n)$ que hace que el polinomio pase por todos los puntos dados. Por lo tanto:

\[
\boxed{\exists f(x) \in \mathbb{K}[x]: \operatorname{grado}(f(x)) \le n,\,\forall i \in \{0,1,\ldots,n\}\,\, f(x_i) = y_i}
\]

\end{ejercicio}

\begin{ejercicio}[title=Ejercicio 38]
Encuentre el único polinomio $f(x) \in \mathbb{Z}_7[x]$ de grado menor o igual que 2 que verifica $f(1)=2$, $f(3)=3$, $f(6)=5$.

\tcblower
\textbf{Resolución:}

Por el Ejercicio 37, sabemos que existe un único polinomio que cumple las condiciones dadas.

Como $\operatorname{grado}(f(x))\le 2$, escribimos
\[
f(x)=ax^2+bx+c,\qquad a,b,c\in\mathbb{Z}_7.
\]

Imponiendo las condiciones:
\[
\begin{aligned}
f(1)=a+b+c=2,\\
f(3)=2a+3b+c=3,\\
f(6)=a+6b+c=5.
\end{aligned}
\]
Luego surge el siguiente sistema:
\[
\begin{cases}
a+b+c=2,\\
2a+3b+c=3,\\
a+6b+c=5.
\end{cases}
\]

Aplicamos Gauss--Jordan en $\mathbb{Z}_7$ con la matriz ampliada:
\[
(A\,|\,B)=
\begin{pmatrix}
1 & 1 & 1 &\Big|& 2\\
2 & 3 & 1 &\Big|& 3\\
1 & 6 & 1 &\Big|& 5
\end{pmatrix}
\overset{\substack{F_2+5F_1\\F_3+6F_1\\{}}}{\sim_f}
\begin{pmatrix}
1 & 1 & 1 &\Big|& 2\\
0 & 1 & 6 &\Big|& 6\\
0 & 5 & 0 &\Big|& 3
\end{pmatrix}
\overset{\substack{F_3+2F_2\\{}}}{\sim_f}
\begin{pmatrix}
1 & 1 & 1 &\Big|& 2\\
0 & 1 & 6 &\Big|& 6\\
0 & 0 & 5 &\Big|& 1
\end{pmatrix}
\]
\[
\overset{\substack{3F_3\\{}}}{\sim_f}
\begin{pmatrix}
1 & 1 & 1 &\Big|& 2\\
0 & 1 & 6 &\Big|& 6\\
0 & 0 & 1 &\Big|& 3
\end{pmatrix}
\overset{\substack{F_1+6F_3\\F_2+F_3\\{}}}{\sim_f}
\begin{pmatrix}
1 & 1 & 0 &\Big|& 6\\
0 & 1 & 0 &\Big|& 2\\
0 & 0 & 1 &\Big|& 3
\end{pmatrix}
\overset{\substack{F_1+6F_2\\{}}}{\sim_f}
\begin{pmatrix}
1 & 0 & 0 &\Big|& 4\\
0 & 1 & 0 &\Big|& 2\\
0 & 0 & 1 &\Big|& 3
\end{pmatrix}
= H_f(A\,|\,B).
\]

Por tanto,
\[
\begin{cases}
    a = 4 \\
    b = 2 \\ 
    c = 3
\end{cases}
\]

Entonces,
\[
\boxed{f(x)=4x^2+2x+3 \in \mathbb{Z}_7[x]}
\]

\end{ejercicio}

\begin{ejercicio}[title=Ejercicio 39]
Encuentre el único polinomio $f(x) \in \mathbb{Z}_7[x]$ de grado menor o igual que 3 que verifica $f(2)=1$, $f(3)=0$, $f(4)=2$, $f(5)=2$.

\tcblower
\textbf{Resolución:}

Por el Ejercicio 37, sabemos que existe un único polinomio que cumple las condiciones dadas.

Como $\operatorname{grado}(f(x))\le 3$, escribimos
\[
f(x)=ax^3+bx^2+cx+d,\qquad a,b,c,d\in\mathbb{Z}_7.
\]

Imponiendo las condiciones y trabajando en $\mathbb{Z}_7$:
\[
\begin{aligned}
f(2)&=8a+4b+2c+d=1 \;\Rightarrow\; a+4b+2c+d=1,\\
f(3)&=27a+9b+3c+d=0 \;\Rightarrow\; 6a+2b+3c+d=0,\\
f(4)&=64a+16b+4c+d=2 \;\Rightarrow\; a+2b+4c+d=2,\\
f(5)&=125a+25b+5c+d=2 \;\Rightarrow\; 6a+4b+5c+d=2.
\end{aligned}
\]

Luego surge el sistema:
\[
\begin{cases}
a+4b+2c+d=1,\\
6a+2b+3c+d=0,\\
a+2b+4c+d=2,\\
6a+4b+5c+d=2.
\end{cases}
\]

Aplicamos Gauss--Jordan en $\mathbb{Z}_7$ con la matriz ampliada:
\[
(A\,|\,B)=
\begin{pmatrix}
1 & 4 & 2 & 1 &\Big|& 1\\
6 & 2 & 3 & 1 &\Big|& 0\\
1 & 2 & 4 & 1 &\Big|& 2\\
6 & 4 & 5 & 1 &\Big|& 2
\end{pmatrix}
\overset{\substack{F_2+F_1\\F_3+6F_1\\F_4+F_1\\{}}}{\sim_f}
\begin{pmatrix}
1 & 4 & 2 & 1 &\Big|& 1\\
0 & 6 & 5 & 2 &\Big|& 1\\
0 & 5 & 2 & 0 &\Big|& 1\\
0 & 1 & 0 & 2 &\Big|& 3
\end{pmatrix}
\overset{\substack{F_3+5F_2\\F_4+F_2\\{}}}{\sim_f}
\begin{pmatrix}
1 & 4 & 2 & 1 &\Big|& 1\\
0 & 6 & 5 & 2 &\Big|& 1\\
0 & 0 & 6 & 3 &\Big|& 6\\
0 & 0 & 5 & 4 &\Big|& 4
\end{pmatrix}
\]
\vspace{0.4em}
\[
\overset{\substack{F_4+5F_3\\{}}}{\sim_f}
\begin{pmatrix}
1 & 4 & 2 & 1 &\Big|& 1\\
0 & 6 & 5 & 2 &\Big|& 1\\
0 & 0 & 6 & 3 &\Big|& 6\\
0 & 0 & 0 & 5 &\Big|& 6
\end{pmatrix}
\overset{\substack{6F_2\\6F_3\\3F_4\\{}}}{\sim_f}
\begin{pmatrix}
1 & 4 & 2 & 1 &\Big|& 1\\
0 & 1 & 2 & 5 &\Big|& 6\\
0 & 0 & 1 & 4 &\Big|& 1\\
0 & 0 & 0 & 1 &\Big|& 4
\end{pmatrix}
\overset{\substack{F_3+3F_4\\{}}}{\sim_f}
\begin{pmatrix}
1 & 4 & 2 & 1 &\Big|& 1\\
0 & 1 & 2 & 5 &\Big|& 6\\
0 & 0 & 1 & 0 &\Big|& 6\\
0 & 0 & 0 & 1 &\Big|& 4
\end{pmatrix}
\]
\vspace{0.4em}
\[
\overset{\substack{F_2+5F_3\\{}}}{\sim_f}
\begin{pmatrix}
1 & 4 & 2 & 1 &\Big|& 1\\
0 & 1 & 0 & 5 &\Big|& 1\\
0 & 0 & 1 & 0 &\Big|& 6\\
0 & 0 & 0 & 1 &\Big|& 4
\end{pmatrix}
\overset{\substack{F_2+2F_4\\{}}}{\sim_f}
\begin{pmatrix}
1 & 4 & 2 & 1 &\Big|& 1\\
0 & 1 & 0 & 0 &\Big|& 2\\
0 & 0 & 1 & 0 &\Big|& 6\\
0 & 0 & 0 & 1 &\Big|& 4
\end{pmatrix}
\overset{\substack{F_1+3F_2\\{}}}{\sim_f}
\begin{pmatrix}
1 & 0 & 2 & 1 &\Big|& 0\\
0 & 1 & 0 & 0 &\Big|& 2\\
0 & 0 & 1 & 0 &\Big|& 6\\
0 & 0 & 0 & 1 &\Big|& 4
\end{pmatrix}
\]
\end{ejercicio}
\begin{ejercicio}[title=Ejercicio 39]
\[
\overset{\substack{F_1+5F_3\\{}}}{\sim_f}
\begin{pmatrix}
1 & 0 & 0 & 1 &\Big|& 2\\
0 & 1 & 0 & 0 &\Big|& 2\\
0 & 0 & 1 & 0 &\Big|& 6\\
0 & 0 & 0 & 1 &\Big|& 4
\end{pmatrix}
\overset{\substack{F_1+6F_4\\{}}}{\sim_f}
\begin{pmatrix}
1 & 0 & 0 & 0 &\Big|& 5\\
0 & 1 & 0 & 0 &\Big|& 2\\
0 & 0 & 1 & 0 &\Big|& 6\\
0 & 0 & 0 & 1 &\Big|& 4
\end{pmatrix}
= H_f(A\,|\,B).
\]

Por tanto,
\[
\begin{cases}
a = 5,\\
b = 2,\\
c = 6,\\
d = 4.
\end{cases}
\]

Entonces,
\[
\boxed{f(x)=5x^3+2x^2+6x+4 \in \mathbb{Z}_7[x]}
\]

\end{ejercicio}

\begin{ejercicio}[title=Ejercicio 41]
Sea $S$ un sistema de ecuaciones lineales sobre un cuerpo $\mathbb{K}$ escrito en forma matricial como $A \cdot X = B$, donde $A \in \mathfrak{M}_{m\times n}(\mathbb{K})$, y sea $S'$ el sistema homogéneo $A \cdot X = (0)_{m\times 1}$.  
Supongamos que $(\alpha_1,\ldots,\alpha_n)$ es una solución de $S$.  
Pruebe que la aplicación $f$ definida del conjunto de soluciones de $S'$ en el conjunto de soluciones de $S$ mediante
\[
f(x_1,\ldots,x_n) = (x_1,\ldots,x_n) + (\alpha_1,\ldots,\alpha_n)
\]
es biyectiva.

\tcblower
\textbf{Resolución:}

Los sistemas vienen definidos por las expresiones matriciales:
\[
S': A \cdot X = (0)_{m\times 1}, \qquad 
S: A \cdot X = B.
\]

Sean $\operatorname{Sol}(S')$ y $\operatorname{Sol}(S)$ los conjuntos de soluciones de los sistemas $S'$ y $S$. Tenemos la función:
\[
f: \text{Sol}(S') \longrightarrow \text{Sol}(S), \quad 
f(x_1,\dots,x_n) = (x_1,\dots,x_n) + (\alpha_1,\dots,\alpha_n)
\]

\vspace{0.5em}

\textbf{Inyectividad.}

Supongamos que 
\[
f(x_1,\ldots,x_n) = f(y_1,\ldots,y_n).
\]
Por definición de $f$:
\[
(x_1,\ldots,x_n) + (\alpha_1,\ldots,\alpha_n) = (y_1,\ldots,y_n) + (\alpha_1,\ldots,\alpha_n).
\]
Restando $(\alpha_1,\ldots,\alpha_n)$ en ambos lados obtenemos:
\[
(x_1,\ldots,x_n) = (y_1,\ldots,y_n).
\]
Por lo tanto, $f$ es inyectiva.

\vspace{0.5em}

\textbf{Sobreyectividad.}

Queremos demostrar que para cualquier solución $(\beta_1,\ldots,\beta_n)$ de $S$ existe una solución $(x_1,\ldots,x_n)$ de $S'$ tal que
\[
f(x_1,\ldots,x_n) = (\beta_1,\ldots,\beta_n).
\]

Entonces sustituyendo la expresión de $f(x_1,\ldots,x_n)$ tenemos:
\[
(x_1,\ldots,x_n) + (\alpha_1,\ldots,\alpha_n) = (\beta_1,\ldots,\beta_n)
\]
\[
(x_1,\ldots,x_n) = (\beta_1,\ldots,\beta_n) - (\alpha_1,\ldots,\alpha_n).
\]

Como $(\beta_1,\ldots,\beta_n)$ es solución de $S$, es decir,
\(
(\beta_1,\ldots,\beta_n) \in \operatorname{Sol}(S),
\)
tenemos
\[
A \cdot (\beta_1,\ldots,\beta_n) = B.
\]
\end{ejercicio}
\begin{ejercicio}[title=Ejercicio 41]
Como $(\alpha_1,\ldots,\alpha_n)$ también es solución de $S$, es decir, 
\(
(\alpha_1,\ldots,\alpha_n) \in \operatorname{Sol}(S),
\)
se cumple
\[
A \cdot (\alpha_1,\ldots,\alpha_n) = B.
\]
Restando estas dos ecuaciones obtenemos
\[
A \cdot \big((\beta_1,\ldots,\beta_n) - (\alpha_1,\ldots,\alpha_n)\big) = B - B = (0)_{m\times 1},
\]
lo que implica que 
\[
(x_1,\ldots,x_n) = (\beta_1,\ldots,\beta_n) - (\alpha_1,\ldots,\alpha_n) \in \operatorname{Sol}(S'),
\]
es decir, $(x_1,\ldots,x_n)$ pertenece al dominio de $f$.

Entonces al aplicar la función $f$ a este elemento del dominio, obtenemos
\[
f(x_1,\ldots,x_n) = (x_1,\ldots,x_n) + (\alpha_1,\ldots,\alpha_n).
\]

Sustituyendo cuánto es $(x_1,\ldots,x_n)$, tenemos
\[
f(x_1,\ldots,x_n) = \big((\beta_1,\ldots,\beta_n) - (\alpha_1,\ldots,\alpha_n)\big) + (\alpha_1,\ldots,\alpha_n) = (\beta_1,\ldots,\beta_n).
\]

Esto significa que para cualquier solución $(\beta_1,\ldots,\beta_n) \in \operatorname{Sol}(S)$ hemos encontrado un elemento $(x_1,\ldots,x_n) \in \operatorname{Sol}(S')$ tal que $f(x_1,\ldots,x_n) = (\beta_1,\ldots,\beta_n)$.  
En otras palabras, cada elemento del codominio pertenece a la imagen de $f$, es decir, el codominio y la imagen son lo mismo ($\operatorname{Sol}(S)$). Esto demuestra que $f$ es sobreyectiva.

\medskip

Como $f$ es inyectiva y sobreyectiva,
\[
\boxed{\text{$f$ es biyectiva.}}
\]

\end{ejercicio}

\section{Matrices Regulares}

\begin{ejercicio}[title=Ejercicio 13]
Aplique operaciones elementales de filas para comprobar que la matriz 
\[
A =
\begin{pmatrix}
4 & 1 & 1\\
5 & 5 & 3\\
9 & 2 & 7
\end{pmatrix}
\]
con coeficientes en $\mathbb{Z}_{11}$ es regular, y obtenga su inversa.

\tcblower
\textbf{Resolución:}

\[
(A \,|\, I_3) =
\begin{pmatrix}
4 & 1 & 1 &\Big|& 1 & 0 & 0\\
5 & 5 & 3 &\Big|& 0 & 1 & 0\\
9 & 2 & 7 &\Big|& 0 & 0 & 1
\end{pmatrix}
\overset{\substack{3F_1 \\ {}}}{\sim_f}
\begin{pmatrix}
1 & 3 & 3 &\Big|& 3 & 0 & 0\\
5 & 5 & 3 &\Big|& 0 & 1 & 0\\
9 & 2 & 7 &\Big|& 0 & 0 & 1
\end{pmatrix}
\]
\vspace{0.5em}
\[
\overset{\substack{F_2 + 6F_1 \\ F_3 + 2F_1}}{\sim_f}
\begin{pmatrix}
1 & 3 & 3 &\Big|& 3 & 0 & 0\\
0 & 1 & -1 &\Big|& 7 & 1 & 0\\
0 & 8 & 2 &\Big|& 6 & 0 & 1
\end{pmatrix}
\overset{\substack{F_1 + F_3 \\ {}}}{\sim_f}
\begin{pmatrix}
1 & 0 & 5 &\Big|& 9 & 0 & 1\\
0 & 1 & -1 &\Big|& 7 & 1 & 0\\
0 & 8 & 2 &\Big|& 6 & 0 & 1
\end{pmatrix}
\]
\[
\overset{\substack{F_3 + 3F_2 \\ {}}}{\sim_f}
\begin{pmatrix}
1 & 0 & 5 &\Big|& 9 & 0 & 1\\
0 & 1 & -1 &\Big|& 7 & 1 & 0\\
0 & 0 & -1 &\Big|& 5 & 3 & 1
\end{pmatrix}.
\]

\noindent
\[
\begin{pmatrix}
1 & 0 & 5\\
0 & 1 & -1\\
0 & 0 & -1
\end{pmatrix}
\text{ está en forma escalonada por filas.}
\]

\[\operatorname{Rg}(A) = 3 = \text{orden de } \mathfrak{M}_3(\mathbb{Z}_{11}) \implies \boxed{\text{$A$ regular}}
\]

\[
\begin{pmatrix}
1 & 0 & 5 &\Big|& 9 & 0 & 1\\
0 & 1 & -1 &\Big|& 7 & 1 & 0\\
0 & 0 & -1 &\Big|& 5 & 3 & 1
\end{pmatrix}
\overset{\substack{F_2 - F_3 \\ F_1 + 5F_3}}{\sim_f}
\begin{pmatrix}
1 & 0 & 0 &\Big|& 1 & 4 & 6\\
0 & 1 & 0 &\Big|& 2 & -2 & -1\\
0 & 0 & -1 &\Big|& 5 & 3 & 1
\end{pmatrix}
\]
\vspace{0.5em}
\[
\overset{\substack{-F_3 \\ {}}}{\sim_f}
\begin{pmatrix}
1 & 0 & 0 &\Big|& 1 & 4 & 6\\
0 & 1 & 0 &\Big|& 2 & 9 & 10\\
0 & 0 & 1 &\Big|& 6 & 8 & 10
\end{pmatrix}
= (I_3 \, | \, A^{-1})
\quad \implies \quad
\boxed{
A^{-1} =
\begin{pmatrix}
1 & 4 & 6\\
2 & 9 & 10\\
6 & 8 & 10
\end{pmatrix}
}
\]
\end{ejercicio}

\begin{ejercicio}[title=Ejercicio 14]
Calcule mediante operaciones elementales de filas la inversa de cada una de las matrices siguientes:
\[
A =
\begin{pmatrix}
1 & 1 & 1 & 1\\
0 & 1 & 1 & 1\\
0 & 0 & 1 & 1\\
0 & 0 & 0 & 1
\end{pmatrix}, \qquad
B =
\begin{pmatrix}
1 & 2 & 3 & 4\\
0 & 1 & 2 & 3\\
0 & 0 & 1 & 2\\
0 & 0 & 0 & 1
\end{pmatrix}.
\]

\tcblower
\textbf{Resolución:}

\[
(A \,|\, I_4) =
\begin{pmatrix}
1 & 1 & 1 & 1 &\Big|& 1 & 0 & 0 & 0\\
0 & 1 & 1 & 1 &\Big|& 0 & 1 & 0 & 0\\
0 & 0 & 1 & 1 &\Big|& 0 & 0 & 1 & 0\\
0 & 0 & 0 & 1 &\Big|& 0 & 0 & 0 & 1
\end{pmatrix}
\overset{\substack{F_1 - F_2 \\ {}}}{\sim_f}
\begin{pmatrix}
1 & 0 & 0 & 0 &\Big|& 1 & -1 & 0 & 0\\
0 & 1 & 1 & 1 &\Big|& 0 & 1 & 0 & 0\\
0 & 0 & 1 & 1 &\Big|& 0 & 0 & 1 & 0\\
0 & 0 & 0 & 1 &\Big|& 0 & 0 & 0 & 1
\end{pmatrix}
\]

\[
\overset{\substack{F_2 - F_3 \\ {}}}{\sim_f}
\begin{pmatrix}
1 & 0 & 0 & 0 &\Big|& 1 & -1 & 0 & 0\\
0 & 1 & 0 & 0 &\Big|& 0 & 1 & -1 & 0\\
0 & 0 & 1 & 1 &\Big|& 0 & 0 & 1 & 0\\
0 & 0 & 0 & 1 &\Big|& 0 & 0 & 0 & 1
\end{pmatrix}
\overset{\substack{F_3 - F_4 \\ {}}}{\sim_f}
\begin{pmatrix}
1 & 0 & 0 & 0 &\Big|& 1 & -1 & 0 & 0\\
0 & 1 & 0 & 0 &\Big|& 0 & 1 & -1 & 0\\
0 & 0 & 1 & 0 &\Big|& 0 & 0 & 1 & -1\\
0 & 0 & 0 & 1 &\Big|& 0 & 0 & 0 & 1
\end{pmatrix}
\]

\[
= (I_4 \,|\, A^{-1}) \quad \implies \quad
\boxed{
A^{-1} =
\begin{pmatrix}
1 & -1 & 0 & 0\\
0 & 1 & -1 & 0\\
0 & 0 & 1 & -1\\
0 & 0 & 0 & 1
\end{pmatrix}
}
\]

\vspace{0.5em}

\[
(B \,|\, I_4) =
\begin{pmatrix}
1 & 2 & 3 & 4 &\Big|& 1 & 0 & 0 & 0\\
0 & 1 & 2 & 3 &\Big|& 0 & 1 & 0 & 0\\
0 & 0 & 1 & 2 &\Big|& 0 & 0 & 1 & 0\\
0 & 0 & 0 & 1 &\Big|& 0 & 0 & 0 & 1
\end{pmatrix}
\overset{\substack{F_3 - 2F_4 \\ {}}}{\sim_f}
\begin{pmatrix}
1 & 2 & 3 & 4 &\Big|& 1 & 0 & 0 & 0\\
0 & 1 & 2 & 3 &\Big|& 0 & 1 & 0 & 0\\
0 & 0 & 1 & 0 &\Big|& 0 & 0 & 1 & -2\\
0 & 0 & 0 & 1 &\Big|& 0 & 0 & 0 & 1
\end{pmatrix}
\]

\[
\overset{\substack{F_2 - 2F_3 \\ {}}}{\sim_f}
\begin{pmatrix}
1 & 2 & 3 & 4 &\Big|& 1 & 0 & 0 & 0\\
0 & 1 & 0 & 3 &\Big|& 0 & 1 & -2 & 4\\
0 & 0 & 1 & 0 &\Big|& 0 & 0 & 1 & -2\\
0 & 0 & 0 & 1 &\Big|& 0 & 0 & 0 & 1
\end{pmatrix}
\overset{\substack{F_2 - 3F_4 \\ {}}}{\sim_f}
\begin{pmatrix}
1 & 2 & 3 & 4 &\Big|& 1 & 0 & 0 & 0\\
0 & 1 & 0 & 0 &\Big|& 0 & 1 & -2 & 1\\
0 & 0 & 1 & 0 &\Big|& 0 & 0 & 1 & -2\\
0 & 0 & 0 & 1 &\Big|& 0 & 0 & 0 & 1
\end{pmatrix}
\]
\end{ejercicio}
\begin{ejercicio}[title=Ejercicio 14]
\[
\overset{\substack{F_1 - 2F_2 \\ {}}}{\sim_f}
\begin{pmatrix}
1 & 0 & 3 & 4 &\Big|& 1 & -2 & 4 & -2\\
0 & 1 & 0 & 0 &\Big|& 0 & 1 & -2 & 1\\
0 & 0 & 1 & 0 &\Big|& 0 & 0 & 1 & -2\\
0 & 0 & 0 & 1 &\Big|& 0 & 0 & 0 & 1
\end{pmatrix}
\overset{\substack{F_1 - 3F_3 \\ {}}}{\sim_f}
\begin{pmatrix}
1 & 0 & 0 & 4 &\Big|& 1 & -2 & 1 & 4\\
0 & 1 & 0 & 0 &\Big|& 0 & 1 & -2 & 1\\
0 & 0 & 1 & 0 &\Big|& 0 & 0 & 1 & -2\\
0 & 0 & 0 & 1 &\Big|& 0 & 0 & 0 & 1
\end{pmatrix}
\]

\[
\overset{\substack{F_1 - 4F_4 \\ {}}}{\sim_f}
\begin{pmatrix}
1 & 0 & 0 & 0 &\Big|& 1 & -2 & 1 & 0\\
0 & 1 & 0 & 0 &\Big|& 0 & 1 & -2 & 1\\
0 & 0 & 1 & 0 &\Big|& 0 & 0 & 1 & -2\\
0 & 0 & 0 & 1 &\Big|& 0 & 0 & 0 & 1
\end{pmatrix}
= (I_4 \,|\, B^{-1})
\]
\vspace{0.5em}
\[
\implies 
\boxed{
B^{-1} =
\begin{pmatrix}
1 & -2 & 1 & 0\\
0 & 1 & -2 & 1\\
0 & 0 & 1 & -2\\
0 & 0 & 0 & 1
\end{pmatrix}
}
\]
\end{ejercicio}

\begin{ejercicio}[title=Ejercicio 18]
En $\mathfrak{M}_2(\mathbb{R})$ obtenga todas las matrices $X$ tales que
\[
X \cdot B \cdot (A + A^6) - (X \cdot B - B^2) \cdot A = A^7 + B^2 \cdot A,
\]
donde
\[
A =
\begin{pmatrix}
1 & 3\\
1 & 2
\end{pmatrix}, \qquad
B =
\begin{pmatrix}
1 & 2\\
0 & 1
\end{pmatrix}.
\]

\tcblower
\textbf{Resolución:}

Simplificamos la expresión y despejamos $X$:

\[
XB(A+A^6) - (XB-B^2)A = A^7 + B^2 A
\implies 
 X(BA+BA^6) - (XBA-B^2 A) = A^7 + B^2 A
\]
\[
\implies 
X(BA+BA^6) - XBA+B^2 = A^7+B^2 A
\implies 
 X(BA+BA^6-BA) = A^7
\]
\[
\implies 
 X(BA^6) = A^7
\implies 
 X(BA^6)(BA^6)^{-1} = A^7 (BA^6)^{-1}
\implies 
 X \cdot I = A^7 (BA^6)^{-1}
\]

\[
\text{Comprobamos si $A$ y $B$ son regulares:}
\]

\vspace{0.5em}

\[
A =
\begin{pmatrix}
1 & 3\\
1 & 2
\end{pmatrix}
\overset{\substack{F_2 - F_1 \\ {}}}{\sim_f}
\begin{pmatrix}
1 & 3\\
0 & -1
\end{pmatrix}
\begin{aligned}{
\text{ está en forma escalonada por filas} \\
$\Rightarrow \operatorname{Rg}(A) = \text{nº filas no nulas} = 2$
\end{aligned}
\]

\vspace{0.5em}

Como \(\operatorname{Rg}(A)=2=\) orden de \(A\Rightarrow A\) es regular.

\vspace{0.5em}

\[
B =
\begin{pmatrix}
1 & 2\\
0 & 1
\end{pmatrix}
\begin{aligned}{
\text{ está en forma escalonada por filas} \\
$\Rightarrow \operatorname{Rg}(B) = \text{nº filas no nulas} = 2$
\end{aligned}
\]

\vspace{0.5em}

Como \(\operatorname{Rg}(B)=2=\) orden de \(B\Rightarrow B\) es regular.

\vspace{0.5em}

\[
\Rightarrow X = A^7 (BA^6)^{-1} = A^7 (A^6)^{-1} B^{-1} = A\cdot A^6 \cdot A^{-6} \cdot B^{-1} = A \cdot I \cdot B^{-1} = AB^{-1}
\]
\[
\Rightarrow X = AB^{-1}.
\]

Hallamos \(B^{-1}\):

\[
(B \,|\, I_2) =
\begin{pmatrix}
1 & 2 &\Big|& 1 & 0\\
0 & 1 &\Big|& 0 & 1
\end{pmatrix}
\overset{\substack{F_1 - 2F_2 \\ {}}}{\sim_f}
\begin{pmatrix}
1 & 0 &\Big|& 1 & -2\\
0 & 1 &\Big|& 0 & 1
\end{pmatrix}
= (I_2 \,|\, B^{-1})
\]
\[
\Rightarrow B^{-1} =
\begin{pmatrix}
1 & -2\\
0 & 1
\end{pmatrix}.
\]
Entonces:
\[
X = AB^{-1} = 
\begin{pmatrix}
1 & 3\\
1 & 2
\end{pmatrix}
\begin{pmatrix}
1 & -2\\
0 & 1
\end{pmatrix}
=
\boxed{\,\begin{pmatrix}1 & 1\\[4pt]1 & 0\end{pmatrix}\,}
\]
\end{ejercicio}

\begin{ejercicio}[title=Ejercicio 23]
Sea $A \in \mathfrak{M}_4(\mathbb{R})$ tal que $A^3 + 5A - I_4$ es la matriz nula.  
Justifique que $A$ es una matriz regular y obtenga una expresión para la inversa de $A$.

\tcblower
\textbf{Resolución:}

Para justificar que $A$ es regular, por definición necesitamos encontrar una matriz $B$ tal que:
\[
A \cdot B = I_4 = B \cdot A.
\]

Despejamos $I_4$ de la ecuación dada:
\[
A^3 + 5A - I_4 = 0_4 \quad \Rightarrow \quad A^3 + 5A = I_4.
\]

Factorizando tanto por la izquierda como por la derecha obtenemos:
\begin{itemize}
\item $A(A^2 + 5 I_4) = I_4$
\item $(A^2 + 5 I_4)A = I_4$
\end{itemize}

Hemos encontrado que $B = A^2 + 5 I_4$ satisface la ecuación de la definición de matriz regular.  
\[
\boxed{\text{A es regular}}.
\]

\vspace{1em}

Si pasamos $A$ al lado contrario en cualquiera de las ecuaciones obtenemos una expresión para la inversa:
\[
\boxed{A^{-1} = A^2 + 5 I_4}.
\]

\end{ejercicio}

\section{Determinantes}

\begin{ejercicio}[title=Ejercicio 16]
En cada apartado determine los valores de $a$ para que la matriz dada sea una matriz regular en $\mathfrak{M}_3(\mathbb{Z}_7)$:

\[
\begin{aligned}
\text{a)}\;&
\begin{pmatrix}
a & 1 & 1\\
1 & a & 1\\
1 & 1 & a
\end{pmatrix}
\qquad\qquad
\text{b)}\;&
\begin{pmatrix}
a & 2 & 1\\
1 & a & 4\\
1 & 1 & a
\end{pmatrix}
\qquad\qquad
\text{c)}\;&
\begin{pmatrix}
a & 2 & 5\\
2 & a & 4\\
1 & 1 & a
\end{pmatrix}
\end{aligned}
\]

\tcblower
\textbf{Resolución:}

$A$ regular $\Leftrightarrow |A| \neq 0$ en $\mathbb{Z}_7$.

\vspace{0.5em}

\textbf{a)}
Calculamos el determinante:
\[
\left|
\begin{matrix}
a & 1 & 1\\
1 & a & 1\\
1 & 1 & a
\end{matrix}
\right|
= a^3 + 4a + 2 \quad (\text{en } \mathbb{Z}_7)
\]

Probamos $a = 1$ usando Ruffini:
\[
\begin{array}{c|cccc}
  & 1 & 0 & 4 & 2 \\
1 &   & 1 & 1 & 5  \\
\hline
  & 1 & 1 & 5 & \boxed{0} \\
1 &   & 1 & 2 \\
\hline
  & 1 & 2 & \boxed{0}
\end{array}
\]

Por lo tanto, $a = 1$ es raíz doble. Factorizando:
\[
|A| = a^3 + 4a + 2 = (a - 1)^2 (a + 2) = (a - 1)^2 (a - 5)
\]

Para que $A$ sea regular, necesitamos $|A| \neq 0$, es decir:
\[
a \neq 1 \quad \text{y} \quad a \neq 5.
\]
\[
\boxed{a \in \mathbb{Z}_7 \setminus \{1,5\}}
\]

\textbf{b)}
Calculamos el determinante:
\[
\left|
\begin{matrix}
a & 2 & 1\\
1 & a & 4\\
1 & 1 & a
\end{matrix}
\right|
= a^3 + 2 \quad (\text{en } \mathbb{Z}_7)
\]
Probamos las raíces en $\mathbb{Z}_7$ para ver si anulan el determinante:

\[
a = 0 \Rightarrow 0 + 2 = 2 \neq 0
\]
\[
a = 1 \Rightarrow 1 + 2 = 3 \neq 0
\]
\[
a = 2 \Rightarrow 8 + 2 = 10 \equiv 3 \neq 0
\]
\[
a = 3 \Rightarrow 27 + 2 = 29 \equiv 1 \neq 0
\]
\end{ejercicio}
\begin{ejercicio}[title=Ejercicio 16]
\[
a = 4 \Rightarrow 64 + 2 = 66 \equiv 3 \neq 0
\]
\[
a = 5 \Rightarrow 125 + 2 = 127 \equiv 1 \neq 0
\]
\[
a = 6 \Rightarrow 216 + 2 = 218 \equiv 1 \neq 0
\]

Como $|A| \neq 0$ para ningún $a \in \mathbb{Z}_7$, $A$ es regular para todos los $a$ en $\mathbb{Z}_7$:
\[
\boxed{a \in \mathbb{Z}_7}
\]

\vspace{0.5em}

\textbf{c)}
Calculamos el determinante:
\[
\left|
\begin{matrix}
a & 2 & 5\\
2 & a & 4\\
1 & 1 & a
\end{matrix}
\right|
= a^3 + a + 4 \quad (\text{en } \mathbb{Z}_7)
\]

Vemos que $a=2$ es raíz. Aplicamos Ruffini:
\[
\begin{array}{c|cccc}
  & 1 & 0 & 1 & 4 \\
2 &   & 2 & 4 & 3 \\
\hline
  & 1 & 2 & 5 & \boxed{0} \\
2 &  & 2 & 1 \\
\hline
  & 1 & 4 & \boxed{6} \\
\end{array}
\]
De ahí:
\[
a^3 + a + 4 = (a - 2)(a^2 + 2a + 5) \quad (\text{en } \mathbb{Z}_7).
\]

Comprobamos si \(a^2 + 2a + 5\) tiene raíces en \(\mathbb{Z}_7\):

\[
a^2+2a+5=0 \quad\rightarrow\quad a=\frac{5\pm\sqrt{4+1}}{2} = \frac{5\pm\sqrt{5}}{2} \notin \mathbb{Z}_7
\]

Conclusión: la única raíz del polinomio es \(a=2\). Para que la matriz sea regular necesitamos \( |A| \neq 0\), luego \(a\) no puede ser \(2\). Por tanto:
\[
\boxed{a \in \mathbb{Z}_7 \setminus \{2\}}
\]

\end{ejercicio}

\begin{ejercicio}[title=Ejercicio 32]
Calcule el determinante de la matriz
\[
A =
\begin{pmatrix}
4 & 2 & 3 & 5\\
7 & -2 & 6 & -8\\
3 & 2 & 2 & -5\\
2 & -7 & 3 & 2
\end{pmatrix}
\]
con coeficientes en $\mathbb{R}$.

\tcblower
\textbf{Resolución:}

Para calcular el determinante de una matriz cuadrada de orden 4 no podemos aplicar la Regla de Sarrus. En su lugar, realizaremos operaciones por filas o columnas hasta transformarla en una matriz cuya determinante sabemos calcular (por ejemplo la identidad, cuyo determinante es 1). Antes de operar recordamos las propiedades que afectan al determinante:

\begin{itemize}
\item Si intercambiamos dos filas (o dos columnas), el determinante cambia de signo.
\item Sacar fuera del determinante un escalar que multiplica a una fila o columna equivale a multiplicar el determinante por ese escalar. En la práctica, multiplicar una fila por $\lambda$ multiplica el determinante por $\lambda$ (y dividir una fila por $\lambda$ divide el determinante por $\lambda$).
\item Si sumamos a una fila (o columna) una combinación lineal de otras filas (o columnas), el determinante no varía.
\end{itemize}

Teniendo esto en cuenta, procedemos con las operaciones elementales por filas:

\[
\left|
\begin{matrix}
4 & 2 & 3 & 5\\
7 & -2 & 6 & -8\\
3 & 2 & 2 & -5\\
2 & -7 & 3 & 2
\end{matrix}
\right|
\overset{\substack{F_1-F_3\\F_2-2F_3\\{}}}{=}
\left|
\begin{matrix}
1 & 0 & 1 & 10\\
1 & -6 & 2 & 2\\
3 & 2 & 2 & -5\\
2 & -7 & 3 & 2
\end{matrix}
\right|
\overset{\substack{F_2-F_1\\F_3-3F_1\\F_4-2F_1\\{}}}{=}
\left|
\begin{matrix}
1 & 0 & 1 & 10\\
0 & -6 & 1 & -8\\
0 & 2 & -1 & -35\\
0 & -7 & 1 & -18
\end{matrix}
\right|
\]

\[
\overset{\substack{F_2-F_4\\{}}}{=}
\left|
\begin{matrix}
1 & 0 & 1 & 10\\
0 & 1 & 0 & 10\\
0 & 2 & -1 & -35\\
0 & -7 & 1 & -18
\end{matrix}
\right|
\overset{\substack{F_3-2F_2\\F_4+7F_2\\{}}}{=}
\left|
\begin{matrix}
1 & 0 & 1 & 10\\
0 & 1 & 0 & 10\\
0 & 0 & -1 & -55\\
0 & 0 & 1 & 52
\end{matrix}
\right|
\overset{\substack{F_3+F_4\\{}}}{=}
\left|
\begin{matrix}
1 & 0 & 1 & 10\\
0 & 1 & 0 & 10\\
0 & 0 & 0 & -3\\
0 & 0 & 1 & 52
\end{matrix}
\right|
\]

Intercambiamos filas \(F_3\leftrightarrow F_4\) (esto cambia el signo del determinante):

\[
\overset{\substack{F_3\leftrightarrow F_4\\{}}}{=}
-\,\left|
\begin{matrix}
1 & 0 & 1 & 10\\
0 & 1 & 0 & 10\\
0 & 0 & 1 & 52\\
0 & 0 & 0 & -3
\end{matrix}
\right|
\]
\end{ejercicio}
\begin{ejercicio}[title=Ejercicio 32]
Multiplicar la fila 4 por $-\tfrac{1}{3}$ es lo mismo que a extraer el factor $-3$ fuera del determinante; por eso aparece el $-(-3)$ multiplicando.

\[
\overset{\substack{F_1-F_3\\-\tfrac{1}{3}F_4\\{}}}{=}
-\,(-3)\,
\left|
\begin{matrix}
1 & 0 & 0 & -42\\
0 & 1 & 0 & 10\\
0 & 0 & 1 & 52\\
0 & 0 & 0 & 1
\end{matrix}
\right|
\overset{\substack{F_1+42F_4\\F_2-10F_4\\F_3-52F_4\\{}}}{=}
-(-3)\,
\left|
\begin{matrix}
1 & 0 & 0 & 0\\
0 & 1 & 0 & 0\\
0 & 0 & 1 & 0\\
0 & 0 & 0 & 1
\end{matrix}
\right|
= 3\,\left|I_4\right| = 3\cdot 1 = \boxed{3}
\]
\end{ejercicio}

\begin{ejercicio}[title=Ejercicio 33]
Sin utilizar la Regla de Sarrus, justifique que cada uno de los determinantes siguientes vale cero:

\[
\begin{tabular}{clclcl}
$a)\begin{vmatrix}
a & d & 2a - 3d\\
b & e & 2b - 3e\\
c & f & 2c - 3f
\end{vmatrix}\quad$ &
$b)\begin{vmatrix}
y + z & z + x & x + y\\
x & y & z\\
1 & 1 & 1
\end{vmatrix}\quad$ &
$c)\begin{vmatrix}
a + 1 & a + 2 & a + 3\\
a + 4 & a + 5 & a + 6\\
a + 7 & a + 8 & a + 9
\end{vmatrix}$
\end{tabular}
\]

\tcblower
\textbf{Resolución:}

\textbf{a)}

\[
\begin{vmatrix}
a & d & 2a-3d\\
b & e & 2b-3e\\
c & f & 2c-3f
\end{vmatrix}
=
\begin{vmatrix}
a & d & 2a\\
b & e & 2b\\
c & f & 2c
\end{vmatrix}
+
\begin{vmatrix}
a & d & -3d\\
b & e & -3e\\
c & f & -3f
\end{vmatrix}
=2
\underbrace{
\begin{vmatrix}
a & d & a\\
b & e & b\\
c & f & c
\end{vmatrix}}_{0}
-3
\underbrace{
\begin{vmatrix}
a & d & d\\
b & e & e\\
c & f & f
\end{vmatrix}}_{0}
= \boxed{0}
\]
porque en el primer determinante las columnas \(1\) y \(3\) son iguales y en el segundo las columnas \(2\) y \(3\) son iguales.

\vspace{0.5em}

\textbf{b)}
Al igual que en el ejercicio 32, aplicamos operaciones elementales de fila o columna teniendo en cuenta las propiedades de los determinantes:

\[
\left|
\begin{matrix}
y+z & z+x & x+y \\
x & y & z \\
1 & 1 & 1
\end{matrix}
\right|
\overset{\substack{F_1+F_2\\{}}}{=}
\left|
\begin{matrix}
x+y+z & x+y+z & x+y+z \\
x & y & z \\
1 & 1 & 1
\end{matrix}
\right|
\]
\[
\overset{\substack{\frac{1}{x+y+z}F_1\\{}}}{=}
(x+y+z)
\underbrace{
\left|
\begin{matrix}
1 & 1 & 1 \\
x & y & z \\
1 & 1 & 1
\end{matrix}
\right|
}_{0}
= \boxed{0}
\]

porque las filas 1 y 3 son iguales.

\vspace{0.5em}

\textbf{c)}  
\[
\begin{vmatrix}
a+1 & a+2 & a+3\\
a+4 & a+5 & a+6\\
a+7 & a+8 & a+9
\end{vmatrix}
=
\begin{vmatrix}
a & a+2 & a+3\\
a & a+5 & a+6\\
a & a+8 & a+9
\end{vmatrix}
+
\begin{vmatrix}
1 & a+2 & a+3\\
4 & a+5 & a+6\\
7 & a+8 & a+9
\end{vmatrix}
\]

\[
=
\left(
\underbrace{
\begin{vmatrix}
a & a & a+3\\
a & a & a+6\\
a & a & a+9
\end{vmatrix}
}_{0}
+
\begin{vmatrix}
a & 2 & a+3\\
a & 5 & a+6\\
a & 8 & a+9
\end{vmatrix}
\right)
+
\left(
\begin{vmatrix}
1 & a & a+3\\
4 & a & a+6\\
7 & a & a+9
\end{vmatrix}
+
\begin{vmatrix}
1 & 2 & a+3\\
4 & 5 & a+6\\
7 & 8 & a+9
\end{vmatrix}
\right)
\]
\[
=
\left(
\underbrace{
\begin{vmatrix}
a & 2 & a\\
a & 5 & a\\
a & 8 & a
\end{vmatrix}
}_{0}
+
\begin{vmatrix}
a & 2 & 3\\
a & 5 & 6\\
a & 8 & 9
\end{vmatrix}
\right)
+
\left(
\underbrace{
\begin{vmatrix}
1 & a & a\\
4 & a & a\\
7 & a & a
\end{vmatrix}
}_{0}
+
\begin{vmatrix}
1 & a & 3\\
4 & a & 6\\
7 & a & 9
\end{vmatrix}
\right)
+
\left(
\begin{vmatrix}
1 & 2 & a\\
4 & 5 & a\\
7 & 8 & a
\end{vmatrix}
+
\begin{vmatrix}
1 & 2 & 3\\
4 & 5 & 6\\
7 & 8 & 9
\end{vmatrix}
\right)
\]

porque todos los determinantes señalados tienen columnas repetidas.
\end{ejercicio}
\begin{ejercicio}[title=Ejercicio 33]
Además, los determinantes restantes también son nulos, porque en todos ellos las columnas son combinaciones lineales de las otras.
\begin{itemize}
    \item Primera: $C_1 = aC_3 - aC_2$
    \item Segunda: $C_1 = C_3 - \frac{2}{a}C_2$
    \item Tercera: $C_1 = C_2 - \frac{1}{a}C_3$
    \item Cuarta: $C_1 = 2C_2-C_3$
\end{itemize}
\[
\underbrace{
\begin{vmatrix}
a & 2 & 3\\
a & 5 & 6\\
a & 8 & 9
\end{vmatrix}
}_{0}
+
\underbrace{
\begin{vmatrix}
1 & a & 3\\
4 & a & 6\\
7 & a & 9
\end{vmatrix}
}_{0}
+
\underbrace{
\begin{vmatrix}
1 & 2 & a\\
4 & 5 & a\\
7 & 8 & a
\end{vmatrix}
}_{0}
+
\underbrace{
\begin{vmatrix}
1 & 2 & 3\\
4 & 5 & 6\\
7 & 8 & 9
\end{vmatrix}
}_{0}
= \boxed{0}
\]
\end{ejercicio}

\subsection{Matriz Adjunta}

\begin{ejercicio}[title=Ejercicio 42 (⋆)]
Las siguientes afirmaciones se refieren a matrices cuadradas de orden $n$, con $n \ge 2$, y con coeficientes en un cuerpo $\mathbb{K}$. Determine las que son ciertas, justificando su respuesta, y ponga un contraejemplo para aquellas que sean falsas:

\begin{enumerate}
\item[a)] $\forall A \in \mathfrak{M}_n(\mathbb{K})$, la matriz $A + A^T$ es simétrica.
\item[b)] $\forall A \in \mathfrak{M}_n(\mathbb{K})$, la matriz $A \cdot A^T$ es simétrica.
\item[c)] $\forall A \in \mathfrak{M}_n(\mathbb{K})$, la matriz $A - A^T$ es antisimétrica.
\item[d)] Si $A$ es simétrica, entonces $A^2$ es simétrica.
\item[e)] Si $A^2$ es simétrica, entonces $A$ es simétrica.
\item[f)] Si $A$ y $B$ son ambas simétricas, entonces $A + B$ es simétrica.
\item[g)] Si $A$ y $B$ son ambas simétricas, entonces $A \cdot B$ es simétrica.
\item[h)] Si $A \cdot B = B \cdot A$, entonces $A^2 \cdot B^2 = B^2 \cdot A^2$.
\item[i)] Si $A^2 \cdot B^2 = B^2 \cdot A^2$, entonces $A \cdot B = B \cdot A$.
\item[j)] Si $A^2$ y $B^2$ son ambas la matriz nula, entonces $(A + B)^2$ es la matriz nula.
\item[k)] Si $A^2 = I_n$, entonces $A = I_n$ o $A = -I_n$.
\item[l)] Si $|A| = 1$, entonces $|2A| = 2$.
\item[m)] Si $|A| = 0$, entonces $|A + I_n| \ne 0$.
\item[n)] Si $A$ es antisimétrica y $n$ es impar, entonces $|A| = 0$.
\item[ñ)] $\forall A \in \mathfrak{M}_n(\mathbb{K}),\ \operatorname{rg}(A) = \operatorname{rg}(A^2)$.
\item[o)] Si $A$ es regular, entonces $\overline{A}$, la matriz adjunta de $A$, es regular.
\item[p)] Si $A$ no es regular, entonces $\overline{A}$ no es regular.
\item[q)] Si $|A| = 1$, entonces $|\overline{A}| = 1$.
\item[r)] Si $|A| = -1$, entonces $|A| + |\overline{A}| = 0$.
\end{enumerate}

\tcblower
\textbf{Resolución:}

Recordamos que:
\[
(A+B)^T = A^T + B^T,\qquad
(AB)^T = B^T A^T.
\]
Una matriz $A$ es \emph{simétrica} si $A^T=A$, y es \emph{antisimétrica} si $A=-A^T$.

\end{ejercicio}
\begin{ejercicio}[title=Ejercicio 42 (⋆)]

\textbf{a)} $\forall A \in \mathfrak{M}_n(\mathbb{K})$, la matriz $A + A^T$ es simétrica.

\vspace{0.5em}

Esto significa que $A + A^T = (A + A^T)^T$.  
Vemos que:
\[
(A + A^T)^T = A^T + (A^T)^T = A^T + A = A + A^T.
\]
\[
\boxed{\textbf{Verdadera. }\quad  \forall A \in \mathfrak{M}_n(\mathbb{K})\,\,\,\, A + A^T \text{ simétrica.}}
\]

\vspace{1em}

\textbf{b)} $\forall A \in \mathfrak{M}_n(\mathbb{K})$, la matriz $A \cdot A^T$ es simétrica.

\vspace{0.5em}

Esto significa que $A \cdot A^T = (A \cdot A^T)^T$.  
Vemos que:
\[
(A \cdot A^T)^T = (A^T)^T \cdot A^T = A \cdot A^T.
\]
\[
\boxed{\textbf{Verdadera. }\quad  \forall A \in \mathfrak{M}_n(\mathbb{K})\,\,\,\, A \cdot A^T \text{ simétrica.}}
\]

\vspace{1em}

\textbf{c)} $\forall A \in \mathfrak{M}_n(\mathbb{K})$, la matriz $A - A^T$ es antisimétrica.

\vspace{0.5em}

Esto significa que $A - A^T = -(A - A^T)^T$.  
Vemos que:
\[
-(A - A^T)^T = -(A^T + (-A^T)^T) = -(A^T - (A^T)^T) = -(A^T - A) = A - A^T.
\]
\[
\boxed{\textbf{Verdadera. }\quad  \forall A \in \mathfrak{M}_n(\mathbb{K})\,\,\,\, A - A^T \text{ antisimétrica.}}
\]

\vspace{1em}

\textbf{d)} Si $A$ es simétrica, entonces $A^2$ es simétrica.

\vspace{0.5em}

Esto significa que $A = A^T \;\Rightarrow\; A^2 = (A^2)^T$.
Vemos que elevando al cuadrado:
\[
A^2 = (A^T)^2 = A^T \cdot A^T = (A \cdot A)^T = (A^2)^T
\]
\[
\boxed{\textbf{Verdadera. }\quad A \text{ simétrica } \Rightarrow A^2 \text{ simétrica.}}
\]

\vspace{1em}

\textbf{e)} Si $A^2$ es simétrica, entonces $A$ es simétrica.

\vspace{0.5em}

Esto significa que $A^2 = (A^2)^T \;\Rightarrow\; A = A^T$.  
Vemos que:
\[
A \cdot A = (A \cdot A)^T
\]
\[
A \cdot A = A^T \cdot A^T
\]
Es imposible deducir que $A = A^T$ de esta afirmación, por lo tanto es falsa.

\vspace{0.5em}

\textbf{Contraejemplo:}  
\[
A = \begin{pmatrix}0 & 1 \\ -1 & 0\end{pmatrix} \neq A^T = \begin{pmatrix}0 & -1 \\ 1 & 0\end{pmatrix}, \quad
A^2 = \begin{pmatrix}-1 & 0 \\ 0 & -1\end{pmatrix} = (A^2)^T
\]

\[
\boxed{\textbf{Falsa. }\quad  A^2 \text{ simétrica } \not\Rightarrow A \text{ simétrica.}}
\]

\end{ejercicio}
\begin{ejercicio}[title=Ejercicio 42 (⋆)]

\textbf{f)} Si $A$ y $B$ son ambas simétricas, entonces $A + B$ es simétrica.

\vspace{0.5em}

Esto significa que $A = A^T$ y $B = B^T \;\Rightarrow\; A + B = (A + B)^T$.  
Vemos que:
\[
(A + B)^T = A^T + B^T = A + B
\]
\[
\boxed{\textbf{Verdadera. }\quad A \text{ y } B \text{ simétricas } \Rightarrow A + B \text{ simétrica.}}
\]

\vspace{1em}

\textbf{g)} Si $A$ y $B$ son ambas simétricas, entonces $A \cdot B$ es simétrica.

\vspace{0.5em}

Esto significa que $A = A^T$ y $B = B^T \;\Rightarrow\; A \cdot B = (A \cdot B)^T$.  
Vemos que:
\[
(A \cdot B)^T = B^T \cdot A^T = B \cdot A \neq A \cdot B
\]
por lo que es falsa.

\textbf{Contraejemplo:}
\[
A = \begin{pmatrix}1 & 0 \\ 0 & 2\end{pmatrix} = A^T, \quad
B = \begin{pmatrix}0 & 1 \\ 1 & 0\end{pmatrix} = B^T
\]

\[
A \cdot B = \begin{pmatrix}1 & 0 \\ 0 & 2\end{pmatrix} \begin{pmatrix}0 & 1 \\ 1 & 0\end{pmatrix} = \begin{pmatrix}0 & 1 \\ 2 & 0\end{pmatrix} \neq (A \cdot B)^T = \begin{pmatrix}0 & 2 \\ 1 & 0\end{pmatrix}
\]

\[
\boxed{\textbf{Falsa. }\quad A \text{ y } B \text{ simétricas } \not\Rightarrow A \cdot B \text{ simétrica.}}
\]

\vspace{1em}

\textbf{h)} Si $A \cdot B = B \cdot A$, entonces $A^2 \cdot B^2 = B^2 \cdot A^2$.

\vspace{0.5em}

Esto significa que $A \cdot B = B \cdot A \;\Rightarrow\; A^2 \cdot B^2 = B^2 \cdot A^2$.  
Vemos que:
\[
A^2 \cdot B^2 = (A A)(B B) = A(AB)B = A(BA)B = (AB)(AB)
\]
\[
= (BA)(BA) = B(AB)A = B(BA)A = (BB)(AA) = B^2 \cdot A^2
\]

\[
\boxed{\textbf{Verdadera. }\quad A \cdot B = B \cdot A \Rightarrow A^2 \cdot B^2 = B^2 \cdot A^2.}
\]

\vspace{1em}

\textbf{i)} Si $A^2 \cdot B^2 = B^2 \cdot A^2$, entonces $A \cdot B = B \cdot A$.

\vspace{0.5em}

Esto significa que $A^2 \cdot B^2 = B^2 \cdot A^2 \;\Rightarrow\; A \cdot B = B \cdot A$.  
Esto es falso.

\vspace{0.5em}

\textbf{Contraejemplo:}
\[
A = \begin{pmatrix}0 & 1 \\ -1 & 0\end{pmatrix}, \quad
B = \begin{pmatrix}1 & 0 \\ 0 & -1\end{pmatrix}
\quad \Rightarrow \quad
A^2 = \begin{pmatrix}-1 & 0 \\ 0 & -1\end{pmatrix}, \quad
B^2 = \begin{pmatrix}1 & 0 \\ 0 & 1\end{pmatrix}
\]

\[
A^2 \cdot B^2 = \begin{pmatrix}-1 & 0 \\ 0 & -1\end{pmatrix}
= B^2 \cdot A^2
\]

\[
A \cdot B = \begin{pmatrix}0 & -1 \\ -1 & 0\end{pmatrix}
\neq
\begin{pmatrix}0 & 1 \\ 1 & 0\end{pmatrix} = B \cdot A
\]

\[
\boxed{\textbf{Falsa. }\quad A^2 \cdot B^2 = B^2 \cdot A^2 \not\Rightarrow A \cdot B = B \cdot A.}
\]
\end{ejercicio}
\begin{ejercicio}[title=Ejercicio 42 (⋆)]
\textbf{j)} Si $A^2$ y $B^2$ son ambas la matriz nula, entonces $(A + B)^2$ es la matriz nula.

\vspace{0.5em}

Esto significa que $A^2 = 0_n$ y $B^2 = 0_n \;\Rightarrow\; (A + B)^2 = 0_n$.  
Vemos que:
\[
(A + B)^2 = A^2 + A \cdot B + B \cdot A + B^2 = 0_n + A \cdot B + B \cdot A + 0_n = A \cdot B + B \cdot A \neq 0_n
\]
Por lo tanto, es falsa.

\vspace{0.5em}

\textbf{Contraejemplo:}
\[
A = \begin{pmatrix}0 & 1 \\ 0 & 0\end{pmatrix}, \quad
B = \begin{pmatrix}0 & 0 \\ 1 & 0\end{pmatrix}
\]

\[
A^2 = \begin{pmatrix}0 & 0 \\ 0 & 0\end{pmatrix}, \quad
B^2 = \begin{pmatrix}0 & 0 \\ 0 & 0\end{pmatrix}
\]

\[
(A + B)^2 = \begin{pmatrix}0 & 1 \\ 1 & 0\end{pmatrix}^2
= \begin{pmatrix}0 & 1 \\ 1 & 0\end{pmatrix}
\begin{pmatrix}0 & 1 \\ 1 & 0\end{pmatrix}
= \begin{pmatrix}1 & 0 \\ 0 & 1\end{pmatrix}
\neq \begin{pmatrix}0 & 0 \\ 0 & 0\end{pmatrix}
\]

\[
\boxed{\textbf{Falsa. }\quad A^2 \text{ y } B^2 \text{ nulas} \not\Rightarrow (A + B)^2 \text{ nula.}}
\]

\vspace{1em}

\textbf{k)} Si $A^2 = I_n$, entonces $A = I_n$ o $A = -I_n$.

\vspace{0.5em}

Esto significa que $A^2 = I_n \;\Rightarrow\; A = I_n \text{ o } A = -I_n$.  
Esto es falso.

\vspace{0.5em}

\textbf{Contraejemplo:}
\[
A = \begin{pmatrix}1 & 0 \\ 0 & -1\end{pmatrix} \neq I_2
\]

\[
A^2 =
\begin{pmatrix}1 & 0 \\ 0 & -1\end{pmatrix}
\begin{pmatrix}1 & 0 \\ 0 & -1\end{pmatrix}
=
\begin{pmatrix}1 & 0 \\ 0 & 1\end{pmatrix}
= I_2
\]

\[
\boxed{\textbf{Falsa. }\quad A^2 = I_n \not\Rightarrow A = I_n \text{ o } A = -I_n.}
\]

\vspace{1em}

\textbf{l)} Si $|A| = 1$, entonces $|2A| = 2$.

\vspace{0.5em}

Esto significa que $|A| = 1 \;\Rightarrow\; |2A| = 2$.  
Vemos que:
\[
|2A| = 2^n |A| = 2^n \neq 2
\]
por lo tanto, es falsa.

\vspace{0.5em}

\textbf{Contraejemplo:}
\[
|A| = \left|\begin{pmatrix}1 & 0 \\ 0 & 1\end{pmatrix}\right| = 1
\]

\[
|2A| = \left|\begin{pmatrix}2 & 0 \\ 0 & 2\end{pmatrix}\right| = 4 \neq 2
\]

\[
\boxed{\textbf{Falsa. }\quad |A| = 1 \not\Rightarrow |2A| = 2.}
\]
\end{ejercicio}
\begin{ejercicio}[title=Ejercicio 42 (⋆)]
\textbf{m)} Si $|A| = 0$, entonces $|A + I_n| \ne 0$.

\vspace{0.5em}

Esto significa que $|A| = 0 \;\Rightarrow\; |A + I_n| \ne 0$.  
Esto es falso.

\vspace{0.5em}

\textbf{Contraejemplo:}
\[
A = \begin{pmatrix}-1 & 0 \\ 0 & 0\end{pmatrix}, \quad |A| = 0
\]

\[
A + I_2 = \begin{pmatrix}0 & 0 \\ 0 & 1\end{pmatrix}, \quad |A + I_2| = 0
\]

\[
\boxed{\textbf{Falsa. }\quad |A| = 0 \not\Rightarrow |A + I_n| \ne 0.}
\]

\vspace{1em}

\textbf{n)} Si $A$ es antisimétrica y $n$ es impar, entonces $|A| = 0$.

\vspace{0.5em}

Esto significa que $A = -A^T$ y $n$ impar $\;\Rightarrow\; |A| = 0$.  
Vemos que:
\[
|A| = |-A^T| = (-1)^n |A^T| = (-1)^n |A|
\]
Como $n$ es impar, se tiene $(-1)^n = -1$. Entonces:
\[
|A| = -|A| \;\Rightarrow\; 2|A| = 0 \;\Rightarrow\; |A| = 0
\]
\[
\boxed{\textbf{Verdadera. }\quad A \text{ antisimétrica y } n \text{ impar } \Rightarrow |A| = 0.}
\]

\vspace{1em}

\textbf{ñ)} $\forall A \in \mathfrak{M}_n(\mathbb{K}),\ \operatorname{rg}(A) = \operatorname{rg}(A^2)$.

\vspace{0.5em}

Esto es falso.

\vspace{0.5em}

\textbf{Contraejemplo:}
\[
A = \begin{pmatrix}
0 & 1\\
0 & 0
\end{pmatrix}, \qquad
A^2 = \begin{pmatrix}
0 & 0\\
0 & 0
\end{pmatrix}
\]

\[
\operatorname{Rg}(A) = 1 \neq \operatorname{Rg}(A^2) = 0
\]
\[
\boxed{\textbf{Falsa. }\quad \exists A \in \mathfrak{M}_n(\mathbb{K}) \text{ tal que } \operatorname{Rg}(A) \neq \operatorname{Rg}(A^2).}
\]

\vspace{1em}

\textbf{o)} Si $A$ es regular, entonces $\overline{A}$, la matriz adjunta de $A$, es regular.

\vspace{0.5em}

Esto significa que $|A| \neq 0 \;\Rightarrow\; |\overline{A}| \neq 0$.  
Sabemos que:
\[
A \cdot \overline{A}^T = |A| \cdot I_n
\]
Tomando determinantes a ambos lados:
\[
|A \cdot \overline{A}^T| = \bigl||A| \cdot I_n\bigr|
\Rightarrow
|A| \cdot |\overline{A}^T| = |A|^n \cdot |I_n|
\Rightarrow
|A| \cdot |\overline{A}| = |A|^n
\]
\end{ejercicio}
\begin{ejercicio}[title=Ejercicio 42 (⋆)]
Como $|A| \neq 0$, podemos dividir:
\[
|\overline{A}| = |A|^{n-1} \neq 0
\]
\[
\boxed{\textbf{Verdadera. }\quad A \text{ regular } \Rightarrow \overline{A} \text{ regular.}}
\]
\textbf{p)} Si $A$ no es regular, entonces $\overline{A}$ no es regular.

\vspace{0.5em}

Esto significa que $|A| = 0 \;\Rightarrow\; |\overline{A}| = 0$.  

Sabemos que:
\[
A \cdot \overline{A}^T = |A| \cdot I_n.
\]

\textbf{Caso 1: $A$ es la matriz nula.}  
Si $A = 0_n$, entonces claramente $A$ no es regular y su matriz adjunta es
\[
\overline{A} = 0_n,
\]
que tampoco es regular. Por lo tanto, la afirmación se cumple trivialmente en este caso.

\vspace{0.5em}

\textbf{Caso 2: $A$ no es la matriz nula.} 

Supongamos que $A$ no es regular y que, por el contrario, $\overline{A}$ es regular.  
Como $\overline{A}^T$ es invertible, multiplicamos por su inversa:
\[
A = \underbrace{|A|}_{0} \cdot (\overline{A}^T)^{-1} = 0_n,
\]
lo cual contradice la hipótesis de que $A$ no es la matriz nula.  

Por lo tanto, $\overline{A}$ no puede ser regular si $A$ no lo es.

\[
\boxed{\textbf{Verdadera. }\quad A \text{ no regular } \Rightarrow \overline{A} \text{ no regular.}}
\]

\textbf{q)} Si $|A| = 1$, entonces $|\overline{A}| = 1$.

\vspace{0.5em}

Esto significa que $|A| = 1 \;\Rightarrow\; |\overline{A}| = 1$.  
Sabemos que:
\[
A \cdot \overline{A}^T = \underbrace{|A|}_{1} \cdot I_n
\]
Tomando determinantes a ambos lados:
\[
|A \cdot \overline{A}^T| = |I_n|
\Rightarrow
\underbrace{|A|}_{1} \cdot |\overline{A}^T| = |I_n|
\Rightarrow
|\overline{A}| = 1
\]
\[
\boxed{\textbf{Verdadera. }\quad |A| = 1 \Rightarrow |\overline{A}| = 1.}
\]

\vspace{1em}

\textbf{r)} Si $|A| = -1$, entonces $|A| + |\overline{A}| = 0$.

\vspace{0.5em}

Esto significa que $|A| = -1 \;\Rightarrow\; |\overline{A}| = 1.$

Sabemos que:
\[
A \cdot \overline{A}^T = \underbrace{|A|}_{-1} \cdot I_n
\]
\end{ejercicio}
\begin{ejercicio}[title=Ejercicio 42 (⋆)]
Tomando determinantes a ambos lados:
\[
|A \cdot \overline{A}^T| = |-I_n|
\Rightarrow
\underbrace{|A|}_{-1} \cdot |\overline{A}^T| = (-1)^n |I_n|
\Rightarrow
-|\overline{A}| = (-1)^n
\Rightarrow
|\overline{A}| = (-1)^{n+1} \ne 1
\]
Por lo tanto es falso.

\vspace{0.5em}
\textbf{Contraejemplo:}

\[
|A| = \begin{vmatrix}-1 & 0 \\ 0 & 1\end{vmatrix} = -1, \quad
|\overline{A}| = \begin{vmatrix}1 & 0 \\ 0 & -1\end{vmatrix} = -1 \neq 1
\]
\[
-1 + (-1) = -2 \ne 0
\]
\[
\boxed{\textbf{Falsa. }\quad |A| = -1 \not\Rightarrow |A| + |\overline{A}| = 0.}
\]
\end{ejercicio}

\subsection{Propiedades consecuencia}

\subsubsection{Desarrollos de Laplace}

\begin{ejercicio}[title=Ejercicio 34]
Calcule los determinantes siguientes:
\[
\begin{vmatrix}
1 & 1\\
a & b
\end{vmatrix},
\quad
\begin{vmatrix}
1 & 1 & 1\\
a & b & c\\
a^2 & b^2 & c^2
\end{vmatrix},
\quad
\begin{vmatrix}
1 & 1 & 1 & 1\\
a & b & c & d\\
a^2 & b^2 & c^2 & d^2\\
a^3 & b^3 & c^3 & d^3
\end{vmatrix}.
\]
¿Cómo se pueden generalizar los resultados obtenidos?

\tcblower
\textbf{Resolución:}

\textbf{Primer determinante.}

Se calcula directamente:
\[
\boxed{\,\begin{vmatrix}
1 & 1\\
a & b
\end{vmatrix} = b-a.}
\]

\bigskip

\textbf{Segundo determinante.}

Partimos de:
\[
\begin{vmatrix}
1 & 1 & 1\\
a & b & c\\
a^2 & b^2 & c^2
\end{vmatrix}.
\]

Aplicamos operaciones elementales por filas:
\[
\overset{F_3-aF_2}{=}
\begin{vmatrix}
1 & 1 & 1\\
a & b & c\\
0 & b^2-ab & c^2-ac
\end{vmatrix}
=
\begin{vmatrix}
1 & 1 & 1\\
a & b & c\\
0 & (b-a)b & (c-a)c
\end{vmatrix}.
\]

A continuación,
\[
\overset{F_2-aF_1}{=}
\begin{vmatrix}
1 & 1 & 1\\
0 & b-a & c-a\\
0 & (b-a)b & (c-a)c
\end{vmatrix}.
\]

Desarrollando por la primera columna, el determinante es igual a:
\[
\begin{vmatrix}
b-a & c-a\\
(b-a)b & (c-a)c
\end{vmatrix}.
\]

Sacamos los escalares que multiplican a las columnas:
\[
\begin{vmatrix}
b-a & c-a\\
(b-a)b & (c-a)c
\end{vmatrix}
=
(b-a)(c-a)
\begin{vmatrix}
1 & 1\\
b & c
\end{vmatrix}.
\]

Finalmente,
\[
\begin{vmatrix}
1 & 1\\
b & c
\end{vmatrix}
= c-b,
\]
\end{ejercicio}
\begin{ejercicio}[title=Ejercicio 34]
y por tanto
\[
\boxed{\,\begin{vmatrix}
1 & 1 & 1\\
a & b & c\\
a^2 & b^2 & c^2
\end{vmatrix}
= (b-a)(c-a)(c-b).}
\]

\textbf{Tercer determinante.}

Consideramos:
\[
\begin{vmatrix}
1 & 1 & 1 & 1\\
a & b & c & d\\
a^2 & b^2 & c^2 & d^2\\
a^3 & b^3 & c^3 & d^3
\end{vmatrix}.
\]

Aplicamos operaciones elementales por filas:
\[
\overset{F_4-aF_3}{=}
\begin{vmatrix}
1 & 1 & 1 & 1\\
a & b & c & d\\
a^2 & b^2 & c^2 & d^2\\
0 & b^3-ab^2 & c^3-ac^2 & d^3-ad^2
\end{vmatrix}
=
\begin{vmatrix}
1 & 1 & 1 & 1\\
a & b & c & d\\
a^2 & b^2 & c^2 & d^2\\
0 & (b-a)b^2 & (c-a)c^2 & (d-a)d^2
\end{vmatrix}.
\]

Después,
\[
\overset{F_3-aF_2}{=}
\begin{vmatrix}
1 & 1 & 1 & 1\\
a & b & c & d\\
0 & b^2-ab & c^2-ac & d^2-ad\\
0 & (b-a)b^2 & (c-a)c^2 & (d-a)d^2
\end{vmatrix}
=
\begin{vmatrix}
1 & 1 & 1 & 1\\
a & b & c & d\\
0 & (b-a)b & (c-a)c & (d-a)d\\
0 & (b-a)b^2 & (c-a)c^2 & (d-a)d^2
\end{vmatrix}.
\]

A continuación,
\[
\overset{F_2-aF_1}{=}
\begin{vmatrix}
1 & 1 & 1 & 1\\
0 & b-a & c-a & d-a\\
0 & (b-a)b & (c-a)c & (d-a)d\\
0 & (b-a)b^2 & (c-a)c^2 & (d-a)d^2
\end{vmatrix}.
\]

Desarrollando por la primera columna, el determinante es igual a:
\[
\begin{vmatrix}
b-a & c-a & d-a\\
(b-a)b & (c-a)c & (d-a)d\\
(b-a)b^2 & (c-a)c^2 & (d-a)d^2
\end{vmatrix}.
\]

Sacamos los factores comunes de cada columna:
\[
= (b-a)(c-a)(d-a)
\begin{vmatrix}
1 & 1 & 1\\
b & c & d\\
b^2 & c^2 & d^2
\end{vmatrix}.
\]

Este determinante es del mismo tipo que el segundo, sustituyendo
\(a\to b\), \(b\to c\), \(c\to d\). Por tanto,
\[
\begin{vmatrix}
1 & 1 & 1\\
b & c & d\\
b^2 & c^2 & d^2
\end{vmatrix}
= (c-b)(d-b)(d-c).
\]
\end{ejercicio}
\begin{ejercicio}[title=Ejercicio 34]
Finalmente,
\[
\boxed{\,\begin{vmatrix}
1 & 1 & 1 & 1\\
a & b & c & d\\
a^2 & b^2 & c^2 & d^2\\
a^3 & b^3 & c^3 & d^3
\end{vmatrix}
= (b-a)(c-a)(d-a)(c-b)(d-b)(d-c).}
\]

\textbf{Generalización.}

Estos determinantes son determinantes de \emph{Vandermonde}. En general, para
\[
\begin{vmatrix}
1 & \cdots & 1\\[4pt]
x_1 & \cdots & x_n\\[4pt]
x_1^2 & \cdots & x_n^2\\[4pt]
\vdots & \ddots & \vdots\\[4pt]
x_1^{n-1} & \cdots & x_n^{n-1}
\end{vmatrix}
\]
se tiene:
\[
\boxed{
\prod_{1\le i<j\le n}(x_j-x_i).
}
\]

\end{ejercicio}

\begin{ejercicio}[title=Ejercicio 35]
Resuelva la ecuación
\[
\begin{vmatrix}
x & a & b & c\\
a & x & b & c\\
a & b & x & c\\
a & b & c & x
\end{vmatrix}
= 0,
\]
cuya incógnita es $x$.

\tcblower
\textbf{Resolución:}

Aplicamos operaciones elementales por columnas:
\[
\begin{vmatrix}
x & a & b & c\\
a & x & b & c\\
a & b & x & c\\
a & b & c & x
\end{vmatrix}
\overset{\substack{C_1+C_2\\{}}}{=}
\begin{vmatrix}
x+a & a & b & c\\
x+a & x & b & c\\
a+b & b & x & c\\
a+b & b & c & x
\end{vmatrix}
\overset{\substack{C_1+C_3\\{}}}{=}
\begin{vmatrix}
x+a+b & a & b & c\\
x+a+b & x & b & c\\
x+a+b & b & x & c\\
x+b+c & b & c & x
\end{vmatrix}
\]
\[
\overset{\substack{C_1+C_4\\{}}}{=}
\begin{vmatrix}
x+a+b+c & a & b & c\\
x+a+b+c & x & b & c\\
x+a+b+c & b & x & c\\
x+a+b+c & b & c & x
\end{vmatrix}
\]
Sacamos $x+a+b+c$ que multiplica a la primera columna:
\[
= (x+a+b+c)
\begin{vmatrix}
1 & a & b & c\\
1 & x & b & c\\
1 & b & x & c\\
1 & b & c & x
\end{vmatrix}
\]
Aplicamos operaciones elementales de filas:
\[
\overset{\substack{F_2-F_1\\F_3-F_1\\F_4-F_1\\{}}}{=}
(x+a+b+c)
\begin{vmatrix}
1 & a & b & c\\
0 & x-a & 0 & 0\\
0 & b-a & x-b & 0\\
0 & b-a & c-b & x-c
\end{vmatrix}
\]
Como todos los elementos de la primera columna son 0 excepto el 1 al principio, es igual a: 
\[
=
(x+a+b+c)
\begin{vmatrix}
x-a & 0 & 0\\
b-a & x-b & 0\\
b-a & c-b & x-c
\end{vmatrix}
= (x+a+b+c)(x-a)(x-b)(x-c)
\]

Por tanto, las soluciones son:
\[
\boxed{x=a \;\;\;\text{o}\;\;\; x=b \;\;\;\text{o}\;\;\; x=c \;\;\;\text{o}\;\;\; x=-a-b-c.}
\]

\end{ejercicio}

\begin{ejercicio}[title=Ejercicio 36]
Compruebe que
\[
\begin{vmatrix}
a & b & c & d\\
b & a & d & c\\
c & d & a & b\\
d & c & b & a
\end{vmatrix}
=
\big((a+b)^2 - (c+d)^2\big)\,\big((a-b)^2 - (c-d)^2\big).
\]

\tcblower
\textbf{Resolución:}

\vspace{0.5em}

Aplicamos operaciones elementales de columnas:

\[
\begin{vmatrix}
a & b & c & d\\
b & a & d & c\\
c & d & a & b\\
d & c & b & a
\end{vmatrix}
\overset{\substack{C_1+C_2\\{}}}{=}
\begin{vmatrix}
a+b & b & c & d\\
a+b & a & d & c\\
c+d & d & a & b\\
c+d & c & b & a
\end{vmatrix}
\overset{\substack{C_1+C_3\\{}}}{=}
\begin{vmatrix}
a+b+c & b & c & d\\
a+b+c & a & d & c\\
a+c+d & d & a & b\\
b+c+d & c & b & a
\end{vmatrix}
\]

\[
\overset{\substack{C_1+C_4\\{}}}{=}
\begin{vmatrix}
a+b+c+d & b & c & d\\
a+b+c+d & a & d & c\\
a+b+c+d & d & a & b\\
a+b+c+d & c & b & a
\end{vmatrix}
\]

Sacamos el factor $a+b+c+d$ que multiplica a la primera columna:

\[
= (a+b+c+d)
\begin{vmatrix}
1 & b & c & d\\
1 & a & d & c\\
1 & d & a & b\\
1 & c & b & a
\end{vmatrix}
\]

Aplicamos operaciones elementales de filas:

\[
\overset{\substack{F_2-F_1\\F_3-F_1\\F_4-F_1\\{}}}{=}
\begin{vmatrix}
1 & b & c & d\\
0 & a-b & d-c & c-d\\
0 & d-b & a-c & b-d\\
0 & c-b & b-c & a-d
\end{vmatrix}
\]

Como la primera columna está compuesta por ceros excepto la primera, es igual a:
\[
= (a+b+c+d)
\begin{vmatrix}
a-b & d-c & c-d\\
d-b & a-c & b-d\\
c-b & b-c & a-d
\end{vmatrix}.
\]

Seguimos aplicando operaciones elementales de columnas:
\[
\overset{\substack{C_2+C_3\\{}}}{=}
(a+b+c+d)
\begin{vmatrix}
a-b & 0 & c-d\\
d-b & a+b-c-d & b-d\\
c-b & a+b-c-d & a-d
\end{vmatrix}
\]

Sacamos el factor $a+b-c-d$ que multiplica a la segunda columna:
\[
= (a+b+c+d)(a+b-c-d)
\begin{vmatrix}
a-b & 0 & c-d\\
d-b & 1 & b-d\\
c-b & 1 & a-d
\end{vmatrix}
\]
\end{ejercicio}
\begin{ejercicio}[title=Ejercicio 36]
Agrupamos cada término dentro de los paréntesis y seguimos aplicando operaciones elementales de filas:
\[
\overset{\substack{F_3-F_2\\{}}}{=}
\big((a+b)+(c+d)\big)\big((a+b)-(c+d)\big)
\begin{vmatrix}
a-b & 0 & c-d\\
d-b & 1 & b-d\\
c-d & 0 & a-b
\end{vmatrix}
\]

Agrupamos los términos como una diferencia de cuadrados. También vemos que la segunda columna son todos ceros excepto la segunda, así que es igual a:
\[
= \big((a+b)^2-(c+d)^2\big)
\begin{vmatrix}
a-b & c-d\\
c-d & a-b
\end{vmatrix}
\]

\[
= \boxed{\big((a+b)^2-(c+d)^2\big)\big((a-b)^2-(c-d)^2\big)}.
\]

\end{ejercicio}

\begin{ejercicio}[title=Ejercicio 40]
Calcule los siguientes determinantes en $\mathbb{R}$:
\[
\Delta_1 =
\begin{vmatrix}
a & 1 & 1 & 1 & 1\\
1 & a & 1 & 1 & 1\\
1 & 1 & a & 1 & 1\\
1 & 1 & 1 & a & 1\\
1 & 1 & 1 & 1 & a
\end{vmatrix}, \quad
\Delta_2 =
\begin{vmatrix}
a & a & a & a & a\\
b & a & a & a & a\\
b & b & a & a & a\\
b & b & b & a & a\\
b & b & b & b & a
\end{vmatrix}, \quad
\Delta_3 =
\begin{vmatrix}
a & 1 & 0 & 0 & 0\\
1 & a & 1 & 0 & 0\\
0 & 1 & a & 1 & 0\\
0 & 0 & 1 & a & 1\\
0 & 0 & 0 & 1 & a
\end{vmatrix}.
\]

\tcblower
\textbf{Resolución:}

\vspace{0.5em}

$\Delta_1$:

\vspace{0.5em}

Aplicamos operaciones elementales de columnas:
\[
\Delta_1 =
\begin{vmatrix}
a & 1 & 1 & 1 & 1\\
1 & a & 1 & 1 & 1\\
1 & 1 & a & 1 & 1\\
1 & 1 & 1 & a & 1\\
1 & 1 & 1 & 1 & a
\end{vmatrix}
\overset{\substack{C_1+C_2\\{}}}{=}
\begin{vmatrix}
a+1 & 1 & 1 & 1 & 1\\
a+1 & a & 1 & 1 & 1\\
2 & 1 & a & 1 & 1\\
2 & 1 & 1 & a & 1\\
2 & 1 & 1 & 1 & a
\end{vmatrix}
\overset{\substack{C_1+C_3\\{}}}{=}
\begin{vmatrix}
a+2 & 1 & 1 & 1 & 1\\
a+2 & a & 1 & 1 & 1\\
a+2 & 1 & a & 1 & 1\\
3 & 1 & 1 & a & 1\\
3 & 1 & 1 & 1 & a
\end{vmatrix}
\]
\[
\overset{\substack{C_1+C_4\\{}}}{=}
\begin{vmatrix}
a+3 & 1 & 1 & 1 & 1\\
a+3 & a & 1 & 1 & 1\\
a+3 & 1 & a & 1 & 1\\
a+3 & 1 & 1 & a & 1\\
4 & 1 & 1 & 1 & a
\end{vmatrix}
\overset{\substack{C_1+C_5\\{}}}{=}
\begin{vmatrix}
a+4 & 1 & 1 & 1 & 1\\
a+4 & a & 1 & 1 & 1\\
a+4 & 1 & a & 1 & 1\\
a+4 & 1 & 1 & a & 1\\
a+4 & 1 & 1 & 1 & a
\end{vmatrix}
\]

Sacamos el factor que multiplica a la primera columna:
\[
= (a+4)
\begin{vmatrix}
1 & 1 & 1 & 1 & 1\\
1 & a & 1 & 1 & 1\\
1 & 1 & a & 1 & 1\\
1 & 1 & 1 & a & 1\\
1 & 1 & 1 & 1 & a
\end{vmatrix}
\]

Aplicamos operaciones elementales de filas:
\[
\overset{\substack{F_2-F_1\\F_3-F_1\\F_4-F_1\\F_5-F_1\\{}}}{=}
(a+4)
\begin{vmatrix}
1 & 1 & 1 & 1 & 1\\
0 & a-1 & 0 & 0 & 0\\
0 & 0 & a-1 & 0 & 0\\
0 & 0 & 0 & a-1 & 0\\
0 & 0 & 0 & 0 & a-1
\end{vmatrix}
\]

Como la primera columna son todos ceros excepto el $1$, entonces es igual a:
\[
= (a+4)
\begin{vmatrix}
a-1 & 0 & 0 & 0\\
0 & a-1 & 0 & 0\\
0 & 0 & a-1 & 0\\
0 & 0 & 0 & a-1
\end{vmatrix}
= \boxed{(a+4)(a-1)^4}.
\]
\end{ejercicio}
\begin{ejercicio}[title=Ejercicio 40]
$\Delta_2$:

\vspace{0.5em}

Aplicamos operaciones elementales de filas:
\[
\Delta_2 =
\begin{vmatrix}
a & a & a & a & a\\
b & a & a & a & a\\
b & b & a & a & a\\
b & b & b & a & a\\
b & b & b & b & a
\end{vmatrix}
\overset{\substack{F_5-F_4\\{}}}{=}
\begin{vmatrix}
a & a & a & a & a\\
b & a & a & a & a\\
b & b & a & a & a\\
b & b & b & a & a\\
0 & 0 & 0 & b-a & 0
\end{vmatrix}
\overset{\substack{F_4-F_3\\{}}}{=}
\begin{vmatrix}
a & a & a & a & a\\
b & a & a & a & a\\
b & b & a & a & a\\
0 & 0 & b-a & 0 & 0\\
0 & 0 & 0 & b-a & 0
\end{vmatrix}
\]

\[
\overset{\substack{F_3-F_2\\{}}}{=}
\begin{vmatrix}
a & a & a & a & a\\
b & a & a & a & a\\
0 & b-a & 0 & 0 & 0\\
0 & 0 & b-a & 0 & 0\\
0 & 0 & 0 & b-a & 0
\end{vmatrix}
\overset{\substack{F_2-F_1\\{}}}{=}
\begin{vmatrix}
a & a & a & a & a\\
b-a & 0 & 0 & 0 & 0\\
0 & b-a & 0 & 0 & 0\\
0 & 0 & b-a & 0 & 0\\
0 & 0 & 0 & b-a & 0
\end{vmatrix}
\]

Sacamos el factor $a$ que multiplica a la primera fila:
\[
= a
\begin{vmatrix}
1 & 1 & 1 & 1 & 1\\
b-a & 0 & 0 & 0 & 0\\
0 & b-a & 0 & 0 & 0\\
0 & 0 & b-a & 0 & 0\\
0 & 0 & 0 & b-a & 0
\end{vmatrix}
\]

Como la última columna son todos ceros excepto el $1$, entonces es igual a:
\[
= a
\begin{vmatrix}
b-a & 0 & 0 & 0\\
0 & b-a & 0 & 0\\
0 & 0 & b-a & 0\\
0 & 0 & 0 & b-a
\end{vmatrix}
= \boxed{a(b-a)^4}.
\]

$\Delta_3$:

\vspace{0.5em}

Para hallar su valor, planteamos una relación de recurrencia con el objetivo de hallar una fórmula para $\Delta(n)$, el determinante de una matriz de la misma forma y tamaño $n \times n$. En nuestro caso, buscamos $\Delta(5)$.

\vspace{0.5em}

Calculamos los primeros casos base:
\[
\Delta(1) = |a| = a, \quad \Delta(2) = \begin{vmatrix} a & 1 \\ 1 & a \end{vmatrix} = a^2 - 1.
\]

Para $n>2$ consideramos el determinante
\[
\Delta(n)=
\begin{vmatrix}
a & 1 & 0 & 0 & \cdots & 0\\
1 & a & 1 & 0 & \cdots & 0\\
0 & 1 & a & 1 & \ddots & \vdots\\
0 & 0 & 1 & a & \ddots & 0\\
\vdots & \vdots & \ddots & \ddots & \ddots & 1\\
0 & 0 & \cdots & 0 & 1 & a
\end{vmatrix}.
\]
\end{ejercicio}
\begin{ejercicio}[title=Ejercicio 40]
Nos fijamos en la \textbf{primera columna}. En ella solo hay dos elementos distintos de cero:
\[
a \text{ (en la posición $(1,1)$)} \quad \text{y} \quad 1 \text{ (en la posición $(2,1)$)}.
\]

El determinante se calcula como, para una columna cualquiera, la suma de los productos:
\[
(\text{elemento})\cdot(\text{determinante de la submatriz resultante de eliminar su fila y columna}),
\]
multiplicados por su signo $(-1)^{i+j}$.  
Como en esta columna solo hay dos elementos no nulos, \textbf{el determinante es la suma de dos términos}.

\bigskip

\textbf{Primer término.}  
Tomamos el elemento $a$ de la primera fila y primera columna. Su signo es
\[
(-1)^{1+1}=+1.
\]
Al eliminar la primera fila y la primera columna queda el determinante
\[
\begin{vmatrix}
\color{gray}{a} & \color{lightgray}{1} & \color{lightgray}{0} & \color{lightgray}{0} & \color{lightgray}\cdots & \color{lightgray}{0}\\
\color{lightgray}{1} & a & 1 & 0 & \cdots & 0\\
\color{lightgray}{0} & 1 & a & 1 & \ddots & \vdots\\
\color{lightgray}{0} & 0 & 1 & a & \ddots & 0\\
\color{lightgray}\vdots & \vdots & \ddots & \ddots & \ddots & 1\\
\color{lightgray}{0} & 0 & \cdots & 0 & 1 & a
\end{vmatrix}
\quad\implies\quad
\begin{vmatrix}
a & 1 & 0 & \cdots & 0\\
1 & a & 1 & \ddots & \vdots\\
0 & 1 & a & \ddots & 0\\
\vdots & \ddots & \ddots & \ddots & 1\\
0 & \cdots & 0 & 1 & a
\end{vmatrix},
\]
que es de tamaño $(n-1)\times(n-1)$ y tiene exactamente la misma forma que el determinante original.  
Por definición, es $\Delta(n-1)$.  
Por tanto, el primer término es
\[
a\,\Delta(n-1).
\]

\medskip

\textbf{Segundo término.}  
Tomamos ahora el elemento $1$ de la segunda fila y primera columna. Su signo es
\[
(-1)^{2+1}=-1.
\]
Al eliminar la segunda fila y la primera columna obtenemos
\[
\begin{vmatrix}
\color{lightgray}{a} & 1 & 0 & 0 & \cdots & 0\\
\color{gray}{1} & \color{lightgray}{a} & \color{lightgray}{1} & \color{lightgray}{0} & \color{lightgray}\cdots & \color{lightgray}{0}\\
\color{lightgray}{0} & 1 & a & 1 & \ddots & \vdots\\
\color{lightgray}{0} & 0 & 1 & a & \ddots & 0\\
\color{lightgray}\vdots & \vdots & \ddots & \ddots & \ddots & 1\\
\color{lightgray}{0} & 0 & \cdots & 0 & 1 & a
\end{vmatrix}
\quad\implies\quad
\begin{vmatrix}
1 & 0 & 0 & \cdots & 0\\
1 & a & 1 & \cdots & 0\\
0 & 1 & a & \ddots & \vdots\\
\vdots & \vdots & \ddots & \ddots & 1\\
0 & 0 & \cdots & 1 & a
\end{vmatrix}.
\]

En esta matriz, la primera fila es $(1,0,0,\dots,0)$, por lo que su determinante coincide con el determinante de la matriz que queda al eliminar esa primera fila y esa primera columna.
\end{ejercicio}
\begin{ejercicio}[title=Ejercicio 40]
\[
\begin{vmatrix}
\color{lightgray}1 & \color{lightgray}0 & \color{lightgray}0 & \color{lightgray}\cdots & \color{lightgray}0\\
\color{lightgray}1 & a & 1 & \cdots & 0\\
\color{lightgray}0 & 1 & a & \ddots & \vdots\\
\color{lightgray}\vdots & \vdots & \ddots & \ddots & 1\\
\color{lightgray}0 & 0 & \cdots & 1 & a
\end{vmatrix}
\quad \implies \quad
\begin{vmatrix}
a & 1 & \cdots & 0\\
1 & a & \ddots & \vdots\\
\vdots & \ddots & \ddots & 1\\
0 & \cdots & 1 & a
\end{vmatrix}.
\]
que es de tamaño $(n-2)\times(n-2)$.  
Por definición, este determinante es $\Delta(n-2)$, y aparece restando debido al signo negativo.

\bigskip

\textbf{Conclusión.}  
Sumando ambos términos obtenemos la relación de recurrencia:
\[
\Delta(n)=a\,\Delta(n-1)-\Delta(n-2).
\]

Aplicamos esta fórmula sucesivamente:

\begin{itemize}
    \item Para $n=3$:
    \[ \Delta(3) = a \Delta(2) - \Delta(1) = a(a^2-1) - a = a^3 - a - a = a^3 - 2a. \]
    
    \item Para $n=4$:
    \[ \Delta(4) = a \Delta(3) - \Delta(2) = a(a^3-2a) - (a^2-1) = a^4 - 2a^2 - a^2 + 1 = a^4 - 3a^2 + 1. \]
    
    \item Para $n=5$:
    \[ \Delta(5) = a \Delta(4) - \Delta(3) = a(a^4 - 3a^2 + 1) - (a^3 - 2a) = a^5 - 3a^3 + a - a^3 + 2a. \]
\end{itemize}

Finalmente, agrupando términos obtenemos:
\[
\Delta_3 = \Delta(5) = \boxed{a^5 - 4a^3 + 3a}.
\]

\end{ejercicio}

\subsubsection{Caracterización de matrices regulares y fórmula para la inversa}

\begin{ejercicio}[title=Ejercicio 17]
Sea $a \in \mathbb{C}$ tal que $a^3 = 1$ y $a \ne 1$.  
Compruebe que la matriz siguiente es regular y calcule su inversa:
\[
A =
\begin{pmatrix}
1 & 1 & 1\\
1 & a & a^2\\
1 & a^2 & a
\end{pmatrix}.
\]

\tcblower
\textbf{Resolución:}

A es regular $\Longleftrightarrow |A| \neq 0$.  
Calculamos:
\[
|A| =
\begin{vmatrix}
1 & 1 & 1\\
1 & a & a^2\\
1 & a^2 & a
\end{vmatrix}
=
a^2 + a^2 + a^2 - a - a^4 - a
= -a^4 + 3a^2 - 2a
\]
\[
= -a \cdot \underbrace{a^3}_{1}
+ 3a^2 - 2a
= -a + 3a^2 - 2a
= 3a^2 - 3a
= 3a(a-1).
\]

Para que $A$ sea regular necesitamos:
\[
3a(a-1) \neq 0 \quad\Longrightarrow\quad a\neq 0 \ \text{o}\ a\neq 1.
\]

\begin{itemize}
    \item \(a\neq 0\): dado que el enunciado establece que \(a^3=1\), si \(a=0\) tendríamos \(0^3=0\neq 1\). Por lo tanto, \(a=0\) no forma parte de los posibles valores de la matriz y no nos afecta.
    
    \item \(a\neq 1\): el propio enunciado especifica que \(a\neq 1\), así que esta condición tampoco representa un problema.
\end{itemize}

Por lo tanto:
\[
\boxed{A \text{ es regular}.}
\]

Para calcular $A^{-1}$, usamos la fórmula de la adjunta. Calculamos $\overline{A}$ y su traspuesta:

\[
\overline{A} =
\begin{pmatrix}
    a^2-a^4 & -(a-a^2) & a^2-a \\
    -(a-a^2) & a-1 & -(a^2-1) \\
    a^2-a & -(a^2-1) & a-1
\end{pmatrix}
= \begin{pmatrix}
    a^2-a & a^2-a & a^2-a \\
    a^2-a & a-1 & -(a^2-1) \\
    a^2-a & -(a^2-1) & a-1
\end{pmatrix}
\]
\[
= (a-1) \begin{pmatrix}
    a & a & a \\
    a & 1 & -(a+1) \\
    a & -(a+1) & 1
\end{pmatrix}
\qquad
\overline{A}^T =
(a-1) \begin{pmatrix}
    a & a & a \\
    a & 1 & -(a+1) \\
    a & -(a+1) & 1
    \end{pmatrix}
\]

\vspace{0.5em}

Entonces la inversa es:

\[
A^{-1} = \frac{\overline{A}^T}{|A|} = \frac{\cancel{(a-1)}}{3a\cancel{(a-1)}} \begin{pmatrix}
    a & a & a \\
    a & 1 & -(a+1) \\
    a & -(a+1) & 1
\end{pmatrix}
= \boxed{\frac{1}{3}\begin{pmatrix}
    1 & 1 & 1 \\
    1 & \frac{1}{a} & -\frac{a+1}{a} \\
    1 & -\frac{a+1}{a} & \frac{1}{a}
\end{pmatrix}}
\]
\end{ejercicio}

\begin{ejercicio}[title=Ejercicio 19]
Calcule todas las matrices $X \in \mathfrak{M}_2(\mathbb{Z}_7)$ tales que $A \cdot X \cdot C = D$, donde
\[
A =
\begin{pmatrix}
2 & 3\\
1 & 2
\end{pmatrix}, \quad
C =
\begin{pmatrix}
3 & 1\\
2 & 5
\end{pmatrix}, \quad
D =
\begin{pmatrix}
2 & 3\\
6 & 2
\end{pmatrix}.
\]
Repita el mismo ejercicio considerando ahora
\[
C =
\begin{pmatrix}
3 & 6\\
4 & 1
\end{pmatrix}.
\]

\tcblower
\textbf{Resolución:}

\[
AXC = D \;\Rightarrow\; A^{-1} A X C C^{-1} = A^{-1} D C^{-1} \;\Rightarrow\; X = A^{-1} D C^{-1}
\]

\vspace{0.5em}

Hallamos $A^{-1}$ con operaciones elementales de fila (1) o con fórmula de la adjunta (2):
\[
(1) \quad (A \,|\, I_2) =
\begin{pmatrix}
2 & 3 &\Big|& 1 & 0\\
1 & 2 &\Big|& 0 & 1
\end{pmatrix}
\overset{\substack{F_1 - F_2 \\ {}}}{\sim_f}
\begin{pmatrix}
1 & 1 &\Big|& 1 & 6\\
1 & 2 &\Big|& 0 & 1
\end{pmatrix}
\overset{\substack{F_2 - F_1 \\ {}}}{\sim_f}
\begin{pmatrix}
1 & 1 &\Big|& 1 & 6\\
0 & 1 &\Big|& 6 & 2
\end{pmatrix}
\]

\vspace{0.5em}

\noindent
\[
\begin{pmatrix}
1 & 1\\
0 & 1
\end{pmatrix}
\text{ está en forma escalonada por filas }
\]
\[
\Rightarrow \operatorname{Rg}(A) =\text{nº filas no nulas }=2 = \text{orden de }A \Rightarrow A \text{ regular.}
\]

\vspace{0.5em}

\[
\begin{pmatrix}
1 & 1 &\Big|& 1 & 6\\
0 & 1 &\Big|& 6 & 2
\end{pmatrix}
\overset{\substack{F_1 - F_2 \\ {}}}{\sim_f}
\begin{pmatrix}
1 & 0 &\Big|& 2 & 4\\
0 & 1 &\Big|& 6 & 2
\end{pmatrix}
= (I_2 \,|\, A^{-1})
\]

\[
\Rightarrow A^{-1} =
\begin{pmatrix}
2 & 4\\
6 & 2
\end{pmatrix}.
\]

\[
(2) \quad |A| =
\left|
\begin{matrix}
2 & 3\\
1 & 2
\end{matrix}
\right| = 1,
\]

\[
\overline{A} =
\begin{pmatrix}
2 & 6\\
4 & 2
\end{pmatrix}
\quad\Rightarrow\quad
(\overline{A})^T =
\begin{pmatrix}
2 & 4\\
6 & 2
\end{pmatrix}
\quad\Rightarrow\quad
A^{-1} = \frac{(\overline{A})^T}{|A|} =
\begin{pmatrix}
2 & 4\\
6 & 2
\end{pmatrix}.
\]

\vspace{0.5em}

De la misma forma, hallamos $C^{-1}$:
\[
(1)\quad (C \,|\, I_2) =
\begin{pmatrix}
3 & 1 &\Big|& 1 & 0\\
2 & 5 &\Big|& 0 & 1
\end{pmatrix}
\overset{\substack{F_1 - F_2 \\ {}}}{\sim_f}
\begin{pmatrix}
1 & 3 &\Big|& 1 & 6\\
2 & 5 &\Big|& 0 & 1
\end{pmatrix}
\overset{\substack{F_2 - 2F_1 \\ {}}}{\sim_f}
\begin{pmatrix}
1 & 3 &\Big|& 1 & 6\\
0 & 6 &\Big|& 5 & 3
\end{pmatrix}
\]

\vspace{0.5em}

\noindent
\[
\begin{pmatrix}
1 & 3\\
0 & 6
\end{pmatrix}
\text{ está en forma escalonada por filas }
\]
\[\Rightarrow \operatorname{Rg}(C) =\text{nº filas no nulas }=2 = \text{orden de }C \Rightarrow C \text{ regular.}
\]
\end{ejercicio}
\begin{ejercicio}[title=Ejercicio 19]
\[
\begin{pmatrix}
1 & 3 &\Big|& 1 & 6\\
0 & 6 &\Big|& 5 & 3
\end{pmatrix}
\overset{\substack{-F_2 \\ {}}}{\sim_f}
\begin{pmatrix}
1 & 3 &\Big|& 1 & 6\\
0 & 1 &\Big|& 2 & 4
\end{pmatrix}
\overset{\substack{F_1 - 3F_2 \\ {}}}{\sim_f}
\begin{pmatrix}
1 & 0 &\Big|& 2 & 1\\
0 & 1 &\Big|& 2 & 4
\end{pmatrix}
= (I_2 \,|\, C^{-1})
\]

\[
\Rightarrow C^{-1} =
\begin{pmatrix}
2 & 1\\
2 & 4
\end{pmatrix}.
\]

\[
(2)\quad |C| =
\left|
\begin{matrix}
3 & 1\\
2 & 5
\end{matrix}
\right|
= 6 \equiv -1,
\]

\[
\overline{C} =
\begin{pmatrix}
5 & 5\\
6 & 3
\end{pmatrix}
\quad\Rightarrow\quad
(\overline{C})^T =
\begin{pmatrix}
5 & 6\\
5 & 3
\end{pmatrix}
\quad\Rightarrow\quad
C^{-1} = \frac{(\overline{C})^T}{|C|} =
\begin{pmatrix}
2 & 1\\
2 & 4
\end{pmatrix}.
\]
Entonces:

\[
X = A^{-1} D C^{-1} =
\begin{pmatrix}
2 & 4\\
6 & 2
\end{pmatrix}
\begin{pmatrix}
2 & 3\\
6 & 2
\end{pmatrix}
\begin{pmatrix}
2 & 1\\
2 & 4
\end{pmatrix}
=
\begin{pmatrix}
0 & 0\\
3 & 1
\end{pmatrix}
\begin{pmatrix}
2 & 1\\
2 & 4
\end{pmatrix}
=
\boxed{\begin{pmatrix}0 & 0\\[4pt]1 & 0\end{pmatrix}\,}
\]

\bigskip

\textbf{Considerando} \(C=\begin{pmatrix}3 & 6\\4 & 1\end{pmatrix}\):

\bigskip

\[
(1) \quad (C \,|\, I_2) =
\begin{pmatrix}
3 & 6 &\Big|& 1 & 0\\
4 & 1 &\Big|& 0 & 1
\end{pmatrix}
\overset{\substack{F_1 + 3F_2 \\ {}}}{\sim_f}
\begin{pmatrix}
1 & 2 &\Big|& 1 & 3\\
4 & 1 &\Big|& 0 & 1
\end{pmatrix}
\overset{\substack{F_2 - 4F_1 \\ {}}}{\sim_f}
\begin{pmatrix}
1 & 2 &\Big|& 1 & 3\\
0 & 0 &\Big|& 3 & 3
\end{pmatrix}
\]

\vspace{0.5em}

\noindent
\[
\begin{pmatrix}
1 & 2\\
0 & 0
\end{pmatrix}
\text{ está en forma escalonada por filas }
\]
\[\Rightarrow \operatorname{Rg}(C)=\text{nº filas no nulas }=1 \neq \text{orden de }C=2 \Rightarrow C \text{ no regular.}
\]

\[
(2)\quad |C| =
\left|
\begin{matrix}
3 & 6\\
4 & 1
\end{matrix}
\right|
= 0
\quad\Rightarrow\quad C \text{ no regular.}
\]

\[
\text{La única forma de averiguar }X\text{ es con un sistema de ecuaciones sin involucrar a $C^{-1}$:}
\]
\[
X \, C = A^{-1} D:
\]

\[
\begin{pmatrix} a & b \\ c & d \end{pmatrix}
\begin{pmatrix} 3 & 6 \\ 4 & 1 \end{pmatrix}
=
\begin{pmatrix} 2 & 4 \\ 6 & 2 \end{pmatrix}
\begin{pmatrix} 2 & 3 \\ 6 & 2 \end{pmatrix}
\Rightarrow
\begin{pmatrix}
3a+4b & 6a+b\\
3c+4d & 6c+d
\end{pmatrix}
=
\begin{pmatrix}
0 & 0\\
3 & 1
\end{pmatrix}
\]

\[
\Rightarrow
\begin{cases}
3a+4b = 0,\\[2pt]
6a+b = 0,\\[2pt]
3c+4d = 3,\\[2pt]
6c+d = 1.
\end{cases}
\]
\end{ejercicio}
\begin{ejercicio}[title=Ejercicio 19]
\[
\begin{aligned}
&3a-3b = 0\\
&-a+b = 0\\
&3c-3d = 3\\
&-c+d = 1
\end{aligned}
\]

\[
\Rightarrow
\begin{aligned}
& a-b = 0 \Rightarrow a = b,\\
& -a+a = 0 \Rightarrow 0 = 0,\\
& c-d = 1 \Rightarrow c = d+1 \Rightarrow (-d+1)+d = 1 \Rightarrow -1 = 1 \Rightarrow 6 = 1.
\end{aligned}
\]

\[
\Rightarrow \boxed{\text{$X$ no tiene solución}}
\]
\end{ejercicio}

\begin{ejercicio}[title=Ejercicio 20]
Encuentre todas las matrices $X \in \mathfrak{M}_2(\mathbb{Z}_{17})$ que verifican
\[
\begin{pmatrix}
2 & 8\\
1 & 14
\end{pmatrix} \cdot X + X \cdot
\begin{pmatrix}
0 & 1\\
2 & 15
\end{pmatrix}
=
\begin{pmatrix}
1 & 2\\
3 & 16
\end{pmatrix}
+ 2X \cdot
\begin{pmatrix}
8 & 9\\
1 & 7
\end{pmatrix}.
\]

\tcblower
\textbf{Resolución:}

Desarrollando paso a paso, despejamos $X$:

\[
\begin{pmatrix}
2 & 8\\
1 & 14
\end{pmatrix} X
+
X
\begin{pmatrix}
0 & 1\\
2 & 15
\end{pmatrix}
=
\begin{pmatrix}
1 & 2\\
3 & 16
\end{pmatrix}
+
2 X
\begin{pmatrix}
8 & 9\\
1 & 7
\end{pmatrix}
\]

\[
\begin{pmatrix}
2 & 8\\
1 & 14
\end{pmatrix} X
+
X
\begin{pmatrix}
0 & 1\\
2 & 15
\end{pmatrix}
=
\begin{pmatrix}
1 & 2\\
3 & 16
\end{pmatrix}
+
X
\begin{pmatrix}
16 & 1\\
2 & 14
\end{pmatrix}
\]

\[
\begin{pmatrix}
2 & 8\\
1 & 14
\end{pmatrix} X
+
X
\begin{pmatrix}
0 & 1\\
2 & 15
\end{pmatrix}
-
X
\begin{pmatrix}
16 & 1\\
2 & 14
\end{pmatrix}
=
\begin{pmatrix}
1 & 2\\
3 & 16
\end{pmatrix}
\]

\[
\begin{pmatrix}
2 & 8\\
1 & 14
\end{pmatrix} X
+
X
\begin{pmatrix}
0 & 1\\
2 & 15
\end{pmatrix}
+
X
\begin{pmatrix}
1 & 16\\
15 & 3
\end{pmatrix}
=
\begin{pmatrix}
1 & 2\\
3 & 16
\end{pmatrix}
\]

\[
\begin{pmatrix}
2 & 8\\
1 & 14
\end{pmatrix} X
+
X
\begin{pmatrix}
1 & 0\\
0 & 1
\end{pmatrix}
=
\begin{pmatrix}
1 & 2\\
3 & 16
\end{pmatrix}
\Rightarrow
\begin{pmatrix}
2 & 8\\
1 & 14
\end{pmatrix} X
+
X \cdot I_2
=
\begin{pmatrix}
1 & 2\\
3 & 16
\end{pmatrix}
\]
\[
\Rightarrow
\begin{pmatrix}
2 & 8\\
1 & 14
\end{pmatrix} X
+
I_2 \cdot X
=
\begin{pmatrix}
1 & 2\\
3 & 16
\end{pmatrix}
\Rightarrow
\begin{pmatrix}
2 & 8\\
1 & 14
\end{pmatrix} X
+
\begin{pmatrix}
1 & 0\\
0 & 1
\end{pmatrix}
X
=
\begin{pmatrix}
1 & 2\\
3 & 16
\end{pmatrix}
\]
\[
\Rightarrow
\begin{pmatrix}
3 & 8\\
1 & 15
\end{pmatrix} X
=
\begin{pmatrix}
1 & 2\\
3 & 16
\end{pmatrix}
\;\Rightarrow\;
X = 
\begin{pmatrix}
3 & 8\\
1 & 15
\end{pmatrix}^{-1}
\begin{pmatrix}
1 & 2\\
3 & 16
\end{pmatrix}
\]

\vspace{0.5em}

Calculamos el inverso:

\[
\left|
\begin{matrix}
3 & 8\\
1 & 15
\end{matrix}
\right| = 3 \ne 0 \quad \implies
\begin{pmatrix}
    3 & 8\\
1 & 15
\end{pmatrix}
\text{ regular}
\]

\[
{\overline{
\begin{pmatrix}
3 & 8\\
1 & 15
\end{pmatrix}
}}^T
=
\begin{pmatrix}
15 & 16\\
9 & 3
\end{pmatrix}^T
=
\begin{pmatrix}
15 & 9\\
16 & 3
\end{pmatrix}
\]

\[
X = 3^{-1} 
\begin{pmatrix}
15 & 9\\
16 & 3
\end{pmatrix}
\begin{pmatrix}
1 & 2\\
3 & 16
\end{pmatrix}
= 6 \cdot
\begin{pmatrix}
-2 & 9\\
-1 & 3
\end{pmatrix}
\begin{pmatrix}
1 & 2\\
3 & -1
\end{pmatrix}
=
\begin{pmatrix}
5 & 3\\
11 & 1
\end{pmatrix}
\begin{pmatrix}
1 & 2\\
3 & -1
\end{pmatrix}
=
\boxed{
\begin{pmatrix}
14 & 7\\
14 & 4
\end{pmatrix}}
\]
\end{ejercicio}

\begin{ejercicio}[title=Ejercicio 21]
En $\mathfrak{M}_2(\mathbb{R})$ encuentre las matrices $X$ que verifican
\[
(A \cdot X^{-1})^{-1} + B = X,
\]
donde
\[
A =
\begin{pmatrix}
3 & 1\\
2 & -1
\end{pmatrix}, \qquad
B =
\begin{pmatrix}
-2 & 1\\
1 & 1
\end{pmatrix}.
\]

\tcblower
\textbf{Resolución:}

Despejamos $X$:

\[
(A X^{-1})^{-1} + B = X
\]
\[
X A^{-1} + B = X
\]
\[
X A^{-1} - X = -B
\]
\[
X (A^{-1} - I_2) = -B
\]
\[
X = -B (A^{-1} - I_2)^{-1}
\]

Calculamos $A^{-1}$:

\[
|A| = \left|\begin{matrix}
    3 & 1 \\ 2 & -1
\end{matrix}\right| = -5 \ne 0 
\implies
A
\text{ regular}
\]

\[
\overline{A} = 
\begin{pmatrix}
-1 & -2\\
-1 & 3
\end{pmatrix}, \quad
\overline{A}^{\,T} = 
\begin{pmatrix}
-1 & -1\\
-2 & 3
\end{pmatrix} \quad \Rightarrow \quad
A^{-1} = \frac{1}{5} 
\begin{pmatrix}
1 & 1\\
2 & -3
\end{pmatrix}
\]

Calculamos $A^{-1}-I_2$ y su inversa:

\[
|A^{-1} - I_2| = \frac{32}{25} - \frac{2}{25} = \frac{6}{5} \ne 0 \implies
A^{-1} - I_2
\text{ regular}
\]

\[
A^{-1} - I_2 = 
\begin{pmatrix}
-4/5 & 1/5\\
2/5 & -8/5
\end{pmatrix}, \quad
\overline{A^{-1} - I_2} = 
\begin{pmatrix}
-8/5 & -2/5\\
-1/5 & -4/5
\end{pmatrix}, 
\]

\[
\overline{A^{-1} - I_2}^{\,T} = 
\begin{pmatrix}
-8/5 & -1/5\\
-2/5 & -4/5
\end{pmatrix} 
\quad\implies\quad
(A^{-1} - I_2)^{-1} = 
\begin{pmatrix}
-4/3 & -1/6\\
-1/3 & -2/3
\end{pmatrix}
\]

\[
X = - 
\begin{pmatrix}-2 & 1\\ 1 & 1\end{pmatrix}
\begin{pmatrix}-4/3 & -1/6\\ -1/3 & -2/3\end{pmatrix}
=
\begin{pmatrix}-2 & 1\\ 1 & 1\end{pmatrix}
\begin{pmatrix}4/3 & 1/6\\ 1/3 & 2/3\end{pmatrix}
=
\boxed{
\begin{pmatrix}
-7/3 & 1/3\\
5/3 & 5/6
\end{pmatrix}
}
\]

\end{ejercicio}

\begin{ejercicio}[title=Ejercicio 22]
En cada apartado obtenga todas las matrices $A \in \mathfrak{M}_2(\mathbb{Z}_{11})$ que verifican las condiciones indicadas:

\begin{enumerate}
\item[a)] 
$A^3 =
\begin{pmatrix}
1 & 8\\
10 & 9
\end{pmatrix}, \quad
A^4 =
\begin{pmatrix}
3 & 3\\
1 & 6
\end{pmatrix}.$

\item[b)] 
$A^4 =
\begin{pmatrix}
8 & 1\\
2 & 0
\end{pmatrix}, \quad
A^9 =
\begin{pmatrix}
5 & 4\\
8 & 6
\end{pmatrix}.$
\end{enumerate}

\tcblower
\textbf{Resolución:}

\vspace{0.5em}

\textbf{a)}

Sabemos que $A^3 \cdot A = A^4$, entonces
\[
\begin{pmatrix}1 & 8\\ 10 & 9\end{pmatrix} A = \begin{pmatrix}3 & 3\\ 1 & 6\end{pmatrix} 
\quad \Rightarrow \quad 
A = \begin{pmatrix}1 & 8\\ 10 & 9\end{pmatrix}^{-1} \begin{pmatrix}3 & 3\\ 1 & 6\end{pmatrix}.
\qquad
\left|\begin{matrix}1 & 8\\ 10 & 9\end{matrix}\right| = 6 \ne 0
\]

\[
\Rightarrow
\begin{pmatrix}
    1 & 8\\ 10 & 9
\end{pmatrix}
\text{ regular.}
\qquad
\overline{\begin{pmatrix}1 & 8\\ 10 & 9\end{pmatrix}}^{\,T} =
\begin{pmatrix}9 & 1\\ 3 & 1\end{pmatrix}^{\,T} =
\begin{pmatrix}9 & 3\\ 1 & 1\end{pmatrix}
\]

\[\Rightarrow
\begin{pmatrix}1 & 8\\ 10 & 9\end{pmatrix}^{-1} = 6^{-1} 
\begin{pmatrix}9 & 3\\ 1 & 1\end{pmatrix}.
\qquad
A = 2 
\begin{pmatrix}9 & 3\\ 1 & 1\end{pmatrix}
\begin{pmatrix}3 & 3\\ 1 & 6\end{pmatrix}
=
\begin{pmatrix}7 & 6\\ 2 & 2\end{pmatrix}
\begin{pmatrix}3 & 3\\ 1 & 6\end{pmatrix}
=
\boxed{ 
\begin{pmatrix}5 & 2\\ 8 & 7\end{pmatrix}}.
\]

\vspace{0.5em}

\textbf{b)}

Sabemos que $(A^4)^2 \cdot A = A^9$, entonces
\[
\begin{pmatrix}8 & 1\\ 2 & 0\end{pmatrix}^2 A = \begin{pmatrix}5 & 4\\ 8 & 6\end{pmatrix}.
\]

Calculamos primero el cuadrado:
\[
\begin{pmatrix}8 & 1\\ 2 & 0\end{pmatrix}^2 = 
\begin{pmatrix}8 & 1\\ 2 & 0\end{pmatrix} 
\begin{pmatrix}8 & 1\\ 2 & 0\end{pmatrix} =
\begin{pmatrix}0 & 8\\ 5 & 2\end{pmatrix}.
\]

Por tanto,
\[
\begin{pmatrix}0 & 8\\ 5 & 2\end{pmatrix} A = \begin{pmatrix}5 & 4\\ 8 & 6\end{pmatrix}
\quad \Rightarrow \quad 
A = \begin{pmatrix}0 & 8\\ 5 & 2\end{pmatrix}^{-1} \begin{pmatrix}5 & 4\\ 8 & 6\end{pmatrix}. \qquad
\left|\begin{matrix}0 & 8\\ 5 & 2\end{matrix}\right| = 4 \ne 0
\]

\[
\Rightarrow \quad
\begin{pmatrix}
    0 & 8\\ 5 & 2
\end{pmatrix}
\text{ regular.} \qquad
\overline{\begin{pmatrix}0 & 8\\ 5 & 2\end{pmatrix}}^{\,T} = 
\begin{pmatrix}2 & 6\\ 3 & 0\end{pmatrix}^{\,T} = 
\begin{pmatrix}2 & 3\\ 6 & 0\end{pmatrix}
\]

\[
\Rightarrow \quad 
\begin{pmatrix}0 & 8\\ 5 & 2\end{pmatrix}^{-1} = 4^{-1} 
\begin{pmatrix}2 & 3\\ 6 & 0\end{pmatrix}.
\quad
A = 3 
\begin{pmatrix}2 & 3\\ 6 & 0\end{pmatrix} 
\begin{pmatrix}5 & 4\\ 8 & 6\end{pmatrix} =
\begin{pmatrix}6 & 9\\ 7 & 0\end{pmatrix} 
\begin{pmatrix}5 & 4\\ 8 & 6\end{pmatrix} =
\boxed{
\begin{pmatrix}3 & 1\\ 2 & 6\end{pmatrix}}.
\]

\end{ejercicio}

\begin{ejercicio}[title=Ejercicio 24]
Calcule todas las soluciones del sistema siguiente, donde $X$ e $Y$ son matrices cuadradas de orden 2 con coeficientes en $\mathbb{R}$:
\[
\begin{cases}
\phantom{
\begin{pmatrix}
-1 & 1\\
1 & 2
\end{pmatrix}
}\,
X +
\,\,\,\,\,
\begin{pmatrix}
1 & 0\\
-1 & 1
\end{pmatrix}
Y =
\begin{pmatrix}
0 & -2\\
0 & 4
\end{pmatrix},\\[1em]
\begin{pmatrix}
-1 & 1\\
1 & 2
\end{pmatrix}
X +
\begin{pmatrix}
0 & -1\\
-1 & 0
\end{pmatrix}
Y =
\begin{pmatrix}
-2 & 2\\
0 & 4
\end{pmatrix}.
\end{cases}
\]

\tcblower
\textbf{Resolución:}

Usamos sustitución.

Despejamos \(X\) en la primera ecuación:
\[
X =
\begin{pmatrix}
0 & -2\\
0 & 4
\end{pmatrix}
-
\begin{pmatrix}
1 & 0\\
-1 & 1
\end{pmatrix} Y.
\]

Sustituimos \(X\) en la segunda ecuación:
\[
\begin{pmatrix}
-1 & 1\\
1 & 2
\end{pmatrix}
\Bigg[
\begin{pmatrix}
0 & -2\\
0 & 4
\end{pmatrix}
-
\begin{pmatrix}
1 & 0\\
-1 & 1
\end{pmatrix} Y
\Bigg]
+
\begin{pmatrix}
0 & -1\\
-1 & 0
\end{pmatrix} Y
=
\begin{pmatrix}
-2 & 2\\
0 & 4
\end{pmatrix}
\]

\[
\Rightarrow
\begin{pmatrix}
0 & 6\\
0 & 6
\end{pmatrix}
-
\begin{pmatrix}
-2 & 1\\
-1 & 2
\end{pmatrix} Y
+
\begin{pmatrix}
0 & -1\\
-1 & 0
\end{pmatrix} Y
=
\begin{pmatrix}
-2 & 2\\
0 & 4
\end{pmatrix}
\]

\[
\Rightarrow
\begin{pmatrix}
2 & -1\\
1 & -2
\end{pmatrix} Y
+
\begin{pmatrix}
0 & -1\\
-1 & 0
\end{pmatrix} Y
=
\begin{pmatrix}
-2 & -4\\
0 & -2
\end{pmatrix}
\]

\[
\Rightarrow
\begin{pmatrix}
2 & -2\\
0 & -2
\end{pmatrix} Y
=
\begin{pmatrix}
-2 & -4\\
0 & -2
\end{pmatrix}
\quad
\Rightarrow
Y = 
\begin{pmatrix}
2 & -2\\
0 & -2
\end{pmatrix}^{-1}
\begin{pmatrix}
-2 & -4\\
0 & -2
\end{pmatrix}.
\]
Calculamos la inversa:
\[
\left|
\begin{matrix}
2 & -2\\
0 & -2
\end{matrix}
\right|
= -4 \neq 0
\qquad \Rightarrow
\begin{pmatrix}
    2 & -2\\
0 & -2
\end{pmatrix}
\text{ regular.}
\]
Entonces:
\[
\overline{\begin{pmatrix}2 & -2\\ 0 & -2\end{pmatrix}}^T
=
\begin{pmatrix}-2 & 0\\ 2 & 2\end{pmatrix}^T
=
\begin{pmatrix}
-2 & 2\\
0 & 2
\end{pmatrix}
\quad
\Rightarrow
\begin{pmatrix}
2 & -2\\
0 & -2
\end{pmatrix}^{-1}
= -\frac{1}{4}
\begin{pmatrix}
-2 & 2\\
0 & 2
\end{pmatrix}.
\]
Ahora multiplicamos y obtenemos $Y$:
\[
Y = -\frac{1}{4}
\begin{pmatrix}
-2 & 2\\
0 & 2
\end{pmatrix}
\begin{pmatrix}
-2 & -4\\
0 & -2
\end{pmatrix}
= -\frac{1}{4}
\begin{pmatrix}
4 & 4\\
0 & -4
\end{pmatrix}
=
\boxed{
\begin{pmatrix}
-1 & -1\\[4pt]
0 & 1
\end{pmatrix}
}
\]
Sustituimos en la primera ecuación para obtener \(X\):
\[
X =
\begin{pmatrix}
0 & -2\\
0 & 4
\end{pmatrix}
-
\begin{pmatrix}
1 & 0\\
-1 & 1
\end{pmatrix}
\begin{pmatrix}
-1 & -1\\
0 & 1
\end{pmatrix}
=
\begin{pmatrix}
0 & -2\\
0 & 4
\end{pmatrix}
-
\begin{pmatrix}
-1 & -1\\
1 & 2
\end{pmatrix}
=
\boxed{
\begin{pmatrix}
1 & -1\\
-1 & 2
\end{pmatrix}
}
\]
\end{ejercicio}

\subsubsection{La fórmula de Cramer}

\begin{ejercicio}[title=Ejercicio 25]
Compruebe que cada uno de los sistemas siguientes sobre $\mathbb{R}$ es de Cramer y calcule sus soluciones aplicando la fórmula de Cramer:

\[
\begin{aligned}
\text{a)}\;
\begin{cases}
x + y + z = 1\\
x - y + z = 0\\
x + y - z = 2
\end{cases}
\qquad\qquad
\text{b)}\;
\begin{cases}
x + y - z = -1\\
2x + y - 3z = 1\\
x - 2y - 2z = 1
\end{cases}
\end{aligned}
\]

A continuación, obtenga la representación matricial de cada uno de los sistemas y resuélvalo usando dicha representación.

\tcblower
\textbf{Resolución:}

\textbf{Comprobación de Sistemas de Cramer}

En ambos sistemas hay 3 ecuaciones y 3 incógnitas.

\bigskip

Para \textbf{a)}:

\[
(A \,|\, B) =
\begin{pmatrix}
1 & 1 & 1 &\Big|& 1\\
1 & -1 & 1 &\Big|& 0\\
1 & 1 & -1 &\Big|& 2
\end{pmatrix}
\]

Entonces
\[
|A| = 1+1+1+1-1+1 = 4 \neq 0 \quad\Rightarrow\quad A \text{ regular.}
\]

\bigskip

Para \textbf{b)}:

\[
(A \,|\, B) =
\begin{pmatrix}
1 & 1 & -1 &\Big|& -1\\
2 & 1 & -3 &\Big|& 1\\
1 & -2 & -2 &\Big|& 1
\end{pmatrix}
\]

Entonces
\[
|A| = -2+4-3+1-6+4 = -2 \neq 0 \quad\Rightarrow\quad A \text{ regular.}
\]

\[
\boxed{\text{Ambos sistemas a) y b) son de Cramer.}}
\]

\bigskip

\textbf{Soluciones aplicando la fórmula de Cramer}

\medskip

\textbf{a)}

\[
x = \frac{|M_x|}{|A|} = \frac{1}{4}
\left|
\begin{matrix}
1 & 1 & 1\\
0 & -1 & 1\\
2 & 1 & -1
\end{matrix}
\right|
= \frac{1}{4}(1+2+2-1) = 1,
\]
\end{ejercicio}
\begin{ejercicio}[title=Ejercicio 25]
\[
y = \frac{|M_y|}{|A|}
= \frac{1}{4}
\left|
\begin{matrix}
1 & 1 & 1\\
1 & 0 & 1\\
1 & 2 & -1
\end{matrix}
\right|
= \frac{1}{4}(2+1-2+1) = \frac{1}{2},
\]
\[
z = \frac{|M_z|}{|A|}
= \frac{1}{4}
\left|
\begin{matrix}
1 & 1 & 1\\
1 & -1 & 0\\
1 & 1 & 2
\end{matrix}
\right|
= \frac{1}{4}(-2+1+1-2) = -\frac{1}{2}.
\]

\[
\boxed{\,x = 1,\quad y = \frac{1}{2},\quad z = -\frac{1}{2}\,}
\]

\bigskip

\textbf{b)}

\[
x = \frac{|M_x|}{|A|} = -\frac{1}{2}
\left|
\begin{matrix}
-1 & 1 & -1\\
1 & 1 & -3\\
1 & -2 & -2
\end{matrix}
\right|
= -\frac{1}{2}(2+2-3+1+6+2) = -5,
\]

\[
y = \frac{|M_y|}{|A|} = -\frac{1}{2}
\left|
\begin{matrix}
1 & -1 & -1\\
2 & 1 & -3\\
1 & 1 & -2
\end{matrix}
\right|
= -\frac{1}{2}(-2-2+3+1+3-4) = \frac{1}{2},
\]

\[
z = \frac{|M_z|}{|A|} = -\frac{1}{2}
\left|
\begin{matrix}
1 & 1 & -1\\
2 & 1 & 1\\
1 & -2 & 1
\end{matrix}
\right|
= -\frac{1}{2}(1+4+1+1+2-2) = -\frac{7}{2}.
\]

\[
\boxed{\,x = -5,\quad y = \frac{1}{2},\quad z = -\frac{7}{2}\,}
\]

\bigskip

\textbf{Resolución matricial}

La solución viene dada por la ecuación $X = A^{-1}B$. Es un SCD.

\textbf{a)}

\[
A =
\begin{pmatrix}
1 & 1 & 1\\
1 & -1 & 1\\
1 & 1 & -1
\end{pmatrix},
\qquad
\overline{A} =
\begin{pmatrix}
0 & 2 & 2\\
2 & -2 & 0\\
2 & 0 & -2
\end{pmatrix},
\qquad
\overline{A}^{\,T} =
\begin{pmatrix}
0 & 2 & 2\\
2 & -2 & 0\\
2 & 0 & -2
\end{pmatrix}.
\]

\[
A^{-1}
= \frac{1}{4}
\begin{pmatrix}
0 & 2 & 2\\
2 & -2 & 0\\
2 & 0 & -2
\end{pmatrix}
= \frac{1}{2}
\begin{pmatrix}
0 & 1 & 1\\
1 & -1 & 0\\
1 & 0 & -1
\end{pmatrix}.
\]

Entonces,

\[
X
= \frac{1}{2}
\begin{pmatrix}
0 & 1 & 1\\
1 & -1 & 0\\
1 & 0 & -1
\end{pmatrix}
\begin{pmatrix}
1\\ 0\\ 2
\end{pmatrix}
= \frac{1}{2}
\begin{pmatrix}
2\\ 1\\ -1
\end{pmatrix}
=
\begin{pmatrix}
1\\[2pt]
\tfrac12\\[2pt]
-\tfrac12
\end{pmatrix}
\]

\[
\boxed{\,x = 1,\quad y = \frac{1}{2},\quad z = -\frac{1}{2}\,}
\]
\end{ejercicio}
\begin{ejercicio}[title=Ejercicio 25]
\textbf{b)}

\[
A =
\begin{pmatrix}
1 & 1 & -1\\
2 & 1 & -3\\
1 & -2 & -2
\end{pmatrix},
\qquad
\overline{A} =
\begin{pmatrix}
-8 & 1 & -5\\
4 & -1 & 3\\
-2 & 1 & -1
\end{pmatrix},
\qquad
\overline{A}^{\,T} =
\begin{pmatrix}
-8 & 4 & -2\\
1 & -1 & 1\\
-5 & 3 & -1
\end{pmatrix}.
\]
\[
A^{-1}
= -\frac{1}{2}
\begin{pmatrix}
-8 & 4 & -2\\
1 & -1 & 1\\
-5 & 3 & -1
\end{pmatrix}.
\]

Entonces,

\[
X
=
-\frac{1}{2}
\begin{pmatrix}
-8 & 4 & -2\\
1 & -1 & 1\\
-5 & 3 & -1
\end{pmatrix}
\begin{pmatrix}
-1\\ 1\\ 1
\end{pmatrix}
=
-\frac{1}{2}
\begin{pmatrix}
10\\ -1\\ 7
\end{pmatrix} =
\begin{pmatrix}
-5\\[2pt]
\tfrac12\\[2pt]
-\tfrac72
\end{pmatrix}
\]

\[
\boxed{\,x = -5,\quad y = \frac{1}{2},\quad z = -\frac{7}{2}\,}
\]

\end{ejercicio}

\begin{ejercicio}[title=Ejercicio 26]
Resuelva cada uno de los sistemas siguientes sobre $\mathbb{Z}_7$:

\[
\begin{aligned}
\text{a)}&
\begin{cases}
x + y + 3z = 1\\
2x + 6y + 3z = 6\\
x + 4y + 6z = 4
\end{cases}
\qquad&
\text{b)}
\begin{cases}
3x + 4y = 2\\
2x + 5y = 4
\end{cases}
\\[1em]
\text{c)}&
\begin{cases}
4x + 5y = 2\\
3x + y = 6\\
6x + 4y = 3\\
4x + 3y = 4\\
4x + 2y = 5
\end{cases}
\qquad&\qquad
\text{d)}
\begin{cases}
x + 2y + 3z + t = 2\\
x + 6y + z + 2t = 1\\
x + 3y + 6z + 3t = 0\\
2x + 2y + 5t = 1
\end{cases}
\end{aligned}
\]

\tcblower
\textbf{Resolución:}

\textbf{a)}  

Vemos que nº ecuaciones = nº incógnitas = 3, así que comprobamos si es sistema de Cramer:

\[
\left|
\begin{matrix}
1 & 1 & 3\\
2 & 6 & 3\\
1 & 4 & 6
\end{matrix}
\right| = 0 \quad(\text{en } \mathbb{Z}_7),
\]
por tanto $A$ no es regular $\Rightarrow$ no es sistema de Cramer.

Aplicamos Gauss–Jordan:

\[
(A\,|\,B)=
\begin{pmatrix}
1 & 1 & 3 &\Big|& 1\\
2 & 6 & 3 &\Big|& 6\\
1 & 4 & 6 &\Big|& 4
\end{pmatrix}
\overset{\substack{F_2-2F_1\\F_3-F_1\\{}}}{\sim_f}
\begin{pmatrix}
1 & 1 & 3 &\Big|& 1\\
0 & 4 & 4 &\Big|& 4\\
0 & 3 & 3 &\Big|& 3
\end{pmatrix}
\overset{\substack{F_3+F_2\\{}}}{\sim_f}
\begin{pmatrix}
1 & 1 & 3 &\Big|& 1\\
0 & 4 & 4 &\Big|& 4\\
0 & 0 & 0 &\Big|& 0
\end{pmatrix}
\]

\[
\overset{\substack{F_1-2F_2\\{}}}{\sim_f}
\begin{pmatrix}
1 & 0 & 2 &\Big|& 0\\
0 & 4 & 4 &\Big|& 4\\
0 & 0 & 0 &\Big|& 0
\end{pmatrix}
\overset{\substack{2F_2\\{}}}{\sim_f}
\begin{pmatrix}
1 & 0 & 2 &\Big|& 0\\
0 & 1 & 1 &\Big|& 1\\
0 & 0 & 0 &\Big|& 0
\end{pmatrix}
= H_f(A\,|\,B)
\Rightarrow
\begin{cases}
x + 2z = 0\\[4pt]
y + z = 1
\end{cases}
\]

\vspace{0.5em}

$\operatorname{Rg}(A) = \operatorname{Rg}(A\,|\,B) \ne $ nº incógnitas $ \Rightarrow$ SCI. Tomando $z=\lambda\in\mathbb{Z}_7$ obtenemos:
\[
x = -2\lambda \equiv 5\lambda,\qquad y = 1-\lambda \equiv 1+6\lambda,\qquad z=\lambda.
\]
\[
\boxed{\,x = 5\lambda,\quad y = 1+6\lambda,\quad z = \lambda,\qquad \lambda\in\mathbb{Z}_7\,}
\]

\textbf{b)}

Vemos que nº ecuaciones = nº incógnitas = 2, así que vemos si es sistema de Cramer:

\[
\left|
\begin{matrix}
3 & 4\\
2 & 5
\end{matrix}
\right| = 0 \quad(\text{en } \mathbb{Z}_7),
\]
por tanto $A$ no es regular $\Rightarrow$ no es sistema de Cramer.
\end{ejercicio}
\begin{ejercicio}[title=Ejercicio 26]
Aplicamos Gauss–Jordan:

\[
(A\,|\,B)=
\begin{pmatrix}
3 & 4 &\Big|& 2\\
2 & 5 &\Big|& 4
\end{pmatrix}
\overset{\substack{5F_1\\4F_2\\{}}}{\sim_f}
\begin{pmatrix}
1 & 6 &\Big|& 3\\
1 & 6 &\Big|& 2
\end{pmatrix}
\overset{\substack{F_2 - F_1\\{}}}{\sim_f}
\begin{pmatrix}
1 & 6 &\Big|& 3\\
0 & 0 &\Big|& 6
\end{pmatrix}
\]

\[
\Rightarrow
\begin{cases}
x + 6y = 3,\\[4pt]
0 = 6.
\end{cases}
\Rightarrow
\text{SI}
\]

\[
\boxed{\text{El sistema no tiene solución.}}
\]

\bigskip

\textbf{c)}

Como nº ecuaciones = 5 $\ne$ nº incógnitas = 2, así que no podemos aplicar la fórmula de Cramer. Aplicamos Gauss-Jordan:

\[
(A\,|\,B)=
\begin{pmatrix}
4 & 5 &\Big|& 2\\
3 & 1 &\Big|& 6\\
6 & 4 &\Big|& 3\\
4 & 3 &\Big|& 4\\
4 & 2 &\Big|& 5
\end{pmatrix}
\overset{\substack{F_1+F_2\\F_3-2F_2\\F_4+F_2\\F_5+F_2\\{}}}{\sim_f}
\begin{pmatrix}
0 & 6 &\Big|& 1\\
3 & 1 &\Big|& 6\\
0 & 2 &\Big|& 5\\
0 & 4 &\Big|& 3\\
0 & 3 &\Big|& 4
\end{pmatrix}
\overset{\substack{5F_2\\{}}}{\sim_f}
\begin{pmatrix}
0 & 6 &\Big|& 1\\
1 & 5 &\Big|& 2\\
0 & 2 &\Big|& 5\\
0 & 4 &\Big|& 3\\
0 & 3 &\Big|& 4
\end{pmatrix}
\overset{\substack{F_1 \leftrightarrow F_2\\{}}}{\sim_f}
\begin{pmatrix}
1 & 5 &\Big|& 2\\
0 & 6 &\Big|& 1\\
0 & 2 &\Big|& 5\\
0 & 4 &\Big|& 3\\
0 & 3 &\Big|& 4
\end{pmatrix}
\]

\[
\overset{\substack{6F_2\\4F_3\\2F_4\\5F_5\\{}}}{\sim_f}
\begin{pmatrix}
1 & 5 &\Big|& 2\\
0 & 1 &\Big|& 6\\
0 & 1 &\Big|& 6\\
0 & 1 &\Big|& 6\\
0 & 1 &\Big|& 6
\end{pmatrix}
\overset{\substack{F_1 - 5F_2\\F_3 - F_2\\F_4 - F_2\\F_5 - F_2\\{}}}{\sim_f}
\begin{pmatrix}
1 & 0 &\Big|& 0\\
0 & 1 &\Big|& 6\\
0 & 0 &\Big|& 0\\
0 & 0 &\Big|& 0\\
0 & 0 &\Big|& 0
\end{pmatrix}
= H_f(A\,|\,B)
\Rightarrow
\begin{cases}
x = 0\\[4pt]
y = 6
\end{cases}
\quad \text{(SCD)}
\]

\vspace{0.5em}

\[
\boxed{\,x = 0,\quad y = 6\,}
\]

\vspace{0.5em}

\textbf{d)}

Vemos que nº ecuaciones = nº incógnitas = 4, pero aplicar la regla de Sarrus para órdenes $\ge 4$ no es buen método. Aplicamos Gauss–Jordan directamente:

\[
(A\,|\,B)=
\begin{pmatrix}
1 & 2 & 3 & 1 &\Big|& 2\\
1 & 6 & 1 & 2 &\Big|& 1\\
1 & 3 & 6 & 3 &\Big|& 0\\
2 & 2 & 0 & 5 &\Big|& 1
\end{pmatrix}
\overset{\substack{F_2 - F_1 \\ F_3 - F_1 \\ F_4 - 2F_1\\{}}}{\sim_f}
\begin{pmatrix}
1 & 2 & 3 & 1 &\Big|& 2\\
0 & 4 & 5 & 1 &\Big|& 6\\
0 & 1 & 3 & 2 &\Big|& 5\\
0 & 5 & 1 & 3 &\Big|& 4
\end{pmatrix}
\]
\end{ejercicio}
\begin{ejercicio}[title=Ejercicio 26]
\[
\overset{\substack{2F_2 \\ 3F_4 \\ {}}}{\sim_f}
\begin{pmatrix}
1 & 2 & 3 & 1 &\Big|& 2\\
0 & 1 & 3 & 2 &\Big|& 5\\
0 & 1 & 3 & 2 &\Big|& 5\\
0 & 1 & 3 & 2 &\Big|& 5
\end{pmatrix}
\overset{\substack{F_1 - 2F_2 \\ F_3 - F_2 \\ F_4 - F_2 \\{}}}{\sim_f}
\begin{pmatrix}
1 & 0 & 4 & 4 &\Big|& 6\\
0 & 1 & 3 & 2 &\Big|& 5\\
0 & 0 & 0 & 0 &\Big|& 0\\
0 & 0 & 0 & 0 &\Big|& 0
\end{pmatrix}
= H_f(A\,|\,B)
\Rightarrow
\begin{cases}
x + 4z + 4t = 6\\[4pt]
y + 3z + 2t = 5
\end{cases}
\]

$\operatorname{Rg}(A) = \operatorname{Rg}(A\,|\,B) \ne $ nº incógnitas $ \Rightarrow$ SCI.

\vspace{0.5em}

Tomando parámetros libres \(z=\lambda_1,\ t=\lambda_2\in\mathbb{Z}_7\) se tiene
\[
x = 6 - 4z - 4t \equiv 6 + 3z + 3t,
\qquad
y = 5 - 3z - 2t \equiv 5 + 4z + 5t.
\]
\[
\boxed{\,x = 6+3\lambda_1+3\lambda_2,\quad y = 5+4\lambda_1+5\lambda_2,\quad z=\lambda_1,\quad t=\lambda_2,\quad
\lambda_1,\lambda_2\in\mathbb{Z}_7\,}
\]
\end{ejercicio}

\begin{ejercicio}[title=Ejercicio 30]
Discuta y resuelva los sistemas siguientes sobre $\mathbb{Z}_7$ según los valores del parámetro $a$:

\begin{enumerate}
\item[a)] 
$\begin{cases}
x + 3y + (a + 1)z = 2\\
x + 4y + (2a + 2)z = 5\\
x + 2y + 2z = a^2
\end{cases}$

\item[b)] 
$\begin{cases}
x + 2y + 3z = 2\\
2x + y + a z = 3\\
4x + y + a z = 6
\end{cases}$

\item[c)] 
$\begin{cases}
5x + 2y + 6z = 1\\
2x + y + a z = 4\\
3x + 4y + 5z = a
\end{cases}$
\end{enumerate}

\tcblower
\textbf{Resolución:}

Usamos el teorema de Rouché–Frobenius para averiguar la compatibilidad del sistema. Esto consiste en analizar los valores de los rangos de $A$ y $(A\,|\,B)$. Calculamos una matriz en forma escalonada por filas en cada apartado.

\textbf{a)}

\[
(A\,|\,B)=
\begin{pmatrix}
1 & 3 & a+1 &\Big|& 2\\
1 & 4 & 2a+2 &\Big|& 5\\
1 & 2 & 2 &\Big|& a^2
\end{pmatrix}
\overset{\substack{F_2 - F_1 \\ F_3 - F_1}}{\sim_f}
\begin{pmatrix}
1 & 3 & a+1 &\Big|& 2\\
0 & 1 & a+1 &\Big|& 3\\
0 & 6 & 6a+1 &\Big|& a^2-2
\end{pmatrix}
\]

\[
\overset{\substack{F_3 + F_2\\{}}}{\sim_f}
\begin{pmatrix}
1 & 3 & a+1 &\Big|& 2\\
0 & 1 & a+1 &\Big|& 3\\
0 & 0 & 2 &\Big|& a^2+1
\end{pmatrix}
\text{ está en forma escalonada por filas.}
\]

De aquí:
\[
\operatorname{Rg}(A)=\text{nº filas no nulas }=3,\qquad
\operatorname{Rg}(A\,|\,B)=\text{nº filas no nulas }=3,
\]
porque en ambos casos hay elementos no nulos en cada una de las filas.

\vspace{0.5em}

Como \(\operatorname{Rg}(A)=\operatorname{Rg}(A\,|\,B)=3=\text{nº incógnitas}\), por el teorema de Rouché–Frobenius:

\[\boxed{\text{El sistema es un SCD.}}\]

\vspace{0.5em}

Para resolver el sistema, comprobamos que \(A\) es regular:
\[
\left|
\begin{matrix}
1 & 3 & a+1\\
1 & 4 & 2a+2\\
1 & 2 & 2
\end{matrix}
\right| = 2 \neq 0 \quad(\text{en }\mathbb{Z}_7),
\]
por tanto \(A\) es regular y es un sistema de Cramer.  
\end{ejercicio}
\begin{ejercicio}[title=Ejercicio 30]
Como \(2^{-1}=4\) en \(\mathbb{Z}_7\), aplicamos la fórmula de Cramer.

\[
x = 4\,
\left|
\begin{matrix}
2 & 3 & a+1\\
5 & 4 & 2a+2\\
a^2 & 2 & 2
\end{matrix}
\right|
= a^3 + a^2 + a + 1,
\]

\[
y = 4\,
\left|
\begin{matrix}
1 & 2 & a+1\\
1 & 5 & 2a+2\\
1 & a^2 & 2
\end{matrix}
\right|
= 3a^3 + 3a^2 + 3a + 6,
\]

\[
z = 4\,
\left|
\begin{matrix}
1 & 3 & 2\\
1 & 4 & 5\\
1 & 2 & a^2
\end{matrix}
\right|
= 4(a^2+1)
= 4a^2 + 4.
\]

Por tanto las soluciones son

\[
\boxed{\,x = a^3 + a^2 + a + 1,\qquad
y = 3a^3 + 3a^2 + 3a + 6,\qquad
z = 4a^2 + 4\,}
\]

\vspace{1em}

\textbf{b)}

La matriz ampliada es
\[
(A\,|\,B)=
\begin{pmatrix}
1 & 2 & 3 &\Big|& 2\\
2 & 1 & a &\Big|& 3\\
4 & 1 & a &\Big|& 6
\end{pmatrix}
\overset{\substack{F_2+5F_1\\ F_3+3F_1\\{}}}{\sim_f}
\begin{pmatrix}
1 & 2 & 3 &\Big|& 2\\
0 & 4 & a+1 &\Big|& 6\\
0 & 0 & a+2 &\Big|& 5
\end{pmatrix}
\text{\shortstack{ está en forma\\ escalonada por filas.}}
\]

De aquí,
\[
\operatorname{Rg}(A\,|\,B)=\text{nº filas no nulas}=3.
\]

Para la matriz de coeficientes \(A\), la tercera fila es nula si y sólo si
\(a+2=0\), es decir, \(a=5\) (en \(\mathbb{Z}_7\)). Por tanto,
\[
\operatorname{Rg}(A)=
\begin{cases}
2 & \text{si } a=5,\\
3 & \text{si } a\neq 5.
\end{cases}
\]

\begin{itemize}
\item \textbf{Caso 1. \(a=5\).}

\[
(A\,|\,B)\sim_f
\begin{pmatrix}
1 & 2 & 3 &\Big|& 2\\
0 & 4 & 6 &\Big|& 6\\
0 & 0 & 0 &\Big|& 5
\end{pmatrix}.
\]

Se tiene \(\operatorname{Rg}(A)=2 \neq \operatorname{Rg}(A\,|\,B)=3\).  
\[
\boxed{\text{El sistema es un SI si } a=5.}
\]
\end{itemize}
\end{ejercicio}
\begin{ejercicio}[title=Ejercicio 30]
\begin{itemize}
\item \textbf{Caso 2. \(a\neq 5\).}

\[
(A\,|\,B)\sim_f
\begin{pmatrix}
1 & 2 & 3 &\Big|& 2\\
0 & 4 & a+1 &\Big|& 6\\
0 & 0 & a+2 &\Big|& 5
\end{pmatrix}.
\]

En este caso,
\[
\operatorname{Rg}(A)=\operatorname{Rg}(A\,|\,B)=\text{nº incógnitas}=3,
\]
luego
\[
\boxed{\text{El sistema es un SCD si } a\neq 5.}
\]

Como el sistema tiene el mismo número de ecuaciones que de incógnitas, es un
sistema de Cramer. El sistema equivalente tiene las mismas soluciones que el original.
\[
|A|=
\begin{vmatrix}
1 & 2 & 3\\
0 & 4 & a+1\\
0 & 0 & a+2
\end{vmatrix}
=4a+1
\]

Aplicamos la fórmula de Cramer.

\[
x=(4a+1)^{-1}
\begin{vmatrix}
2 & 2 & 3\\
6 & 4 & a+1\\
5 & 0 & a+2
\end{vmatrix}
=(4a+1)^{-1}(6a+5)
=\boxed{(6a+5)(4a+1)^{-1}}.
\]

\[
y=(4a+1)^{-1}
\begin{vmatrix}
1 & 2 & 3\\
0 & 6 & a+1\\
0 & 5 & a+2
\end{vmatrix}
=(4a+1)^{-1}\,a
=\boxed{a(4a+1)^{-1}}.
\]

\[
z=(4a+1)^{-1}
\begin{vmatrix}
1 & 2 & 2\\
0 & 4 & 6\\
0 & 0 & 5
\end{vmatrix}
=\boxed{6(4a+1)^{-1}}.
\]
\end{itemize}

\vspace{1em}

\textbf{c)}

\[
(A\,|\,B)=
\begin{pmatrix}
5 & 2 & 6 &\Big|& 1\\
2 & 1 & a &\Big|& 4\\
3 & 4 & 5 &\Big|& a
\end{pmatrix}
\overset{\substack{F_2+F_1\\F_3+5F_1\\{}}}{\sim_f}
\begin{pmatrix}
5 & 2 & 6 &\Big|& 1\\
0 & 3 & a+6 &\Big|& 5\\
0 & 0 & 0 &\Big|& a+5
\end{pmatrix}
\text{\shortstack{ está en forma\\ escalonada por filas.}}
\]

\[
\operatorname{Rg}(A)=\text{nº filas no nulas }=2.
\]
\end{ejercicio}
\begin{ejercicio}[title=Ejercicio 30]
Para la matriz ampliada \((A\,|\,B)\), la tercera fila se anula si y sólo si \(a+5 = 0\), es decir \(a=2\). Entonces
\[
\operatorname{Rg}(A\,|\,B)=
\begin{cases}
2 & \text{si } a=2,\\[4pt]
3 & \text{si } a\neq 2.
\end{cases}
\]

\vspace{0.5em}

Teniendo en cuenta que nº incógnitas = 3, por el teorema de Rouché–Frobenius se obtiene:
\[
\boxed{\;
\text{El sistema es }\begin{cases}
\text{SCI} & \text{si } a=2,\\[4pt]
\text{SI} & \text{si } a\neq 2.
\end{cases}
}
\]

\bigskip

\noindent Para resolver el caso \(a=2\) aplicamos Gauss–Jordan sobre \((A\,|\,B)\) con \(a=2\):

\[
\begin{pmatrix}
5 & 2 & 6 &\Big|& 1\\
0 & 3 & 1 &\Big|& 5\\
0 & 0 & 0 &\Big|& 0
\end{pmatrix}
\overset{\substack{3F_1\\5F_2\\{}}}{\sim_f}
\begin{pmatrix}
1 & 6 & 4 &\Big|& 3\\
0 & 1 & 5 &\Big|& 4\\
0 & 0 & 0 &\Big|& 0
\end{pmatrix}
\overset{\substack{F_1+F_2\\{}}}{\sim_f}
\begin{pmatrix}
1 & 0 & 2 &\Big|& 0\\
0 & 1 & 5 &\Big|& 4\\
0 & 0 & 0 &\Big|& 0
\end{pmatrix}
= H_f(A\,|\,B).
\]

\[
\begin{cases}
x + 2z = 0,\\[4pt]
y + 5z = 4.
\end{cases}
\]

Tomando \(z=\lambda\in\mathbb{Z}_7\) (parámetro libre) se obtiene
\[
x = -2\lambda \equiv 5\lambda,\qquad
y = 4 - 5\lambda \equiv 4 + 2\lambda,\qquad
z = \lambda.
\]
\[
\boxed{\,x = 5\lambda,\quad y = 4+2\lambda,\quad z=\lambda,\qquad \lambda\in\mathbb{Z}_7\,}
\]
\end{ejercicio}

\begin{ejercicio}[title=Ejercicio 31]
Discuta y resuelva los sistemas siguientes sobre $\mathbb{R}$ según los valores de $a$ y $b$:

\[
\begin{tabular}{c @{\hspace{2cm}} c}
$a)\begin{cases}
a x + y + z = 1\\
x + a y + z = 1\\
x + y + a z = b
\end{cases}$ &
$b)\begin{cases}
x + b y + a z = 1\\
a x + b y + z = a\\
x + a b y + z = b
\end{cases}$
\end{tabular}
\]

\tcblower
\textbf{Resolución:}

La discusión de la compatibilidad de los sistemas en este ejercicio se estudia constantemente mediante el \textbf{Teorema de Rouché-Frobenius}, el cual determina la compatiblidad del sistema a partir del rango de la matriz de coeficientes y de la matriz ampliada.

\vspace{0.5em}

\textbf{a)}
En primer lugar, calculamos el determinante de la matriz de coeficientes, $|{A}|$, y vemos cuándo se anula para determinar cuándo podemos o no resolverlo por la fórmula de Cramer.
La matriz ampliada es:
$$({A}\,|\,{B}) =
\begin{pmatrix}
a & 1 & 1 &\Big|& 1\\
1 & a & 1 &\Big|& 1\\
1 & 1 & a &\Big|& b
\end{pmatrix}$$
Entonces, el determinante de ${A}$ es:
$$|{A}| =
\begin{vmatrix}
a & 1 & 1\\
1 & a & 1\\
1 & 1 & a
\end{vmatrix}
= a^3 - 3a + 2$$

Vemos que el número de ecuaciones es igual al número de incógnitas, que es 3.

El sistema es de Cramer si ${A}$ es regular, es decir, si $|{A}| \neq 0$.

Vemos qué valores anulan el determinante. Probamos con $a=1$ en el polinomio con Ruffini:
$$
\begin{array}{c|cccc}
  & 1 & 0 & -3 & 2 \\
1 &   & 1 & 1 & -2 \\
\hline
  & 1 & 1 & -2 & \boxed{0} \\
1 &   & 1 & 2 \\
\hline
  & 1 & 2 & \boxed{0}
\end{array}
$$

$$a^3 - 3a + 2 = (a-1)^2(a+2)$$

El determinante se anula si $a=1$ o $a=-2$.

\begin{itemize}
    \item \textbf{Caso 1. $a \neq 1$ y $a \neq -2$.}
    
Al no anularse el determinante $|{A}|$, la matriz ${A}$ es regular y el sistema es de Cramer. Además:
\[
\text{Rg}({A}) = \text{Rg}({A}|{B}) = \text{nº incógnitas} = 3
\]
\[
\boxed{\text{El sistema es un SCD si } a \ne 1 \text{ y } a \ne -2.}
\]
\end{itemize}
\end{ejercicio}
\begin{ejercicio}[title=Ejercicio 31]
\begin{itemize}
Para resolver el sistema, es más sencillo en este caso aplicar la fórmula ${X} = {A}^{-1} {B}$ que la fórmula de Cramer. Así que calculamos $A^{-1}$.
$$\overline{A} =
\begin{pmatrix}
a^2-1 & -a+1 & -a+1\\
-a+1 & a^2-1 & -a+1\\
-a+1 & -a+1 & a^2-1
\end{pmatrix}$$
$$\overline{{A}}^T =
\begin{pmatrix}
a^2-1 & -a+1 & -a+1\\
-a+1 & a^2-1 & -a+1\\
-a+1 & -a+1 & a^2-1
\end{pmatrix} = (a-1)
\begin{pmatrix}
a+1 & -1 & -1\\
-1 & a+1 & -1\\
-1 & -1 & a+1
\end{pmatrix}$$
\[\]

La inversa $A^{-1}$ es:

\[
A^{-1}
= \frac{1}{|A|}\,\overline{A}^{\,T}
= \frac{a-1}{(a-1)^2(a+2)}
\begin{pmatrix}
a+1 & -1 & -1\\
-1 & a+1 & -1\\
-1 & -1 & a+1
\end{pmatrix}
\]

\[
= \frac{1}{(a-1)(a+2)}
\begin{pmatrix}
a+1 & -1 & -1\\
-1 & a+1 & -1\\
-1 & -1 & a+1
\end{pmatrix}
= \frac{1}{a^2 + a - 2}
\begin{pmatrix}
a+1 & -1 & -1\\
-1 & a+1 & -1\\
-1 & -1 & a+1
\end{pmatrix}.
\]
\[\]

Entonces, la solución ${X} = {A}^{-1} {B}$ es:

\[
\begin{pmatrix}
x\\
y\\
z
\end{pmatrix}
= \frac{1}{a^2+a-2}
\begin{pmatrix}
a+1 & -1 & -1\\
-1 & a+1 & -1\\
-1 & -1 & a+1
\end{pmatrix}
\begin{pmatrix}
1\\
1\\
b
\end{pmatrix}
\]

\[
= \frac{1}{a^2+a-2}
\begin{pmatrix}
(a+1) - 1 - b\\
-1 + (a+1) - b\\
-1 - 1 + b(a+1)
\end{pmatrix}
= \frac{1}{a^2+a-2}
\begin{pmatrix}
a-b\\
a-b\\
ab+b-2
\end{pmatrix}
=
\boxed{
\begin{pmatrix}
\dfrac{a-b}{a^2+a-2}\\[12pt]
\dfrac{a-b}{a^2+a-2}\\[12pt]
\dfrac{ab+b-2}{a^2+a-2}
\end{pmatrix}}
\]

\item \textbf{Caso 2. $a=1$.}

La matriz ampliada es
\[
(A\,|\,B)=
\begin{pmatrix}
1 & 1 & 1 &\Big|& 1\\
1 & 1 & 1 &\Big|& 1\\
1 & 1 & 1 &\Big|& b
\end{pmatrix}
\overset{\substack{F_2-F_1\\ F_3-F_1 \\{}}}{\sim_f}
\begin{pmatrix}
1 & 1 & 1 &\Big|& 1\\
0 & 0 & 0 &\Big|& 0\\
0 & 0 & 0 &\Big|& b-1
\end{pmatrix}
\overset{F_2 \leftrightarrow F_3 \\{}}{\sim_f}
\begin{pmatrix}
1 & 1 & 1 &\Big|& 1\\
0 & 0 & 0 &\Big|& b-1\\
0 & 0 & 0 &\Big|& 0
\end{pmatrix},
\]
\[
\text{que está en forma escalonada por filas.}
\]
\end{itemize}
\end{ejercicio}
\begin{ejercicio}[title=Ejercicio 31]
\begin{itemize}
\[
\operatorname{Rg}(A)=\text{nº filas no nulas}=1,
\qquad
\operatorname{Rg}(A\,|\,B)=
\begin{cases}
1 & \text{si } b=1,\\[2mm]
2 & \text{si } b\neq 1.
\end{cases}
\]

\begin{itemize}
\item \textbf{Subcaso 2.1. $a=1$ y $b=1$.}
\[
(A\,|\,B)\sim_f
\begin{pmatrix}
1 & 1 & 1 &\Big|& 1\\
0 & 0 & 0 &\Big|& 0\\
0 & 0 & 0 &\Big|& 0
\end{pmatrix}
\text{ está en forma escalonada por filas.}
\]

Se tiene:
\[
\operatorname{Rg}(A)=\operatorname{Rg}(A\,|\,B)=1 \neq \text{nº incógnitas}=3.
\]
\[
\boxed{\text{El sistema es un SCI si } a=1 \text{ y } b=1.}
\]

Resolvemos:
\[
\begin{cases}
    x+y+z=1
\end{cases}
\]

Tomamos \(y=\lambda_1\), \(z=\lambda_2\) parámetros libres. Entonces:
\[
x=1-\lambda_1-\lambda_2.
\]
\[
\boxed{
x = 1-\lambda_1-\lambda_2,\qquad
y=\lambda_1,\qquad
z=\lambda_2,\qquad
\lambda_1,\lambda_2\in\mathbb{R}.
}
\]

\vspace{0.5em}

\item \textbf{Subcaso 2.2. $a=1$ y $b\neq 1$.}

En este caso:
\[
\operatorname{Rg}(A)=1 \neq \operatorname{Rg}(A\,|\,B)=2.
\]
\[
\boxed{\text{El sistema es un SI si } a=1 \text{ y } b\neq 1.}
\]

\end{itemize}

\item \textbf{Caso 3. \(a=-2\).}

\[
(A\,|\,B)=
\begin{pmatrix}
-2 & 1 & 1 &\Big|& 1\\
1 & -2 & 1 &\Big|& 1\\
1 & 1 & -2 &\Big|& b
\end{pmatrix}
\overset{\substack{F_1 \leftrightarrow F_3\\{}}}{\sim_f}
\begin{pmatrix}
1 & 1 & -2 &\Big|& b\\
1 & -2 & 1 &\Big|& 1\\
-2 & 1 & 1 &\Big|& 1
\end{pmatrix}
\]

\[
\overset{\substack{F_2 - F_1 \\ F_3 + 2F_1\\{}}}{\sim_f}
\begin{pmatrix}
1 & 1 & -2 &\Big|& b\\
0 & -3 & 3 &\Big|& 1-b\\
0 & 3 & -3 &\Big|& 2b+1
\end{pmatrix}
\overset{\substack{F_3 + F_2\\{}}}{\sim_f}
\begin{pmatrix}
1 & 1 & -2 &\Big|& b\\
0 & -3 & 3 &\Big|& -b+1\\
0 & 0 & 0 &\Big|& b+2
\end{pmatrix},
\]

que está en forma escalonada por filas.
\[
\operatorname{Rg}(A)=\text{nº filas no nulas}=2.
\]
\end{itemize}
\end{ejercicio}
\begin{ejercicio}[title=Ejercicio 31]
\begin{itemize}
Para la matriz aumentada \((A\,|\,B)\) la tercera fila se anula si y sólo si \(b+2=0\), es decir \(b=-2\). Por tanto
\[
\operatorname{Rg}(A\,|\,B)=
\begin{cases}
2 & \text{si } b=-2,\\[4pt]
3 & \text{si } b\neq -2.
\end{cases}
\]

\begin{itemize}
\item \textbf{Subcaso 3.1. \(a=-2\) y \(b=-2\).}

En este caso \(\operatorname{Rg}(A)=\operatorname{Rg}(A\,|\,B)=2 \neq\) nº incógnitas \(=3\).  
\[
\boxed{\text{El sistema es un SCI si } a=-2 \text{ y } b=-2.}
\]

Resolvemos aplicando Gauss–Jordan sobre la ampliada con \(b=-2\):
\[
(A\,|\,B)
\sim_f
\begin{pmatrix}
1 & 1 & -2 &\Big|& -2\\
0 & -3 & 3 &\Big|& 3\\
0 & 0 & 0 &\Big|& 0
\end{pmatrix}
\overset{\substack{-\tfrac{1}{3}F_2\\{}}}{\sim_f}
\begin{pmatrix}
1 & 1 & -2 &\Big|& -2\\
0 & 1 & -1 &\Big|& -1\\
0 & 0 & 0 &\Big|& 0
\end{pmatrix}
\]
\[
\overset{\substack{F_1 - F_2\\{}}}{\sim_f}
\begin{pmatrix}
1 & 0 & -1 &\Big|& -1\\
0 & 1 & -1 &\Big|& -1\\
0 & 0 & 0 &\Big|& 0
\end{pmatrix} = H_f(A\,|\,B) \Rightarrow
\begin{cases}
x - z = -1,\\[4pt]
y - z = -1.
\end{cases}
\]

Tomando \(z=\lambda\in\mathbb{R}\) (parámetro libre) obtenemos
\[
\boxed{\,x = -1+\lambda,\quad y = -1+\lambda,\quad z=\lambda,\qquad \lambda\in\mathbb{R}\,}
\]

\bigskip

\item \textbf{Subcaso 3.2. \(a=-2\) y \(b\neq -2\).}

Aquí \(\operatorname{Rg}(A)=2 \neq \operatorname{Rg}(A\,|\,B)=3\), por lo que el sistema es incompatible.

\[
\boxed{\text{El sistema es un SI si } a=-2 \text{ y } b\neq -2.}
\]
\end{itemize}
\end{itemize}

\textbf{b)}
En primer lugar, calculamos el determinante de la matriz de coeficientes, $|{A}|$, y vemos cuándo se anula para determinar cuándo podemos o no resolverlo por la fórmula de Cramer.
La matriz ampliada es:
$$({A}\,|\,{B}) =
\begin{pmatrix}
1 & b & a &\Big|& 1\\
a & b & 1 &\Big|& a\\
1 & ab & 1 &\Big|& b
\end{pmatrix}$$
\end{ejercicio}
\begin{ejercicio}[title=Ejercicio 31]
Entonces, el determinante de ${A}$ es:
$$|{A}| =
\begin{vmatrix}
1 & b & a\\
a & b & 1\\
1 & ab & 1
\end{vmatrix}
\begin{aligned}
& = b+a^3b+b-ab-ab-ab \\
&= (a^3 - 3a + 2)b
\end{aligned}
$$

Vemos que el número de ecuaciones es igual al número de incógnitas, que es 3.

El sistema es de Cramer si ${A}$ es regular, es decir, si $|{A}| \neq 0$.

$b=0$ anula el determinante. Vemos qué valores de $a$ anulan el determinante. Probamos con $a=1$ en el polinomio con Ruffini:
$$
\begin{array}{c|cccc}
  & 1 & 0 & -3 & 2 \\
1 &   & 1 & 1 & -2 \\
\hline
  & 1 & 1 & -2 & \boxed{0} \\
1 &   & 1 & 2 \\
\hline
  & 1 & 2 & \boxed{0}
\end{array}
$$

$$|A| = (a^3 - 3a + 2)b = (a-1)^2(a+2)b$$

El determinante se anula si $a=1$, $a=-2$ o $b=0$.

\begin{itemize}
    \item \textbf{Caso 1. $a \neq 1$, $a \neq -2$ y $b \neq 0$}
    
Al no anularse el determinante $|{A}|$, la matriz ${A}$ es regular y el sistema es de Cramer. Además:
\[
\text{Rg}({A}) = \text{Rg}({A}|{B}) = \text{nº incógnitas} = 3
\]
\[
\boxed{\text{El sistema es un SCD si } a \ne 1 \text{, } a \ne -2 \text{ y } b \ne 0.}
\]

Para resolver el sistema, es más sencillo en este caso aplicar la fórmula ${X} = {A}^{-1} {B}$ que la fórmula de Cramer. Así que calculamos $A^{-1}$.
$$\overline{A} =
\begin{pmatrix}
-ab+b & -a+1 & a^2b-b\\
a^2b-b & -a+1 & -ab+b\\
-ab+b & a-1 & -ab+b
\end{pmatrix}$$
$$\overline{{A}}^T =
\begin{pmatrix}
-ab+b & a^2b-b & -ab+b\\
-a+1 & -a+1 & a-1\\
a^2b-b & -ab+b & -ab+b
\end{pmatrix} = (a-1)
\begin{pmatrix}
-b & (a+1)b & -b\\
-1 & -1 & 1\\
(a+1)b & -b & -b
\end{pmatrix}$$
\[\]

La inversa $A^{-1}$ es:

\[
A^{-1}
= \frac{1}{|A|}\,\overline{A}^{\,T}
= \frac{a-1}{(a-1)^2(a+2)b}
\begin{pmatrix}
-b & (a+1)b & -b\\
-1 & -1 & 1\\
(a+1)b & -b & -b
\end{pmatrix}
\]
\end{itemize}
\end{ejercicio}
\begin{ejercicio}[title=Ejercicio 31]
\begin{itemize}

\[
= \frac{1}{(a-1)(a+2)b}
\begin{pmatrix}
-b & (a+1)b & -b\\
-1 & -1 & 1\\
(a+1)b & -b & -b
\end{pmatrix}
\]
\[
= \frac{1}{(a^2 + a - 2)b}
\begin{pmatrix}
-b & (a+1)b & -b\\
-1 & -1 & 1\\
(a+1)b & -b & -b
\end{pmatrix}.
\]
\[\]

Entonces, la solución ${X} = {A}^{-1} {B}$ es:

\[
\begin{pmatrix}
x\\
y\\
z
\end{pmatrix}
= \frac{1}{(a^2+a-2)b}
\begin{pmatrix}
-b & (a+1)b & -b\\
-1 & -1 & 1\\
(a+1)b & -b & -b
\end{pmatrix}
\begin{pmatrix}
1\\
a\\
b
\end{pmatrix}
\]

\[
= \frac{1}{(a^2+a-2)b}
\begin{pmatrix}
-b + (a+1)ab -b^2\\
-1 -a +b\\
(a+1)b -ab -b^2
\end{pmatrix}
= \frac{1}{(a^2+a-2)b}
\begin{pmatrix}
a^2b+ab-b^2-b\\
-a +b + 1\\
-b^2 +b
\end{pmatrix}
\]
\[
\begin{pmatrix}
x\\
y\\
z
\end{pmatrix}
=
\begin{pmatrix}
\dfrac{a^2+a-b-1}{(a^2+a-2)b}\\[12pt]
\dfrac{-a+b-1}{(a^2+a-2)b}\\[12pt]
\dfrac{-b+1}{a^2+a-2}
\end{pmatrix}
=
\boxed{
\begin{pmatrix}
\dfrac{a^2+a-b-1}{a^2b+ab-2b}\\[12pt]
\dfrac{-a+b-1}{a^2b+ab-2b}\\[12pt]
\dfrac{-b+1}{a^2+a-2}
\end{pmatrix}}
\]

\item \textbf{Caso 2. $a=1$.}

La matriz ampliada es
\[
(A\,|\,B)=
\begin{pmatrix}
1 & b & 1 &\Big|& 1\\
1 & b & 1 &\Big|& 1\\
1 & b & 1 &\Big|& b
\end{pmatrix}
\overset{\substack{F_2-F_1\\ F_3-F_1 \\{}}}{\sim_f}
\begin{pmatrix}
1 & b & 1 &\Big|& 1\\
0 & 0 & 0 &\Big|& 0\\
0 & 0 & 0 &\Big|& b-1
\end{pmatrix}
\overset{\substack{F_2 \leftrightarrow F_3 \\{}}}{\sim_f}
\begin{pmatrix}
1 & b & 1 &\Big|& 1\\
0 & 0 & 0 &\Big|& b-1\\
0 & 0 & 0 &\Big|& 0
\end{pmatrix},
\]
\[
\text{que está en forma escalonada por filas.}
\]
\[
\operatorname{Rg}(A)=\text{nº filas no nulas}=1,
\qquad
\operatorname{Rg}(A\,|\,B)=
\begin{cases}
1 & \text{si } b=1,\\[2mm]
2 & \text{si } b\neq 1.
\end{cases}
\]

\begin{itemize}
\item \textbf{Subcaso 2.1. $a=1$ y $b=1$.}
\[
(A\,|\,B)\sim_f
\begin{pmatrix}
1 & 1 & 1 &\Big|& 1\\
0 & 0 & 0 &\Big|& 0\\
0 & 0 & 0 &\Big|& 0
\end{pmatrix}
\text{ está en forma escalonada por filas.}
\]
\end{itemize}
\end{itemize}
\end{ejercicio}
\begin{ejercicio}[title=Ejercicio 31]
\begin{itemize}
\begin{itemize}
Se tiene:
\[
\operatorname{Rg}(A)=\operatorname{Rg}(A\,|\,B)=1 \neq \text{nº incógnitas}=3.
\]
\[
\boxed{\text{El sistema es un SCI si } a=1 \text{ y } b=1.}
\]

Resolvemos:
\[
\begin{cases}
    x+y+z=1
\end{cases}
\]

Tomamos \(y=\lambda_1\), \(z=\lambda_2\) parámetros libres. Entonces:
\[
x=1-\lambda_1-\lambda_2.
\]
\[
\boxed{
x = 1-\lambda_1-\lambda_2,\qquad
y=\lambda_1,\qquad
z=\lambda_2,\qquad
\lambda_1,\lambda_2\in\mathbb{R}.
}
\]

\item \textbf{Subcaso 2.2. $a=1$ y $b\neq 1$.}

En este caso:
\[
\operatorname{Rg}(A)=1 \neq \operatorname{Rg}(A\,|\,B)=2.
\]
\[
\boxed{\text{El sistema es un SI si } a=1 \text{ y } b\neq 1.}
\]

\end{itemize}

\item \textbf{Caso 3. \(a=-2\).}

\[
(A\,|\,B)=
\begin{pmatrix}
1 & b & -2 &\Big|& 1\\
-2 & b & 1 &\Big|& -2\\
1 & -2b & 1 &\Big|& b
\end{pmatrix}
\overset{\substack{F_2+2F_1\\F_3-F_1\\{}}}{\sim_f}
\begin{pmatrix}
1 & b & -2 &\Big|& 1\\
0 & 3b & -3 &\Big|& 0\\
0 & -3b & 3 &\Big|& b-1
\end{pmatrix}
\]
\[
\overset{\substack{F_3 + F_2\\{}}}{\sim_f}
\begin{pmatrix}
1 & b & -2 &\Big|& 1\\
0 & 3b & -3 &\Big|& 0\\
0 & 0 & 0 &\Big|& b-1
\end{pmatrix}
\text{ está en forma escalonada por filas.}
\]
\[
\operatorname{Rg}(A)=\text{nº filas no nulas}=2.
\]
Para la matriz aumentada \((A\,|\,B)\) la tercera fila se anula si y sólo si \(b-1=0\), es decir \(b=1\). Por tanto
\[
\operatorname{Rg}(A\,|\,B)=
\begin{cases}
2 & \text{si } b=1,\\[4pt]
3 & \text{si } b\neq 1.
\end{cases}
\]

\begin{itemize}
\item \textbf{Subcaso 3.1. \(a=-2\) y \(b=1\).}

En este caso \(\operatorname{Rg}(A)=\operatorname{Rg}(A\,|\,B)=2 \neq\) nº incógnitas \(=3\).  
\[
\boxed{\text{El sistema es un SCI si } a=-2 \text{ y } b=1.}
\]
\end{itemize}
\end{itemize}
\end{ejercicio}
\begin{ejercicio}[title=Ejercicio 31]
\begin{itemize}
\begin{itemize}
Resolvemos aplicando Gauss–Jordan sobre la ampliada con \(b=1\):
\[
(A\,|\,B)
\sim_f
\begin{pmatrix}
1 & 1 & -2 &\Big|& 1\\
0 & 3 & -3 &\Big|& 0\\
0 & 0 & 0 &\Big|& 0
\end{pmatrix}
\overset{\substack{\tfrac{1}{3}F_2\\{}}}{\sim_f}
\begin{pmatrix}
1 & 1 & -2 &\Big|& 1\\
0 & 1 & -1 &\Big|& 0\\
0 & 0 & 0 &\Big|& 0
\end{pmatrix}
\]
\[
\overset{\substack{F_1 - F_2\\{}}}{\sim_f}
\begin{pmatrix}
1 & 0 & -1 &\Big|& 1\\
0 & 1 & -1 &\Big|& 0\\
0 & 0 & 0 &\Big|& 0
\end{pmatrix} = H_f(A\,|\,B) \Rightarrow
\begin{cases}
x - z = 1,\\[4pt]
y - z = 0.
\end{cases}
\]

Tomando \(z=\lambda\in\mathbb{R}\) (parámetro libre) obtenemos
\[
\boxed{\,x = 1+\lambda,\quad y = \lambda,\quad z=\lambda,\qquad \lambda\in\mathbb{R}\,}
\]

\item \textbf{Subcaso 3.2. \(a=-2\) y \(b\neq 1\).}

Aquí \(\operatorname{Rg}(A)=2 \neq \operatorname{Rg}(A\,|\,B)=3\), por lo que el sistema es incompatible.

\[
\boxed{\text{El sistema es un SI si } a=-2 \text{ y } b\neq 1.}
\]
\end{itemize}

\item \textbf{Caso 4. \(b=0\).}

La matriz ampliada es
\[
(A\,|\,B)=
\begin{pmatrix}
1 & 0 & a &\Big|& 1\\
a & 0 & 1 &\Big|& a\\
1 & 0 & 1 &\Big|& 0
\end{pmatrix}
\overset{\substack{F_2-aF_1\\ F_3-F_1\\{}}}{\sim_f}
\begin{pmatrix}
1 & 0 & a &\Big|& 1\\
0 & 0 & 1-a^2 &\Big|& 0\\
0 & 0 & 1-a &\Big|& -1
\end{pmatrix}.
\]

La segunda fila es nula si y sólo si \(1-a^2=0\), es decir, \(a=\pm1\).

Para la matriz de coeficientes \(A\):
\begin{itemize}
\item La tercera fila es nula si y sólo si \(1-a=0\), es decir, \(a=1\).
\end{itemize}

Por tanto,
\[
\operatorname{Rg}(A)=
\begin{cases}
1, & a=1\\
2, & a=-1 \text{ o } (a\neq 1 \text{ y } a\neq -1)
\end{cases}
\]
\[
\operatorname{Rg}(A\,|\,B)=
\begin{cases}
2, & a=1 \text{ o } a=-1\\
3, & a\neq 1 \text{ y } a\neq -1
\end{cases}
\]
\end{itemize}
\end{ejercicio}
\begin{ejercicio}[title=Ejercicio 31]
\begin{itemize}
\begin{itemize}
\item \textbf{Subcaso 4.1. \(b=0\) y \(a=1\).}
\[
(A\,|\,B)\sim_f
\begin{pmatrix}
1 & 0 & 1 &\Big|& 1\\
0 & 0 & 0 &\Big|& 0\\
0 & 0 & 0 &\Big|& -1
\end{pmatrix}
\overset{\substack{F_2\leftrightarrow F_3\\{}}}{\sim_f}
\begin{pmatrix}
1 & 0 & 1 &\Big|& 1\\
0 & 0 & 0 &\Big|& -1\\
0 & 0 & 0 &\Big|& 0
\end{pmatrix},
\]
que está en forma escalonada por filas.

Entonces \(\operatorname{Rg}(A)=1 \neq \operatorname{Rg}(A\,|\,B)=2\).  
\[
\boxed{\text{El sistema es un SI si } b=0 \text{ y } a=1.}
\]

\item \textbf{Subcaso 4.2. \(b=0\) y \(a=-1\).}
\[
(A\,|\,B)\sim_f
\begin{pmatrix}
1 & 0 & -1 &\Big|& 1\\
0 & 0 & 0 &\Big|& 0\\
0 & 0 & 2 &\Big|& -1
\end{pmatrix}
\overset{\substack{F_2\leftrightarrow F_3\\{}}}{\sim_f}
\begin{pmatrix}
1 & 0 & -1 &\Big|& 1\\
0 & 0 & 2 &\Big|& -1\\
0 & 0 & 0 &\Big|& 0
\end{pmatrix}.
\]
que está en forma escalonada por filas.

Aquí \(\operatorname{Rg}(A)=\operatorname{Rg}(A\,|\,B)=2 \neq\) nº incógnitas \(=3\).  
\[
\boxed{\text{El sistema es un SCI si } b=0 \text{ y } a=-1.}
\]

Aplicando Gauss–Jordan:
\[
\overset{\tfrac12 F_2}{\sim_f}
\begin{pmatrix}
1 & 0 & -1 &\Big|& 1\\
0 & 0 & 1 &\Big|& -\tfrac12\\
0 & 0 & 0 &\Big|& 0
\end{pmatrix}
\overset{F_1+F_2}{\sim_f}
\begin{pmatrix}
1 & 0 & 0 &\Big|& \tfrac12\\
0 & 0 & 1 &\Big|& -\tfrac12\\
0 & 0 & 0 &\Big|& 0
\end{pmatrix}=H_f(A\,|\,B).
\]
Por tanto,
\[
\boxed{\,x=\frac12,\quad y=\lambda,\quad z=-\frac12,\qquad \lambda\in\mathbb{R}\,}
\]

\item \textbf{Subcaso 4.3. \(b=0\), \(a\neq 1\) y \(a\neq -1\).}
\[
(A\,|\,B)\sim_f
\begin{pmatrix}
1 & 0 & a &\Big|& 1\\
0 & 0 & 1-a^2 &\Big|& 0\\
0 & 0 & 1-a &\Big|& -1
\end{pmatrix}
\overset{\substack{F_2-(1+a)F_3\\{}}}{\sim_f}
\begin{pmatrix}
1 & 0 & a &\Big|& 1\\
0 & 0 & 0 &\Big|& 0\\
0 & 0 & 1-a &\Big|& -1
\end{pmatrix}
\]
\end{itemize}
\end{itemize}
\end{ejercicio}
\begin{ejercicio}[title=Ejercicio 31]
\begin{itemize}
\begin{itemize}
\[
\overset{\substack{F_2\leftrightarrow F_3\\{}}}{\sim_f}
\begin{pmatrix}
1 & 0 & a &\Big|& 1\\
0 & 0 & 1-a &\Big|& -1\\
0 & 0 & 0 &\Big|& 1+a
\end{pmatrix}
\text{ está en forma escalonada por filas.}
\]
Se cumple \(\operatorname{Rg}(A)=2 \neq \operatorname{Rg}(A\,|\,B)=3\).  
\[
\boxed{\text{El sistema es un SI si } b=0,\ a\neq 1 \text{ y } a\neq -1.}
\]
\end{itemize}
\end{itemize}
\end{ejercicio}

\section{Conceptos Adicionales}

\subsection{Traza}

\begin{ejercicio}[title=Ejercicio 7]
La traza de una matriz cuadrada $A$, denotada por $\operatorname{tr}(A)$, es la suma de los elementos que aparecen en la diagonal principal de $A$.  
Por tanto, si $A = (a_{ij}) \in \mathfrak{M}_n(\mathbb{K})$, entonces $\operatorname{tr}(A) = a_{1,1} + \cdots + a_{n,n}$.  
Si $A,B \in \mathfrak{M}_n(\mathbb{K})$ y $\lambda \in \mathbb{K}$, justifique las siguientes propiedades:

\begin{enumerate}
\item[a)] $\operatorname{tr}(A+B) = \operatorname{tr}(A) + \operatorname{tr}(B)$,
\item[b)] $\operatorname{tr}(\lambda A) = \lambda \operatorname{tr}(A)$,
\item[c)] $\operatorname{tr}(A \cdot B) = \operatorname{tr}(B \cdot A)$.
\end{enumerate}

\tcblower
\textbf{Resolución:}

\textbf{a)} $\operatorname{tr}(A+B) = \operatorname{tr}(A) + \operatorname{tr}(B)$

\[
A =
\begin{pmatrix}
a_{11} & a_{12} & \cdots & a_{1n} \\
a_{21} & a_{22} & \cdots & a_{2n} \\
\vdots & \vdots & \ddots & \vdots \\
a_{n1} & a_{n2} & \cdots & a_{nn}
\end{pmatrix},
\qquad
\operatorname{tr}(A)=a_{11}+a_{22}+\cdots+a_{nn}.
\]

\[
B =
\begin{pmatrix}
b_{11} & b_{12} & \cdots & b_{1n} \\
b_{21} & b_{22} & \cdots & b_{2n} \\
\vdots & \vdots & \ddots & \vdots \\
b_{n1} & b_{n2} & \cdots & b_{nn}
\end{pmatrix},
\qquad
\operatorname{tr}(B)=b_{11}+b_{22}+\cdots+b_{nn}.
\]

Entonces
\[
A+B =
\begin{pmatrix}
a_{11}+b_{11} & a_{12}+b_{12} & \cdots & a_{1n}+b_{1n} \\
a_{21}+b_{21} & a_{22}+b_{22} & \cdots & a_{2n}+b_{2n} \\
\vdots & \vdots & \ddots & \vdots \\
a_{n1}+b_{n1} & a_{n2}+b_{n2} & \cdots & a_{nn}+b_{nn}
\end{pmatrix}.
\]

Por definición de la traza,
\[
\operatorname{tr}(A+B)
= (a_{11}+b_{11}) + (a_{22}+b_{22}) + \cdots + (a_{nn}+b_{nn})
\]
\[
= (a_{11}+a_{22}+\cdots+a_{nn}) + (b_{11}+b_{22}+\cdots+b_{nn})
= \boxed{\operatorname{tr}(A)+\operatorname{tr}(B)}.
\]

\textbf{b)} $\operatorname{tr}(\lambda A) = \lambda \operatorname{tr}(A)$

\[
A =
\begin{pmatrix}
a_{11} & a_{12} & \cdots & a_{1n} \\
a_{21} & a_{22} & \cdots & a_{2n} \\
\vdots & \vdots & \ddots & \vdots \\
a_{n1} & a_{n2} & \cdots & a_{nn}
\end{pmatrix}.
\]

\end{ejercicio}
\begin{ejercicio}[title=Ejercicio 7]

Entonces
\[
\lambda A
=
\lambda
\begin{pmatrix}
a_{11} & a_{12} & \cdots & a_{1n} \\
a_{21} & a_{22} & \cdots & a_{2n} \\
\vdots & \vdots & \ddots & \vdots \\
a_{n1} & a_{n2} & \cdots & a_{nn}
\end{pmatrix}
=
\begin{pmatrix}
\lambda a_{11} & \lambda a_{12} & \cdots & \lambda a_{1n} \\
\lambda a_{21} & \lambda a_{22} & \cdots & \lambda a_{2n} \\
\vdots & \vdots & \ddots & \vdots \\
\lambda a_{n1} & \lambda a_{n2} & \cdots & \lambda a_{nn}
\end{pmatrix}.
\]

Por tanto,
\(
\operatorname{tr}(\lambda A)
= \lambda a_{11} + \lambda a_{22} + \cdots + \lambda a_{nn}
= \lambda (a_{11} + a_{22} + \cdots + a_{nn})
= \boxed{\lambda\,\operatorname{tr}(A)}.
\)

\textbf{c)} $\operatorname{tr}(A \cdot B) = \operatorname{tr}(B \cdot A)$

\[
A =
\begin{pmatrix}
a_{11} & a_{12} & \cdots & a_{1n} \\
a_{21} & a_{22} & \cdots & a_{2n} \\
\vdots & \vdots & \ddots & \vdots \\
a_{n1} & a_{n2} & \cdots & a_{nn}
\end{pmatrix},
\qquad
B =
\begin{pmatrix}
b_{11} & b_{12} & \cdots & b_{1n} \\
b_{21} & b_{22} & \cdots & b_{2n} \\
\vdots & \vdots & \ddots & \vdots \\
b_{n1} & b_{n2} & \cdots & b_{nn}
\end{pmatrix}.
\]

% necesita \usepackage{graphicx} en el preámbulo
\[
\scalebox{0.75}{$
A\cdot B =
\begin{pmatrix}
a_{11}b_{11} + a_{12}b_{21} + \cdots + a_{1n}b_{n1} & 
a_{11}b_{12} + a_{12}b_{22} + \cdots + a_{1n}b_{n2} & 
\cdots & 
a_{11}b_{1n} + a_{12}b_{2n} + \cdots + a_{1n}b_{nn} \\[6pt]
a_{21}b_{11} + a_{22}b_{21} + \cdots + a_{2n}b_{n1} & 
a_{21}b_{12} + a_{22}b_{22} + \cdots + a_{2n}b_{n2} & 
\cdots & 
a_{21}b_{1n} + a_{22}b_{2n} + \cdots + a_{2n}b_{nn} \\[6pt]
\vdots & \vdots & \ddots & \vdots \\[6pt]
a_{n1}b_{11} + a_{n2}b_{21} + \cdots + a_{nn}b_{n1} & 
a_{n1}b_{12} + a_{n2}b_{22} + \cdots + a_{nn}b_{n2} & 
\cdots & 
a_{n1}b_{1n} + a_{n2}b_{2n} + \cdots + a_{nn}b_{nn}
\end{pmatrix}
$}
\]

\[
\scalebox{0.75}{$
B\cdot A =
\begin{pmatrix}
b_{11}a_{11} + b_{12}a_{21} + \cdots + b_{1n}a_{n1} & 
b_{11}a_{12} + b_{12}a_{22} + \cdots + b_{1n}a_{n2} & 
\cdots & 
b_{11}a_{1n} + b_{12}a_{2n} + \cdots + b_{1n}a_{nn} \\[6pt]
b_{21}a_{11} + b_{22}a_{21} + \cdots + b_{2n}a_{n1} & 
b_{21}a_{12} + b_{22}a_{22} + \cdots + b_{2n}a_{n2} & 
\cdots & 
b_{21}a_{1n} + b_{22}a_{2n} + \cdots + b_{2n}a_{nn} \\[6pt]
\vdots & \vdots & \ddots & \vdots \\[6pt]
b_{n1}a_{11} + b_{n2}a_{21} + \cdots + b_{nn}a_{n1} & 
b_{n1}a_{12} + b_{n2}a_{22} + \cdots + b_{nn}a_{n2} & 
\cdots & 
b_{n1}a_{1n} + b_{n2}a_{2n} + \cdots + b_{nn}a_{nn}
\end{pmatrix}
$}
\]

La traza es la suma de la diagonal:
\[
\operatorname{tr}(A\cdot B)
=
\bigl(a_{11}b_{11} + a_{12}b_{21} + \cdots + a_{1n}b_{n1}\bigr)
+ \bigl(a_{21}b_{12} + a_{22}b_{22} + \cdots + a_{2n}b_{n2}\bigr)
\]
\[
+ \cdots +
\bigl(a_{n1}b_{1n} + a_{n2}b_{2n} + \cdots + a_{nn}b_{nn}\bigr).
\]

Para $B\cdot A$:
\[
\operatorname{tr}(B\cdot A)
=
\bigl(b_{11}a_{11} + b_{12}a_{21} + \cdots + b_{1n}a_{n1}\bigr)
+ \bigl(b_{21}a_{12} + b_{22}a_{22} + \cdots + b_{2n}a_{n2}\bigr)
\]
\[
+ \cdots +
\bigl(b_{n1}a_{1n} + b_{n2}a_{2n} + \cdots + b_{nn}a_{nn}\bigr)
\]
\vspace{0.2em}
\end{ejercicio}
\begin{ejercicio}[title=Ejercicio 7]
Reescribimos cada término ($b_{ji}a_{ij} \rightarrow a_{ij}b_{ji}$):
\[
\qquad\qquad=
\Bigl(
\colorbox{Red!30}{$a_{11}b_{11}$} +
\colorbox{Orange!30}{$a_{21}b_{12}$} +
\cdots +
\colorbox{Yellow!30}{$a_{n1}b_{1n}$}
\Bigr)
+
\Bigl(
\colorbox{LightGreen!30}{$a_{12}b_{21}$} +
\colorbox{Green!30}{$a_{22}b_{22}$} +
\cdots +
\colorbox{LightBlue!30}{$a_{n2}b_{2n}$}
\Bigr)
\]
\[
\qquad\qquad
+ \cdots +
\Bigl(
\colorbox{Blue!30}{$a_{1n}b_{n1}$} +
\colorbox{Purple!30}{$a_{2n}b_{n2}$} +
\cdots +
\colorbox{Violet!30}{$a_{nn}b_{nn}$}
\Bigr).
\]

\vspace{0.2em}

\[
\qquad\qquad=
\Bigl(
\colorbox{Red!30}{$a_{11}b_{11}$} +
\colorbox{LightGreen!30}{$a_{12}b_{21}$} +
\cdots +
\colorbox{Blue!30}{$a_{1n}b_{n1}$}
\Bigr)
+
\Bigl(
\colorbox{Orange!30}{$a_{21}b_{12}$} +
\colorbox{Green!30}{$a_{22}b_{22}$} +
\cdots +
\colorbox{Purple!30}{$a_{2n}b_{n2}$}
\Bigr)
\]
\[
+ \cdots +
\Bigl(
\colorbox{Yellow!30}{$a_{n1}b_{1n}$} +
\colorbox{LightBlue!30}{$a_{n2}b_{2n}$} +
\cdots +
\colorbox{Violet!30}{$a_{nn}b_{nn}$}
\Bigr).
\]
\[
= \boxed{\operatorname{tr}(A \cdot B)}
\]
\end{ejercicio}

\begin{ejercicio}[title=Ejercicio 15]
Dada la matriz regular 
\[
A =
\begin{pmatrix}
1 & 1 & 2\\
0 & 1 & 8\\
1 & 2 & 4
\end{pmatrix} \in \mathfrak{M}_3(\mathbb{Z}_{11}),
\]
la traza de $A^{-1}$ es igual a:
\begin{enumerate}
\item[(a)] 2 \quad (b) 7 \quad (c) 8 \quad (d) 9
\end{enumerate}

\tcblower
\textbf{Resolución:}

\[
(A \,|\, I_3) =
\begin{pmatrix}
1 & 1 & 2 &\Big|& 1 & 0 & 0\\
0 & 1 & 8 &\Big|& 0 & 1 & 0\\
1 & 2 & 4 &\Big|& 0 & 0 & 1
\end{pmatrix}
\overset{\substack{F_3 - F_1 \\ {}}}{\sim_f}
\begin{pmatrix}
1 & 1 & 2 &\Big|& 1 & 0 & 0\\
0 & 1 & 8 &\Big|& 0 & 1 & 0\\
0 & 1 & 2 &\Big|& -1 & 0 & 1
\end{pmatrix}
\]

\[
\overset{\substack{F_3 - F_2 \\ {}}}{\sim_f}
\begin{pmatrix}
1 & 1 & 2 &\Big|& 1 & 0 & 0\\
0 & 1 & 8 &\Big|& 0 & 1 & 0\\
0 & 0 & 5 &\Big|& -1 & -1 & 1
\end{pmatrix}
\overset{\substack{F_1 - F_2 \\ -2F_3\\ {}}}{\sim_f}
\begin{pmatrix}
1 & 0 & 5 &\Big|& 1 & -1 & 0\\
0 & 1 & -3 &\Big|& 0 & 1 & 0\\
0 & 0 & 1 &\Big|& 2 & 2 & -2
\end{pmatrix}
\]

\[
\overset{\substack{F_1 - 5F_3 \\ F_2 + 3F_3\\ {}}}{\sim_f}
\begin{pmatrix}
1 & 0 & 0 &\Big|& 2 & 0 & -1\\
0 & 1 & 0 &\Big|& 6 & 7 & 5\\
0 & 0 & 1 &\Big|& 2 & 2 & -2
\end{pmatrix}
= (I_3 \,|\, A^{-1})
\]

\noindent
La traza de $A^{-1}$ es
\[
\operatorname{tr}(A^{-1}) = 2 + 7 + (-2) = 7.
\]

\[
\boxed{(b) \, 7}
\]
\end{ejercicio}

\subsection{Matrices estocásticas o de probabilidad}

\begin{ejercicio}[title=Ejercicio 43 (⋆)]
Una matriz cuadrada $A = (a_{i,j}) \in \mathfrak{M}_n(\mathbb{R})$ se dice que es una \textit{matriz estocástica o matriz de probabilidad}, si verifica las condiciones siguientes:
\begin{itemize}
\item $a_{i,j} \ge 0$ para todo $i,j \in \{1,2,\ldots,n\}$.
\item $\displaystyle \sum_{j=1}^n a_{i,j} = 1$ para todo $i \in \{1,2,\ldots,n\}$.
\end{itemize}
Una matriz estocástica $A$ se denomina \emph{doblemente estocástica} si además
\[
\displaystyle \sum_{i=1}^n a_{i,j} = 1 \quad \text{para todo } j \in \{1,2,\ldots,n\}.
\]
Justifique las siguientes afirmaciones:

\begin{enumerate}
\item[a)] Si $A$ y $B$ son matrices estocásticas del mismo orden, entonces $A \cdot B$ es una matriz estocástica. Como consecuencia, deduzca que si $A$ y $B$ son matrices doblemente estocásticas del mismo orden, entonces $A \cdot B$ también es doblemente estocástica.
\item[b)] $A$ es doblemente estocástica si y sólo si $A$ y $A^T$ son estocásticas.
\item[c)] Si $A$ es simétrica y estocástica, entonces $A$ es doblemente estocástica.
\end{enumerate}

Por último, ponga un ejemplo de una matriz estocástica e invertible cuya inversa no sea estocástica.

\tcblower
\textbf{Resolución:}

\vspace{0.5em}

\textbf{a)}

Sean $A=(a_{i,k})$ y $B=(b_{k,j})$ matrices estocásticas por columnas del mismo orden.  
El elemento $c_{i,j}$ del producto $A\cdot B$ es
\[
c_{i,j}=\sum_{k=1}^n a_{i,k} b_{k,j}.
\]

Comprobamos que $A\cdot B$ es estocástica (la suma de cada columna $j$ debe ser 1):

\begin{itemize}
\item \textbf{No negatividad:} Como $a_{i,k}\ge 0$ y $b_{k,j}\ge 0$ para todo $i,k,j$, entonces $c_{i,j} = \sum a_{i,k} b_{k,j} \ge 0$.
\item \textbf{Suma por columnas:} Se tiene que cumplir que para cada columna $j$ lo siguiente. Desarrollamos usando la definición del producto de matrices:
\[
\sum_{i=1}^n c_{i,j} = 1
\]
\[
\sum_{i=1}^n \left( \sum_{k=1}^n a_{i,k} b_{k,j} \right) = 1
\]
\end{itemize}
\end{ejercicio}
\begin{ejercicio}[title=Ejercicio 43 (⋆)]
\begin{itemize}
Intercambiamos el orden de las sumatorias, sacando fuera $b_{k,j}$ ya que no depende del índice de suma $i$:
\[
\sum_{k=1}^n \left( b_{k,j} \underbrace{\sum_{i=1}^n a_{i,k}}_{=1} \right) = 1
\]
Como $A$ es estocástica, la suma de su columna $k$ es 1. Nos queda:
\[
\underbrace{\sum_{k=1}^n b_{k,j}}_{1} \cdot 1 = 1
\]
Como $B$ es estocástica, la suma de su columna $j$ es 1, por lo que:
\[
1 = 1
\]
\end{itemize}

Por tanto,

\[
\boxed{\text{$A\cdot B$ es estocástica.}}
\]

Si además $A$ y $B$ son doblemente estocásticas, entonces también se verifica la suma por filas. Para cada fila $i$ de $A \cdot B$:
\[
\sum_{j=1}^n c_{i,j} = \sum_{j=1}^n \left( \sum_{k=1}^n a_{i,k} b_{k,j} \right)
\]
Intercambiamos el orden de las sumatorias, sacando $a_{i,k}$ fuera de la suma en $j$:
\[
\sum_{k=1}^n a_{i,k} \left( \underbrace{\sum_{j=1}^n b_{k,j}}_{=1} \right) = \sum_{k=1}^n a_{i,k} \cdot 1
\]
Como $B$ es estocástica por filas, la suma interior es 1. Finalmente, como $A$ también es estocástica por filas:
\[
\underbrace{\sum_{k=1}^n a_{i,k}}_{1} = 1 \quad \Rightarrow \quad 1 = 1
\]
Al cumplirse ambas condiciones (suma de columnas y suma de filas igual a 1), concluimos que:
\[
\boxed{\text{$A \cdot B$ es doblemente estocástica.}}
\]

\textbf{b)} $A$ es doblemente estocástica si y sólo si $A$ y $A^T$ son estocásticas.

\vspace{0.5em}

Una matriz es estocástica si sus elementos son no negativos y la suma de sus columnas es 1.
\end{ejercicio}
\begin{ejercicio}[title=Ejercicio 43 (⋆)]
\begin{itemize}
    \item[$\boxed{\Rightarrow}$] Si $A$ es doblemente estocástica, por definición sus columnas suman 1 (luego $A$ es estocástica) y sus filas también suman 1. Como las filas de $A$ son las columnas de $A^T$, entonces las columnas de $A^T$ suman 1. Por tanto, $A^T$ es estocástica.
    
    \item[$\boxed{\Leftarrow}$] Si $A$ es estocástica, sus columnas suman 1. Si $A^T$ también es estocástica, sus columnas suman 1; pero las columnas de $A^T$ son las filas de $A$. Por tanto, las filas de $A$ suman 1. Al sumar 1 tanto columnas como filas, $A$ es doblemente estocástica.
\end{itemize}
\[
\boxed{\text{$A$ doblemente estocástica $\iff A$ y $A^T$ estocásticas.}}
\]

\vspace{1em}

\textbf{c)} Si $A$ es simétrica y estocástica, entonces $A$ es doblemente estocástica.

\vspace{0.5em}

Sea $A$ una matriz simétrica ($A = A^T$) y estocástica (sus columnas suman 1). 
\begin{itemize}
    \item Por ser estocástica, sabemos que la suma de cualquier columna $j$ es 1: $\displaystyle\sum_{i=1}^n a_{i,j} = 1$.
    \item Para comprobar si es doblemente estocástica, debemos ver si sus filas suman 1. Sumamos los elementos de una fila $i$ cualquiera:
    \[
    \sum_{j=1}^n a_{i,j} \overset{\substack{\text{$A$ simétrica}\\{}}}{=} \sum_{j=1}^n a_{j,i}
    \]
    \item Esta última expresión es, precisamente, la suma de los elementos de la columna $i$ de $A$. Como $A$ es estocástica, esta suma es igual a 1.
\end{itemize}
Como sus columnas y sus filas suman 1, $A$ es doblemente estocástica.

\[
\boxed{\text{$A$ simétrica y estocástica $\Rightarrow$ $A$ doblemente estocástica.}}
\]

\vspace{1em}

\textbf{Ejemplo de matriz estocástica e invertible cuya inversa no es estocástica:}

Sea la matriz
\[
A = \begin{pmatrix}
0.8 & 0.1 \\
0.2 & 0.9
\end{pmatrix}.
\]

\vspace{0.5em}

\textbf{Comprobación de que $A$ es estocástica:}  
\begin{itemize}
\item Suma de la primera columna: $0.8 + 0.2 = 1$  
\item Suma de la segunda columna: $0.1 + 0.9 = 1$  
\end{itemize}

Todas los elementos son no negativos y las columnas suman 1, por lo tanto $A$ es estocástica por columnas.

\vspace{0.5em}

\textbf{Comprobación de que $A$ es invertible:}  

\[
|A| = (0.8)(0.9) - (0.1)(0.2) = 0.72 - 0.02 = 0.7 \neq 0
\]
Por lo tanto, $A$ es invertible.
\end{ejercicio}
\begin{ejercicio}[title=Ejercicio 43 (⋆)]
\textbf{Cálculo de la inversa de $A$:}
\[
A^{-1} = \frac{1}{|A|} \overline{\begin{pmatrix} 0.8 & 0.1 \\ 0.2 & 0.9 \end{pmatrix}}^T = \frac{1}{0.7} \begin{pmatrix} 0.9 & -0.1 \\ -0.2 & 0.8 \end{pmatrix}
= \begin{pmatrix} 9/7 & -1/7 \\ -2/7 & 8/7 \end{pmatrix}.
\]

\vspace{0.5em}

\textbf{Comprobación de que $A^{-1}$ no es estocástica por columnas:}  
\begin{itemize}
\item Primera columna: $9/7 + (-2/7) = 7/7 = 1$ 
\item Segunda columna: $-1/7 + 8/7 = 7/7 = 1$  
\end{itemize}

Las columnas suman 1, pero hay elementos negativos ($-1/7$ y $-2/7$), por lo tanto $A^{-1}$ no es estocástica.

\vspace{0.5em}

\textbf{Conclusión:}  
\[
\boxed{A = \begin{pmatrix}
0.8 & 0.1 \\
0.2 & 0.9
\end{pmatrix}} \text{ es estocástica e invertible, pero } A^{-1} \text{ no es estocástica.}
\]

\end{ejercicio}

\end{document}
